%% Draft-only owner placeholder (ruling R15, 2026-09-09). Prints ONLY when the
%% build defines \DRAFTTODO (e.g. pdflatex "\def\DRAFTTODO{1}%% Working draft. Structure is generated; prose is authored here.
%% Build from the repository root with: bash scripts/build_paper.sh <this paper dir>
\documentclass[11pt]{article}

\usepackage[utf8]{inputenc}
\usepackage[T1]{fontenc}
\usepackage{amsmath,amssymb,amsthm}
\usepackage{array}
\usepackage{booktabs}
\usepackage{graphicx}
\usepackage{tikz}
\usepackage{placeins}
\usepackage{caption}
\usepackage[margin=1in]{geometry}
% For every included figure, pdfTeX records the absolute path it was read from
% in a /PTEX.FileName object.  Those objects are not rendered, so a text dump of
% the PDF does not show them, but they travel with the file and spell out the
% build machine's home directory, user name and worktree layout.  Suppressing
% them costs nothing: they exist for debugging, not for the reader.
\pdfsuppressptexinfo=-1

\usepackage[hidelinks]{hyperref}
% Blind-review copies are compiled as
%   pdflatex "\def\ANONYMOUS{1}%% Working draft. Structure is generated; prose is authored here.
%% Build from the repository root with: bash scripts/build_paper.sh <this paper dir>
\documentclass[11pt]{article}

\usepackage[utf8]{inputenc}
\usepackage[T1]{fontenc}
\usepackage{amsmath,amssymb,amsthm}
\usepackage{array}
\usepackage{booktabs}
\usepackage{graphicx}
\usepackage{tikz}
\usepackage{placeins}
\usepackage{caption}
\usepackage[margin=1in]{geometry}
% For every included figure, pdfTeX records the absolute path it was read from
% in a /PTEX.FileName object.  Those objects are not rendered, so a text dump of
% the PDF does not show them, but they travel with the file and spell out the
% build machine's home directory, user name and worktree layout.  Suppressing
% them costs nothing: they exist for debugging, not for the reader.
\pdfsuppressptexinfo=-1

\usepackage[hidelinks]{hyperref}
% Blind-review copies are compiled as
%   pdflatex "\def\ANONYMOUS{1}%% Working draft. Structure is generated; prose is authored here.
%% Build from the repository root with: bash scripts/build_paper.sh <this paper dir>
\documentclass[11pt]{article}

\usepackage[utf8]{inputenc}
\usepackage[T1]{fontenc}
\usepackage{amsmath,amssymb,amsthm}
\usepackage{array}
\usepackage{booktabs}
\usepackage{graphicx}
\usepackage{tikz}
\usepackage{placeins}
\usepackage{caption}
\usepackage[margin=1in]{geometry}
% For every included figure, pdfTeX records the absolute path it was read from
% in a /PTEX.FileName object.  Those objects are not rendered, so a text dump of
% the PDF does not show them, but they travel with the file and spell out the
% build machine's home directory, user name and worktree layout.  Suppressing
% them costs nothing: they exist for debugging, not for the reader.
\pdfsuppressptexinfo=-1

\usepackage[hidelinks]{hyperref}
% Blind-review copies are compiled as
%   pdflatex "\def\ANONYMOUS{1}\input{main}"
% Identified output is the default.  hyperref writes the name into the PDF
% /Author key as well as the title page, so both go through this one macro.
\ifdefined\ANONYMOUS\newcommand{\PDFAUTHOR}{}\else\newcommand{\PDFAUTHOR}{Zeyu Fu}\fi
\hypersetup{
  pdftitle={From directional sensing to task completion: the operating envelope of target-position feedforward against a slow integral in wind-disturbed quadrotor waypoint capture},
  pdfauthor={\PDFAUTHOR},
  pdfsubject={Operating envelope of target-position feedforward against a slow integral},
  pdfkeywords={disturbance feedforward, integral control, wind estimation, thresholded endpoint, task allocation, quadrotor swarm, reproducibility}
}
\numberwithin{equation}{section}

\graphicspath{{../figs/out/}}
\setlength{\emergencystretch}{2em}

%% Every float here is small enough to share a page with text; the defaults are
%% conservative enough to push each one onto a page of its own.
\renewcommand{\topfraction}{0.9}
\renewcommand{\bottomfraction}{0.7}
\renewcommand{\textfraction}{0.08}
\renewcommand{\floatpagefraction}{0.8}
\setcounter{topnumber}{2}
\setcounter{totalnumber}{3}
\captionsetup[table]{skip=5pt}
\captionsetup[figure]{skip=4pt}

%% Numbers and tables are emitted by the figure code from hash-bound evidence.
\input{generated_numbers}
\input{captions}
\input{theorems}
\newcommand{\generatedtable}[1]{\input{#1}}

\title{From directional sensing to task completion: the operating envelope of
target-position feedforward against a slow integral in wind-disturbed quadrotor
waypoint capture}
\ifdefined\ANONYMOUS
\author{Author and affiliation withheld for blind review}
\else
\author{Zeyu Fu\\
Army Medical University, Chongqing, China\\
\texttt{fuzeyu99@126.com}\\
ORCID: 0009-0001-8329-0108}
\fi
\date{}

\begin{document}
\maketitle

\begin{abstract}
\input{abstract}
\end{abstract}

\paragraph{Keywords} disturbance feedforward, integral control, wind estimation, thresholded endpoint, task allocation, quadrotor swarm, reproducibility

\input{body}

{\footnotesize
\bibliographystyle{unsrt}
\bibliography{references}
}

\end{document}
"
% Identified output is the default.  hyperref writes the name into the PDF
% /Author key as well as the title page, so both go through this one macro.
\ifdefined\ANONYMOUS\newcommand{\PDFAUTHOR}{}\else\newcommand{\PDFAUTHOR}{Zeyu Fu}\fi
\hypersetup{
  pdftitle={From directional sensing to task completion: the operating envelope of target-position feedforward against a slow integral in wind-disturbed quadrotor waypoint capture},
  pdfauthor={\PDFAUTHOR},
  pdfsubject={Operating envelope of target-position feedforward against a slow integral},
  pdfkeywords={disturbance feedforward, integral control, wind estimation, thresholded endpoint, task allocation, quadrotor swarm, reproducibility}
}
\numberwithin{equation}{section}

\graphicspath{{../figs/out/}}
\setlength{\emergencystretch}{2em}

%% Every float here is small enough to share a page with text; the defaults are
%% conservative enough to push each one onto a page of its own.
\renewcommand{\topfraction}{0.9}
\renewcommand{\bottomfraction}{0.7}
\renewcommand{\textfraction}{0.08}
\renewcommand{\floatpagefraction}{0.8}
\setcounter{topnumber}{2}
\setcounter{totalnumber}{3}
\captionsetup[table]{skip=5pt}
\captionsetup[figure]{skip=4pt}

%% Numbers and tables are emitted by the figure code from hash-bound evidence.
% Generated by the figure script from hash-bound evidence. Do not hand-edit.
\newcommand{\NSeeds}{10}
\newcommand{\NWindLevels}{7}
\newcommand{\NEpisodesOracle}{140}
\newcommand{\NEpisodesNoise}{150}
\newcommand{\NDrones}{4}
\newcommand{\SigmaGrid}{0, 0.1, 0.3, 0.5, 1}
\newcommand{\CaptureRadius}{0.15}
\newcommand{\HoverThrust}{0.265}
\newcommand{\DroneMass}{27}
\newcommand{\KneeLevel}{0.15}
\newcommand{\KneePctHover}{57}
\newcommand{\SaturationLevel}{0.10}
\newcommand{\CollapseLevel}{0.20}
\newcommand{\KneeAgnostic}{0.05}
\newcommand{\KneeAgnosticSD}{0.10}
\newcommand{\KneeAware}{0.60}
\newcommand{\KneeAwareSD}{0.25}
\newcommand{\KneeDelta}{0.55}
\newcommand{\KneeP}{0.00390625}
\newcommand{\BonferroniAlpha}{0.0071}
\newcommand{\PreregSigma}{0.25}
\newcommand{\NearestSigma}{0.30}
\newcommand{\NearestSigmaDelta}{0.55}
\newcommand{\NearestSigmaLower}{0.44}
\newcommand{\NearestSigmaP}{0.00195}
\newcommand{\SigmaZeroDelta}{0.55}
\newcommand{\SigmaZeroAware}{0.60}
\newcommand{\BestSigma}{0.50}
\newcommand{\BestSigmaDelta}{0.65}
\newcommand{\MaxSigma}{1.00}
\newcommand{\MaxSigmaDelta}{0.30}
\newcommand{\MaxSigmaLower}{0.09}
\newcommand{\MaxSigmaP}{0.0312}
\newcommand{\WeakestLower}{0.09}
\newcommand{\NoiseComparisons}{15}
\newcommand{\NoiseBonferroni}{0.0033}
\newcommand{\BridgeLowerPairs}{8}
\newcommand{\BridgeDistanceShift}{-0.018}
\newcommand{\BridgeCompletionGain}{0.55}
\newcommand{\AwareTopSeeds}{7}
\newcommand{\AwareZeroSeeds}{1}
\newcommand{\AgnosticZeroSeeds}{8}
\newcommand{\MissAgnosticKnee}{0.244}
\newcommand{\MissAwareKnee}{0.226}
\newcommand{\MissAgnosticCollapse}{0.597}
\newcommand{\MissAwareCollapse}{0.595}
\newcommand{\NHeadings}{6}
\newcommand{\LargestHeadingGroup}{4}
\newcommand{\HeadingGroupSpan}{1.60}
\newcommand{\UnsampledArc}{152}
\newcommand{\NRepeats}{7}
\newcommand{\RepeatTol}{2}
\newcommand{\ClosestRepeat}{0.05}
\newcommand{\WidestRepeat}{1.60}
\newcommand{\NearestDistinct}{25.2}
\newcommand{\NInformativeLevels}{3}
\newcommand{\NDegenerateLevels}{4}
\newcommand{\RepeatGapKnee}{0.75}
\newcommand{\RepeatGapKneeDrones}{3}
\newcommand{\RepeatGapShoulder}{0.25}
\newcommand{\RepeatGapShoulderDrones}{1}
\newcommand{\NDisagreeingLevels}{2}
\newcommand{\RepeatDisagreeKnee}{6}
\newcommand{\RepeatDisagreeShoulder}{4}
\newcommand{\RepeatDisagreeCollapse}{0}
\newcommand{\CollapseSeedSpread}{0.25}
\newcommand{\NEpisodesDesigned}{4032}
\newcommand{\NAssignedHeadings}{24}
\newcommand{\HeadingStep}{15}
\newcommand{\NScenes}{12}
\newcommand{\NPairsPerForce}{288}
\newcommand{\DesignedArc}{15}
\newcommand{\DesignedKneeDelta}{0.62}
\newcommand{\DesignedKneeLower}{0.57}
\newcommand{\DesignedKneeUpper}{0.69}
\newcommand{\DesignedKneeAgnostic}{0.08}
\newcommand{\DesignedKneeAware}{0.70}
\newcommand{\DesignedKneeAdverse}{1}
\newcommand{\DesignedShareHeading}{20}
\newcommand{\DesignedShareScene}{15}
\newcommand{\DesignedShareResidual}{65}
\newcommand{\ZeroSeparationGap}{1.25}
\newcommand{\ZeroSeparationLow}{-0.25}
\newcommand{\ZeroSeparationHigh}{+1.00}
\newcommand{\HeadingOnlyGap}{1.25}
\newcommand{\DesignedShoulderDelta}{0.030}
\newcommand{\DesignedCollapseDelta}{0.044}
\newcommand{\NReproEpisodes}{140}
\newcommand{\NReproFields}{6}
\newcommand{\NReplanQuiet}{6}
\newcommand{\NReplanLoud}{4}
\newcommand{\ReplanQuietDelta}{0.67}
\newcommand{\ReplanLoudDelta}{0.38}
\newcommand{\ReplanKneeAgnostic}{3}
\newcommand{\ReplanKneeAware}{6}
\newcommand{\NTemporalEpisodes}{1080}
\newcommand{\NTemporalProfiles}{3}
\newcommand{\NTemporalEstimators}{3}
\newcommand{\NTemporalSeeds}{30}
\newcommand{\NTemporalLegacyRows}{60}
\newcommand{\TemporalConstantDelta}{0.63}
\newcommand{\TemporalConstantLower}{0.52}
\newcommand{\TemporalConstantUpper}{0.75}
\newcommand{\TemporalConstantP}{7.45e-9}
\newcommand{\TemporalSlowDelta}{0.42}
\newcommand{\TemporalSlowLower}{0.29}
\newcommand{\TemporalSlowUpper}{0.56}
\newcommand{\TemporalSlowP}{5.13e-6}
\newcommand{\TemporalFastDelta}{0.60}
\newcommand{\TemporalFastLower}{0.48}
\newcommand{\TemporalFastUpper}{0.72}
\newcommand{\TemporalFastP}{7.45e-9}
\newcommand{\TemporalTrackingConstant}{3.9}
\newcommand{\TemporalTrackingSlow}{9.6}
\newcommand{\TemporalTrackingFast}{11.9}
\newcommand{\TemporalShoulderMax}{0.11}
\newcommand{\NAblationEpisodes}{810}
\newcommand{\AllocOnlyConstant}{0.01}
\newcommand{\AllocOnlySlow}{0.00}
\newcommand{\AllocOnlyFast}{0.01}
\newcommand{\CompleteConstant}{0.63}
\newcommand{\CompleteSlow}{0.42}
\newcommand{\CompleteFast}{0.60}
\newcommand{\NGainEpisodes}{1890}
\newcommand{\FeedforwardGainDefault}{0.35}
\newcommand{\GainPlateauMin}{0.50}
\newcommand{\GainPlateauMax}{0.71}
\newcommand{\GainDefaultConstant}{0.63}
\newcommand{\GainDefaultSlow}{0.42}
\newcommand{\GainDefaultFast}{0.60}
\newcommand{\GainEdgeConstant}{0.67}
\newcommand{\GainEdgeSlow}{0.52}
\newcommand{\GainEdgeFast}{0.58}
\newcommand{\NGeometryEpisodes}{1440}
\newcommand{\NGeometryConditions}{8}
\newcommand{\NGeometrySeeds}{30}
\newcommand{\NGeometryClusters}{30}
\newcommand{\NGeometryTermSensitivity}{120}
\newcommand{\NCrossedClusters}{12}
\newcommand{\NCrossedCells}{288}
\newcommand{\NCrossedBootstrap}{10000}
\newcommand{\GeometryKneeCrossingP}{0.00423}
\newcommand{\GeometryKneeCrossingHolm}{0.0381}
\newcommand{\GeometryKneeSpread}{-0.06}
\newcommand{\GeometryKneeSpreadLower}{-0.14}
\newcommand{\GeometryKneeSpreadUpper}{0.01}
\newcommand{\GeometryKneeRange}{-0.05}
\newcommand{\GeometryKneeRangeLower}{-0.10}
\newcommand{\GeometryKneeRangeUpper}{0.00}
\newcommand{\GeometryKneeCrossing}{0.12}
\newcommand{\GeometryKneeCrossingLower}{0.04}
\newcommand{\GeometryKneeCrossingUpper}{0.20}
\newcommand{\GeometryMDEFamily}{9}
\newcommand{\GeometryMDEAlpha}{0.0056}
\newcommand{\GeometryMDEPower}{0.80}
\newcommand{\GeometryRangeClusterN}{30}
\newcommand{\GeometryRangeClusterSD}{0.139}
\newcommand{\GeometryRangeSensitivityLower}{-0.11}
\newcommand{\GeometryRangeSensitivityUpper}{0.01}
\newcommand{\GeometryRangeMDE}{0.098}
\newcommand{\GeometryKneeNominal}{0.54}
\newcommand{\GeometryKneeBest}{0.66}
\newcommand{\NScaleEpisodes}{720}
\newcommand{\NScaleSizes}{4}
\newcommand{\NScaleForces}{3}
\newcommand{\NScaleSeeds}{30}
\newcommand{\ScaleSizeGrid}{2, 4, 6, 8}
\newcommand{\ScaleForceGrid}{0.10, 0.15, 0.20}
\newcommand{\ScaleMinSize}{2}
\newcommand{\ScaleMaxSize}{8}
\newcommand{\ScaleKneeMin}{0.61}
\newcommand{\ScaleKneeMax}{0.67}
\newcommand{\ScaleKneeTwo}{0.67}
\newcommand{\ScaleKneeFour}{0.63}
\newcommand{\ScaleKneeSix}{0.61}
\newcommand{\ScaleKneeEight}{0.62}
\newcommand{\ScaleKneeTwoLower}{0.55}
\newcommand{\ScaleKneeTwoUpper}{0.78}
\newcommand{\ScaleKneeFourLower}{0.52}
\newcommand{\ScaleKneeFourUpper}{0.75}
\newcommand{\ScaleKneeSixLower}{0.51}
\newcommand{\ScaleKneeSixUpper}{0.71}
\newcommand{\ScaleKneeEightLower}{0.53}
\newcommand{\ScaleKneeEightUpper}{0.72}
\newcommand{\ScaleKneePMax}{2.61e-8}
\newcommand{\ScaleKneePMin}{3.73e-9}
\newcommand{\ScaleKneeHolmMax}{2.35e-7}
\newcommand{\NRadiusEpisodes}{420}
\newcommand{\NRadiusSeeds}{10}
\newcommand{\NRadiusLevels}{3}
\newcommand{\RadiusGrid}{0.10, 0.15, 0.20}
\newcommand{\RadiusTightLevel}{0.10}
\newcommand{\RadiusWideLevel}{0.20}
\newcommand{\RadiusSpanM}{0.10}
\newcommand{\RadiusHolmFamily}{21}
\newcommand{\RadiusTightPeakForce}{0.10}
\newcommand{\RadiusTightPeakDelta}{0.55}
\newcommand{\RadiusAnchorPeakForce}{0.15}
\newcommand{\RadiusAnchorPeakDelta}{0.55}
\newcommand{\RadiusWidePeakForce}{0.20}
\newcommand{\RadiusWidePeakDelta}{0.42}
\newcommand{\RadiusPeakDeltaMin}{0.42}
\newcommand{\RadiusPeakDeltaMax}{0.55}
\newcommand{\RadiusPeakForceSpan}{0.10}
\newcommand{\RadiusHolmSurvivors}{1}
\newcommand{\RadiusAnchorHolmP}{0.0781}
\newcommand{\RadiusContinuousCloser}{16}
\newcommand{\RadiusContinuousCells}{21}
\newcommand{\RadiusContinuousHolmSurvivors}{2}
\newcommand{\NFactorialEpisodes}{270}
\newcommand{\NFactorialSeeds}{30}
\newcommand{\FactorialFeedforwardMin}{0.42}
\newcommand{\FactorialFeedforwardMax}{0.63}
\newcommand{\FactorialFeedforwardHolmMax}{1.28e-4}
\newcommand{\FactorialAssignmentAbsMax}{0.004}
\newcommand{\FactorialInteractionAbsMax}{0.008}
\newcommand{\FactorialOtherAbsMax}{0.008}
\newcommand{\NumFactorialCells}{27}
\newcommand{\NFactorialHolmSignificant}{3}
\newcommand{\FactorialContFeedforwardMin}{-0.034}
\newcommand{\FactorialContFeedforwardMax}{-0.010}
\newcommand{\FactorialContOtherAbsMax}{0.004}
\newcommand{\NSafetyEpisodes}{60}
\newcommand{\NSafetySeeds}{10}
\newcommand{\NSafetyMetrics}{6}
\newcommand{\SafetyForceGrid}{0.10, 0.15, 0.20}
\newcommand{\SafetyWorstSeparation}{0.049}
\newcommand{\SafetyWorstSeparationArm}{wind-aware}
\newcommand{\SafetyWorstSeparationForce}{0.20}
\newcommand{\SafetyAgnosticSeparationKnee}{0.100}
\newcommand{\SafetyAwareSeparationKnee}{0.085}
\newcommand{\SafetyMotorMax}{0.85}
\newcommand{\SafetyNonfiniteTotal}{0}
\newcommand{\SafetyStepMaxMs}{13.9}
\newcommand{\SafetyGeomTargetCorr}{0.76}
\newcommand{\SafetyGeomCompletionCorr}{-0.09}
\newcommand{\SafetyGeomHullContacts}{10}
\newcommand{\SafetyGeomHullContactM}{0.12}
\newcommand{\SafetyGeomCloseTargetContacts}{9}
\newcommand{\SafetyGeomFarTargetContacts}{1}
\newcommand{\SafetyGeomPairedCloser}{9}
\newcommand{\SafetyGeomPairedFarther}{2}
\newcommand{\SafetyGeomPairedTied}{19}
\newcommand{\SafetyGeomOverlapPrimary}{30.9}
\newcommand{\SafetyGeomOverlapMin}{11.5}
\newcommand{\SafetyGeomOverlapMax}{54.5}
\newcommand{\NEvidenceBlocks}{10}
\newcommand{\NEvidenceCampaigns}{11}
\newcommand{\NEvidenceRows}{11,012}
\newcommand{\StockKi}{0.05}
\newcommand{\IntegralTauS}{8}
\newcommand{\EpisodeHorizonS}{6}
\newcommand{\IntegralClampXy}{2}
\newcommand{\MaxSteadyIntegralForceN}{0.1}
\newcommand{\AuditStockMatch}{140}
\newcommand{\AuditKiTenGain}{+0.008}
\newcommand{\AuditKiTenWilcoxonP}{0.84}
\newcommand{\AuditKiTenSignP}{0.73}
\newcommand{\AuditKiTenCiLo}{-0.049}
\newcommand{\AuditKiTenCiHi}{+0.066}
\newcommand{\AuditKiTenBase}{0.79}
\newcommand{\AuditKiTenAware}{0.80}
\newcommand{\AuditKiTenKnee}{0.20}
\newcommand{\AuditPOneOverall}{mixed}
\newcommand{\AuditClampChanged}{0}
\newcommand{\AuditStockGain}{+0.54}
\newcommand{\AuditAOneBTwelveBase}{0.82}
\newcommand{\AuditClampMattersLongHorizon}{does}
\newcommand{\AuditGainTFour}{+0.04}
\newcommand{\AuditGainTSix}{+0.54}
\newcommand{\AuditGainTEight}{+0.14}
\newcommand{\AuditGainTTen}{-0.07}
\newcommand{\AuditGainTTwelve}{-0.07}
\newcommand{\AuditWidthTenNinety}{0.057}
\newcommand{\AuditWidthCiLo}{0.045}
\newcommand{\AuditWidthCiHi}{0.068}
\newcommand{\AuditNeverCloseBase}{0.53}
\newcommand{\AuditNeverCloseAware}{0.53}
\newcommand{\AuditGroundContactAfter}{0.87}
\newcommand{\AuditAltLossAfter}{0.94}
\newcommand{\AuditRadiusSlope}{0.68}
\newcommand{\AuditRadiusSlopeLo}{0.54}
\newcommand{\AuditRadiusSlopeHi}{0.82}
\newcommand{\AuditInitialDivergence}{0}
\newcommand{\AuditAwareEpisodes}{210}
\newcommand{\AuditUsedDivergence}{7}
\newcommand{\AuditUsedDivergencePrimary}{2}
\newcommand{\AuditUsedDivergenceKnee}{1}
\newcommand{\LambdaGrid}{0, 0.05, 0.1, 0.25, 0.5, 1, 2}
\newcommand{\LambdaGridSize}{7}
\newcommand{\LambdaGridMax}{2}
\newcommand{\LambdaReference}{0.25}
\newcommand{\LambdaInertMax}{0.5}
\newcommand{\LambdaFirstDivergence}{1.0}
\newcommand{\LambdaInitialDivAtOne}{3}
\newcommand{\LambdaInitialDivAtTwo}{19}
\newcommand{\LambdaUsedDivAtZero}{0}
\newcommand{\LambdaUsedDivAtQuarter}{7}
\newcommand{\LambdaUsedDivAtTwo}{80}
\newcommand{\LambdaAwareEpisodes}{210}
\newcommand{\LambdaEpisodes}{1,680}
\newcommand{\LambdaCompletionFlat}{0.63}
\newcommand{\AFiveBSlope}{0.705}
\newcommand{\AFiveBSlopeLo}{0.63}
\newcommand{\AFiveBSlopeHi}{0.78}
\newcommand{\AFiveBForceStepN}{0.01}
\newcommand{\AFiveBEpisodes}{8,820}

%% Draft-only owner placeholder (ruling R15, 2026-09-09). Prints ONLY when the
%% build defines \DRAFTTODO (e.g. pdflatex "\def\DRAFTTODO{1}\input{main}");
%% every packet build leaves it undefined, so no unconfirmed declaration text
%% reaches a PDF. The owner replaces each use with the confirmed sentence.
\providecommand{\ownertodo}[1]{\ifdefined\DRAFTTODO[TODO-OWNER: #1]\fi}

%% Independent caption file. Every caption is self-contained: readable without
%% the body text, without cross-references, and without float numbers.

\newcommand{\CapSystemOverview}{%
Research object and endpoint bridge for the wind-aware allocation--control
measurement. Panel A is a drawn schematic of the computational contract: four
matched quadrotors, four waypoint targets and a horizontal disturbance vector.
Dashed paths represent the shared
assignment-to-control task; the two labelled paths distinguish the Euclidean
wind-agnostic cost from the wind-aware upwind cost plus target-position bias.
Panel B is a descriptive seed-paired view at the transition force
\KneeLevel{}~N: each arrow connects the same seed under the two arms in mean
closest-approach versus completion-rate coordinates. The vertical reference is
the per-airframe capture radius used to define completion. Displayed arrows use
the paired seed record as the unit; the panel reports a descriptive endpoint view
with the seed-level inferential unit preserved. %
}

\newcommand{\CapBand}{%
Operating envelope of task completion against steady horizontal wind force for a
wind-agnostic and a wind-aware allocator flying the same \NDrones{} quadrotors on
the same \NSeeds{} seeds per force level, with both panels on one shared force
axis. Panel A is the completion ladder: points are means over seeds and bars are
one standard deviation, clipped at the zero and one bounds of the rate, and the
upper axis restates force as a percentage of the \HoverThrust{}~N hover thrust of
one airframe. Panel B is the seed-paired advantage of the wind-aware over the
wind-agnostic allocator, in completion-rate points, with bars giving the
95~percent $t$ interval over the \NSeeds{} paired seeds. Panel B carries the
quantity the significance test is computed on, and its interval clears zero at
\KneeLevel{}~N alone, where the paired gain is $+\KneeDelta{}$ at an exact
two-sided signed-rank $p$ of \KneeP{}. The two panels are drawn together because
the Panel A bars are dispersions of two marginal means and overlap at the
transition even though every seed moves in the same direction, so the paired
contrast is displayed rather than left to be inferred from that overlap. Shading
separates the near-complete shoulder at and below \SaturationLevel{}~N and the
high-force edge at and above \CollapseLevel{}~N from the unshaded transition,
whose \KneeLevel{}~N level the dashed line marks and where directional
information has the largest operational leverage. Panel A's red and blue legend
identifies the two allocation arms; Panel B is a single paired-difference series
and therefore has no second arm legend.%
}

\newcommand{\CapNoise}{%
Profile of the directional signal as wind-estimate quality changes. The horizontal
axis is relative magnitude noise redrawn independently at every control step. Panel
A shows the paired advantage in completion rate at the transition force
\KneeLevel{}~N, with bars giving the 95~percent $t$ interval over \NSeeds{} paired
seeds and labels giving the exact two-sided signed-rank $p$ value; the dashed line
marks the decision level fixed before execution. Panel B places the two neighbouring
force levels on the same vertical scale, showing that the operating transition stays
localized while estimate quality changes. A paired advantage of zero marks equal
completion for the two arms. In both panels, paired advantage means aware-minus-
agnostic completion rate for the same seed. %
}

\newcommand{\CapSpread}{%
Distribution behind the two completion means at the transition force
\KneeLevel{}~N. Each episode assigns \NDrones{} quadrotors to \NDrones{} waypoints,
so a single episode can score zero, one quarter, one half, three quarters or one.
The bars count how many of the \NSeeds{} seeds land on each value for the two arms;
the distribution displays the episode-level margin carried by the aggregate
operating-envelope gain. %
}

\newcommand{\CapHeadings}{%
Coverage of the force headings drawn by the \NSeeds{} seeds and the repeat groups
in the displayed record. Panel A places each seed's heading on the circle, one
spoke per seed; the heading is fixed by the seed and reused at every force level,
in both allocators and in the estimate-noise sweep. The draws resolve to
\NHeadings{} distinct headings, with \NRepeats{} near-matched pairs and a largest
unvisited arc of \UnsampledArc{} degrees. Panel B orders all
$\NSeeds{}(\NSeeds{}-1)/2$ pairwise separations on a logarithmic axis; the empty
gap between the retained group and the next separation makes the stated grouping
stable across the displayed tolerance interval. %
}

\newcommand{\CapRepeats}{%
Context spread under near-matched force headings. At each force level the bar gives
the largest completion difference between two episodes of the same allocator whose
headings differ by less than \RepeatTol{} degrees, expressed as the number of
airframes, out of \NDrones{}, by which the episodes differ. The strip below records
whether any seeds at that force level varied in completion; a zero bar at a fully
uniform level marks a saturated shoulder or floor, while the informative levels
show where waypoint context changes the outcome under nearly identical directional
input. %
}

\newcommand{\CapDesigned}{%
Heading-by-layout operating-envelope map from a crossed design. \NAssignedHeadings{}
headings \HeadingStep{} degrees apart are each flown with all \NScenes{} waypoint
layouts, at every force level and under both allocators, for \NEpisodesDesigned{}
episodes and \NPairsPerForce{} paired comparisons per level. Panel A shows paired advantage
against force, with the band giving an ordinary-bootstrap 95~percent interval over
the \NCrossedClusters{} layout-level means after averaging the fixed headings
(\NCrossedBootstrap{} deterministic resamples); crosses mark the \NSeeds{} seed
estimates from the independently drawn-heading condition.
Panel B opens the transition force into its heading-by-layout rectangle, one cell
per paired episode; rows are assigned headings and columns are scene layouts, and
colour is the paired completion advantage. The heat map displays the same cells
used to compute the heading, layout and interaction shares. %
}

\newcommand{\CapMiss}{%
Continuous tracking quantity behind the thresholded completion endpoint. Points give
the mean closest approach of each quadrotor to its assigned waypoint, averaged over
seeds, with bars of one standard deviation; the dashed line is the
\CaptureRadius{}~m capture radius that defines completion. Both allocators are
measured at the same forces, and the two series are offset horizontally by a fixed
0.0055~N purely so that their overlapping bars remain separable; the axis ticks
mark the true force levels. The two allocators occupy
nearby trajectories across the force ladder, and the measured distance shift at the
transition moves aircraft across the capture boundary to produce the task-level
operating gain. %
}

\newcommand{\CapTemporal}{%
Temporal operating-envelope extension for a causal directional estimate. Panel A
shows paired completion gain against true force for constant, slow-gust and fast-gust
profiles, with exact, causal and biased-causal estimate conditions and 95~percent
intervals over \NTemporalSeeds{} paired seeds. Panel B shows mean relative vector-tracking error
for the same profiles and estimate conditions. The transition remains concentrated
near \KneeLevel{}~N, while the tracking-error profile quantifies the cost of temporal
rate and links estimator dynamics to the measured task-level response. Colours and
symbols identify the three estimate conditions consistently across both panels. %
}

\newcommand{\CapMechanism}{%
Mechanism and gain-sensitivity extension at the transition force \KneeLevel{}~N.
Panel A gives paired completion gain across the tested target-feedforward gains for
each true-force profile; the dashed line marks the inherited gain
\FeedforwardGainDefault{}~m/N. Panel B compares the complete arm with an
allocation-cost-only ablation that retains directional assignment and selectively
removes the target-position bias. The separation identifies the pathway through
which the directional estimate is converted into task completion in this benchmark.
Intervals in both panels are paired 95~percent $t$ intervals over the declared
profile-specific seeds; the upper gain-grid point marks the measured edge of the
tested gain envelope. %
}

\newcommand{\CapGeometry}{%
Factorised scene-geometry extension of the directional operating envelope. Panel A
shows high-minus-low contrasts for target spread, target range and crossing
orientation at each tested force; intervals are 95~percent cluster-aware intervals
over 30 seed-level factor contrasts after averaging the four crossed settings
within seed; the 120-term version is sensitivity-only. Raw exact signed-rank
values and nine-contrast Holm adjustments define the support-family readout, and
the zero line marks no change in the wind-aware advantage. Panel B shows the mean
paired advantage in all eight scene conditions
across the three force levels, with each row defined by one spread, range and
crossing setting. The nominal--near--aligned cell is the legacy compatibility
condition; the remaining cells are deterministic scene transformations. Together,
the panels define the tested simulator geometry envelope and its named transfer
axes. %
}

\newcommand{\CapScale}{%
Homogeneous swarm-size extension of the directional operating envelope. Panel A
shows the paired completion gain against team size for two, four, six and eight
airframes at the three tested force levels; points are means over the same 30
paired seeds per cell and bars are 95~percent $t$ intervals. Raw exact signed-rank
values and their Holm adjustments are reported over the 12 size--force cells.
Panel B opens the transition force at \KneeLevel{}~N and shows the wind-agnostic and
wind-aware completion endpoints separately, with means and 95~percent intervals over
the same seed pairs. The constant-profile four-airframe slice contains 180 rows:
30 seeds, two arms and three force levels with the causal estimator. The ten-seed
exact-vector condition is a distinct oracle slice, and the other team sizes are
declared simulator cells. Together the panels define the measured homogeneous-size
transfer envelope. %
}

\newcommand{\CapRadius}{%
Sensitivity of the operating envelope to the capture radius that defines task
completion, over \NRadiusEpisodes{} episodes at radii of \RadiusGrid{}~m and
\NRadiusSeeds{} paired seeds per radius and force. Panel A gives the paired
completion advantage of the wind-aware over the wind-agnostic allocator against
steady horizontal wind force, one series per radius, with bars showing 95~percent
intervals over the paired seeds and the series offset horizontally so that
overlapping bars remain separable. The labelled point on each series is that
radius' largest advantage. Panel B gives the two completion ladders those
advantages are differences of, one facet per radius, with the dashed rule marking
the same labelled force. The transition survives every tested radius at a
comparable size, between $+\RadiusPeakDeltaMin{}$ and $+\RadiusPeakDeltaMax{}$
completion, and its location moves with the radius, from
\RadiusTightPeakForce{}~N at the tightest radius to \RadiusWidePeakForce{}~N at
the widest, so the force at which directional information pays is a property of
the operating point. Within this \RadiusHolmFamily{}-test validation family
\RadiusHolmSurvivors{} cell clears a Holm-corrected 5~percent threshold, a
resolution set by the \NRadiusSeeds{} seeds per cell; the family is a
sensitivity readout beside the larger primary sweep. %
}

\newcommand{\CapPathwaySafety}{%
Complete pathway decomposition and the descriptive safety diagnostics, from two
separate campaigns of \NFactorialEpisodes{} and \NSafetyEpisodes{} episodes.
Panel A gives all three terms of the assignment-by-feedforward factorial as
paired completion effects against true horizontal force, one facet per temporal
profile, with bars showing 95~percent intervals over \NFactorialSeeds{} paired
seeds and the terms offset horizontally for legibility. The target-feedforward
main effect carries the whole gain at the transition force, between
$+\FactorialFeedforwardMin{}$ and $+\FactorialFeedforwardMax{}$ completion across
profiles; the assignment main effect and the interaction stay within
\FactorialOtherAbsMax{} of zero and neither clears its Holm-corrected threshold at
any force. Panel B is descriptive only and carries no test: for each allocator and
force it gives the range of the per-episode minimum pairwise airframe separation
over \NSafetySeeds{} seeds, with the point at the mean of that range and the
dashed rule at the closest airframe pair recorded anywhere in the campaign,
\SafetyWorstSeparation{}~m. The panel is a benchmark-level feasibility diagnostic;
because target placement and arrival status contribute to the metric, it is not
interpreted as an allocator-causal clearance effect. %
}

\newcommand{\CapEvidence}{%
Cross-campaign evidence synthesis for the directional operating envelope. Panel A
places the transition paired-completion estimates from the primary, estimate-noise,
crossed-layout, causal-temporal and homogeneous-size views on a common rate axis;
bars are the view-specific 95~percent intervals and colours identify the evidence
block. The crossed-layout estimate is shown with its interval because that record
uses a layout-cluster bootstrap. Panel B shows the episode-row volume of each of the
eight views on a logarithmic axis, with the text next to each point naming the seed
or layout-cluster unit that supports its analysis; pathway ablation and gain
sensitivity are two views of one pathway/gain synthesis block. The figure is
descriptive by design and preserves each view's inferential unit. %
}

\newcommand{\CapIntegralGain}{%
Paired completion gain at the \KneeLevel{}~N force against the baseline integral
gain, from the controller-identification sweep at the stock
\EpisodeHorizonS{}~s horizon. Points are the seed-paired aware-minus-agnostic
completion difference; bars are 95~percent $t$ intervals. Clamp settings of
\IntegralClampXy{}, 4 and 8~m/s are drawn as distinct markers at the same Ki,
so a clamp that leaves the 0.15~N completions unchanged sits on top of the
stock marker. Doubling the stock Ki of \StockKi{}~N/m/s collapses the gain from
\AuditStockGain{} to \AuditKiTenGain{} (interval
[\AuditKiTenCiLo{}, \AuditKiTenCiHi{}]). %
}

\newcommand{\CapHorizon}{%
Horizon sweep under the stock controller (Ki~\StockKi{}~N/m/s, clamp
\IntegralClampXy{}~m/s). Panel A is the paired completion gain at
\KneeLevel{}~N against episode duration, with bars giving the 95~percent $t$
interval; the gain is \AuditGainTSix{} at the stock \EpisodeHorizonS{}~s
horizon and \AuditGainTTen{} once the horizon reaches 10~s. Panel B is the
matching wind-agnostic completion at the same force, which rises as the
episode lengthens and the slow integral is given more time to wind up. Both
panels share one horizon axis. %
}

\newcommand{\CapLambdaSweep}{%
Pre-declared sweep of the assignment penalty weight
$\lambda\in\{\LambdaGrid{}\}$~m/N under the stock controller at the
\EpisodeHorizonS{}~s horizon, \LambdaAwareEpisodes{} aware-arm episodes per
weight. Circles count episodes whose initial ($t=0$) assignment differs from
the Euclidean assignment; squares count episodes in which any call,
initial or replan, differs. The dashed line marks the inherited weight
$\lambda=\LambdaReference{}$, at which the counts reproduce the bound
assignment audit exactly. Initial divergence is zero through
$\lambda=\LambdaInertMax{}$ and first appears at
$\lambda=\LambdaFirstDivergence{}$; the horizontal axis is pseudo-logarithmic
so that $\lambda=0$ sits on the same axis. The sweep is descriptive by
declaration and carries no confirmatory test. %
}

%% Numbered environments, shared by every binding of this paper for the same
%% reason the body is: a statement that is Proposition 2 in one binding and
%% Proposition 3 in another cannot be cross-referenced by a reader holding the
%% other one.
%% Equation labels remain authored in methods.tex and are checked for uniqueness;
%% proofs use amsthm's automatic QED marker with no hand-added square.

\theoremstyle{plain}
\newtheorem{proposition}{Proposition}
\newtheorem{corollary}{Corollary}

\theoremstyle{definition}
\newtheorem{definition}{Definition}

\newcommand{\generatedtable}[1]{\input{#1}}

\title{From directional sensing to task completion: the operating envelope of
target-position feedforward against a slow integral in wind-disturbed quadrotor
waypoint capture}
\ifdefined\ANONYMOUS
\author{Author and affiliation withheld for blind review}
\else
\author{Zeyu Fu\\
Army Medical University, Chongqing, China\\
\texttt{fuzeyu99@126.com}\\
ORCID: 0009-0001-8329-0108}
\fi
\date{}

\begin{document}
\maketitle

\begin{abstract}
%% One paragraph, self-contained. No cross-references, section numbers,
%% structural labels or citations.
%% Every numerical value is emitted from the hash-bound evidence pass.

Directional disturbance information can change both which aircraft receives a
waypoint and how it flies there. We map and quantify a reproducible computational
operating envelope for this information-to-completion interface in a
multi-quadrotor allocation--control stack. The system is a Hungarian
allocation--control stack in a rigid-body simulator: a wind-agnostic arm uses
Euclidean assignment and the stock position controller, while a wind-aware arm
adds an upwind assignment cost and target-position feedforward. Shared seeds,
paired contrasts and a seven-level force ladder define the primary envelope, and
estimate-noise, temporal, pathway, gain, geometry, swarm-size, capture-radius
and feasibility campaigns extend the same endpoint. The envelope has a sharp
transition at \KneeLevel{}~N (\KneePctHover{}~percent of hover thrust), where
completion rises from \KneeAgnostic{} to \KneeAware{} (paired gain
+\KneeDelta{}, surviving the pre-specified signed-rank family correction). Relative
noise through \BestSigma{} preserves the transition, and causal estimates retain
positive gains under constant and gusting profiles. A complete
assignment-by-feedforward factorial attributes the gain to target feedforward,
while paired gains stay positive from \ScaleMinSize{} through \ScaleMaxSize{}
airframes. Rerunning the ladder across capture radii keeps a transition of
comparable size but moves it from \RadiusTightPeakForce{} to
\RadiusWidePeakForce{}~N, so force describes an operating point rather than an
interface constant. The study therefore contributes a simulator operating
envelope, a transparent computational validation contract, and a parameterized
transfer protocol for hardware and field validation; clearance and collision
quantities are retained as benchmark-level feasibility diagnostics rather than
as deployment claims.

\end{abstract}

\paragraph{Keywords} disturbance feedforward, integral control, wind estimation, thresholded endpoint, task allocation, quadrotor swarm, reproducibility

%% The manuscript body, shared by every binding of this paper.
%% The class, the front matter and the bibliography style belong to the
%% binding; everything a reader reads belongs here, once.
\section{Introduction}
A quadrotor that knows which way the wind is blowing can fly an assigned waypoint
with a more useful target bias. At the
vehicle level, directional disturbance information is now estimated, modelled and
used for control: unknown-input filters have been validated for wind-gust
compensation~\cite{azid2021windgust}, aerodynamic tracking laws have been tested
under wind perturbations~\cite{perozzi2018wind}, and onboard wind measurement and
simulation practices have been systematized~\cite{abichandani2020wind,wi2024windstress,
wetz2021distributed,schiano2014wind,gonzalezrocha2019sensingwind,
wang2018windestimation}.
Wind-aware planners likewise make direction and reachability explicit in their
costs~\cite{thanellas2019spatialwind,heidari2021windplanning,techy2009minimum};
recent work extends that line to flow prediction, safety and formation adaptation
~\cite{seal2025,wespr2026,knodedw2025}. A natural team-level hypothesis is that
the same directional information has operational leverage when several aircraft
compete for tasks, because it can enter an upwind assignment cost. This paper
puts that hypothesis to a paired test on a multi-quadrotor benchmark and
identifies where the directional information acts: at the inherited penalty
weight the initial assignment coincides with the Euclidean solution in every
audited episode, and the completion gain travels through target-position
feedforward compensating a slow stock integral.

The instance is a multi-quadrotor stack in which \NDrones{} airframes must reach
\NDrones{} waypoints. A wind-agnostic arm uses Euclidean assignment and the stock
position controller; a wind-aware arm can add an upwind assignment cost and
target-position feedforward. Shared seeds, paired contrasts and a
\NWindLevels{}-level force ladder define the primary envelope. The contribution
is a scoped identification of the mechanism: against the stock slow integral
(\StockKi{}~N/m/s) and the \EpisodeHorizonS{}~s horizon, target-position
feedforward changes short-episode completion, and the campaigns map the
operating envelope in which it does so.
Initial assignment is inert at the inherited $\lambda=0.25$ and, on a
pre-declared \LambdaGridSize{}-point sweep, at every tested weight through
$\lambda=\LambdaInertMax{}$~m/N; it first diverges at
$\lambda=\LambdaFirstDivergence{}$~m/N. Doubling that stock Ki collapses the 0.15~N paired advantage to
\AuditKiTenGain{}. A complete assignment-by-feedforward factorial attributes the
gain to the feedforward path. Allocation taxonomies and aerial-swarm
surveys~\cite{gerkey2004formal,dias2006market,korsah2013taxonomy,chung2018swarm,
skaltsis2023uavallocation} frame the team-level allocation question that the
benchmark puts to the test.
Panels A and B of Fig.~\ref{fig:system-overview} make the measured object
explicit: the same directional input enters assignment and feedforward, and the
paired endpoint keeps thresholded completion beside closest approach. Panel A
is a drawn schematic of the research object.

\begin{figure*}[!t]
\centering
\includegraphics[width=\linewidth]{fig0_system_overview.pdf}
\caption{\CapSystemOverview}
\label{fig:system-overview}
\end{figure*}

The study follows the information path from capability to mechanism. It first
locates the transition and tests estimate-noise robustness, then resolves scene
context with seed-level repeats and a crossed heading--layout design. A continuous
tracking endpoint shows how the thresholded gain is formed. Temporal profiles,
pathway and feedforward-gain controls, scene-geometry factors and a homogeneous
team-size axis then open the same estimand along progressively more operational
dimensions. The exact-vector oracle is the reference for every extension, so the
\NEvidenceBlocks{} synthesis blocks form one measured operating envelope with a
continuous information-to-completion trace. The result is a computational
baseline: the tested simulator envelope is the calibration reference against
which a hardware extension is to be validated.

\section{Methods}
\input{methods}
\input{transfer_protocol}

\section{Results}

\subsection{Operating envelope}

Table~\ref{tab:band} and Fig.~\ref{fig:band}A give the full sweep. The continuous
closest-approach companion is reported below with the tracking endpoint and is
read alongside the thresholded rate.
At \KneeLevel{}~N, which is \KneePctHover{} percent of one airframe's hover thrust, the
wind-agnostic allocator completes \KneeAgnostic{} of waypoints on average and the
wind-aware allocator completes \KneeAware{}. The paired difference is
$+\KneeDelta{}$ with an exact two-sided signed-rank $p$ of \KneeP{}, which clears
the Bonferroni-corrected threshold of \BonferroniAlpha{} for \NWindLevels{} levels.
The two marginal standard deviations overlap at that level, so the ladder on its
own understates the separation: every seed moves in the same direction even though
the arms' dispersions do not separate. Fig.~\ref{fig:band}B therefore plots the
seed-paired difference the test is actually computed on, and its 95~percent
interval excludes zero at \KneeLevel{}~N and at no other force in the sweep.
That one-level reading is scoped to the stock integral gain \StockKi{}~N/m/s,
the \EpisodeHorizonS{}~s horizon, the \CaptureRadius{}~m capture radius and the
declared 0.05~N force grid; the controller-identification audit shows how the
same endpoint moves when those settings change.
This force is specific to the \CaptureRadius{}~m capture radius that defines
completion here; rerunning the ladder at two further radii finds a transition
of comparable size at a different force in each, so \KneeLevel{}~N is the
transition of this operating point and moves with the capture radius.

The full sweep gives the envelope around that transition. At and below
\SaturationLevel{}~N both allocators complete nearly every waypoint, establishing a
saturated shoulder; at and above \CollapseLevel{}~N both arms occupy the measured
high-force edge. The four all-zero levels and the two levels with $p=1$ provide
explicit convergence boundaries for the endpoint, alongside the transition where
the arms separate. Taken together, the shoulder, transition and high-force edge
map where directional information has operational leverage; the complete ladder
makes each regime visible.

\begin{table}[tbp]
\centering
\caption{Completion rate by wind force for the wind-agnostic and wind-aware arms,
over \NSeeds{} paired seeds per level. Entries are mean $\pm$ standard deviation;
intervals and exact signed-rank $p$ values use seed-level paired differences. A dash
identifies a saturated all-zero level where no signed-rank statistic is defined.}
\label{tab:band}
\generatedtable{generated_table_band}
\end{table}

\begin{figure}[tbp]
\centering
\includegraphics[width=\linewidth]{fig1_wind_band.pdf}
\caption{\CapBand}
\label{fig:band}
\end{figure}

\subsection{Estimate-noise robustness}

Table~\ref{tab:noise}, Fig.~\ref{fig:noise}A and Fig.~\ref{fig:noise}B report the
controlled estimate perturbation that motivates this robustness analysis.
Panel A tracks the transition force, while Panel B verifies its neighbouring
force levels. Introducing controlled perturbations to the wind vector supplied to
the wind-aware allocator preserves the band. The paired advantage is
$+\SigmaZeroDelta{}$ with an exact input, and remains statistically aligned with
that value across relative noise of \NearestSigma{} and \BestSigma{}, where
the point estimates are $+\NearestSigmaDelta{}$ and $+\BestSigmaDelta{}$.

At the largest tested relative noise, \MaxSigma{}, where the corruption is as large
as the signal and the heading error has a standard deviation of roughly a radian,
the measured advantage reaches its robustness edge at $+\MaxSigmaDelta{}$ with an
exact $p$ of \MaxSigmaP{}. Every 95~percent interval in the table excludes zero, the
weakest lower bound being \WeakestLower{} at that edge.

The decision point fixed before the runs was relative noise of \PreregSigma{};
the nearest executed grid level is \NearestSigma{}, so we evaluate the pre-set
decision on that grid-aligned condition, where the interval is bounded away from
zero. Counting the complete disturbance-by-noise family gives a strict
Bonferroni threshold of \NoiseBonferroni{}. The 0.30 and 0.50 relative-noise
conditions clear that family threshold; the 0.10 condition is retained as a
positive descriptive estimate below that threshold, and the 1.00 condition
($p=\MaxSigmaP{}$) marks the measured edge without a multiplicity-adjusted
significance claim. Read together, the curve is a profile of how directional information translates
into allocation benefit as estimate quality changes.
The neighbouring disturbance levels stay flat at every noise level, which answers a
question the point estimates alone would leave open: degrading the estimate leaves
the transition location stable and changes the advantage within the measured band.
That stability is a useful calibration result for selecting the next estimator
quality regime.

\begin{table}[tbp]
\centering
\caption{Effect of relative wind-estimate noise on the paired advantage at
\KneeLevel{}~N over \NSeeds{} paired seeds. The reference arm is deterministic given
a seed, so its recorded outcome is reused exactly for each noise level. Intervals are 95~percent $t$ intervals on seed-level differences;
$p$ values are exact two-sided signed-rank tests.}
\label{tab:noise}
\generatedtable{generated_table_noise}
\end{table}

\begin{figure}[tbp]
\centering
\includegraphics[width=\linewidth]{fig2_noise_sensitivity.pdf}
\caption{\CapNoise}
\label{fig:noise}
\end{figure}

\subsection{Episode-level distribution}

Because an episode assigns \NDrones{} quadrotors to \NDrones{} waypoints, its
completion rate takes one of five values, and a mean over \NSeeds{} seeds can hide
very different distributions. Fig.~\ref{fig:spread} opens the two means at the
discriminating level. The wind-aware mean of \KneeAware{} is concentrated rather
than diffuse: \AwareTopSeeds{} of \NSeeds{} seeds finish three of the four
waypoints, while \AwareZeroSeeds{} seed finishes none. The wind-agnostic mean of
\KneeAgnostic{} is similarly concentrated, with \AgnosticZeroSeeds{} of
\NSeeds{} seeds finishing nothing.

This distribution is part of the contribution. The gain is a near-categorical
change in outcome for most seeds, with a smaller set of matched scenes producing
the same endpoint under both arms. Reporting the mean together with this
distribution makes the operating envelope useful for aggregate planning and
episode-level expectations.

\begin{figure}[tbp]
\centering
\includegraphics[width=\linewidth]{fig3_seed_spread.pdf}
\caption{\CapSpread}
\label{fig:spread}
\end{figure}

\subsection{Continuous tracking endpoint}

The completion rate is a step function of a continuous quantity: an aircraft counts
as finished when its closest approach falls inside the \CaptureRadius{}~m capture
radius. Fig.~\ref{fig:miss} plots that continuous quantity directly and shows
how the operational gain is formed. At the discriminating level
the mean closest approach is \MissAgnosticKnee{}~m for the wind-agnostic allocator
and \MissAwareKnee{}~m for the wind-aware one, a difference of under two
centimetres, while the completion rates at the same level differ by
$\KneeDelta{}$. At the largest disturbance the two arms are closer still,
\MissAgnosticCollapse{}~m against \MissAwareCollapse{}~m, and neither completes
anything.

The wind-aware allocator flies slightly more accurately at the disturbance where
that small improvement moves most aircraft across the capture boundary. This is the
operational leverage of the method: target-position feedforward turns a modest
continuous tracking shift into a substantial task-level change at the transition.
The continuous endpoint makes that mechanism explicit and links the completion
jump to its physical displacement. This yields a general measurement principle
that extends beyond wind: when an endpoint thresholds a continuous quantity, the
reported effect is determined by where the boundary sits relative to the two
distributions as well as by the intervention. Pairing the rate with its continuous
companion therefore shows both the operational payoff and the physical scale of
the tracking change.

The object-level bridge in Fig.~\ref{fig:system-overview}B makes that relationship
visible at the seed level. At the transition, \BridgeLowerPairs{} of the
\NSeeds{} matched seeds move to a lower mean closest approach under the aware
arm, with a mean shift of $\BridgeDistanceShift{}$~m and a completion gain of
$+\BridgeCompletionGain{}$ rate units. This bridge is a descriptive reuse of
the primary record; the inferential unit and the single paired test are
unchanged.

\begin{figure}[tbp]
\centering
\includegraphics[width=\linewidth]{fig4_miss_distance.pdf}
\caption{\CapMiss}
\label{fig:miss}
\end{figure}

\subsection{Scene-context variation}

The episode-level distribution leaves an important systems question. If the
wind-aware mean is a mixture over seeds, which part of the episode context
carries that spread? A seed fixes both the force heading and the waypoint layout, so near
repetitions provide a direct probe of how the operating envelope depends on
context. The record supports a joint reading: the force level locates the
transition, while the scene determines how strongly an individual episode
expresses it.
Table~\ref{tab:seedpairs} prints the discriminating cell one seed at a time.
The heading and replanning columns expose context that is absent from aggregate
summaries.
The heading is drawn once per seed and reused at every force level, in both arms
and in the second sweep; near-matched seeds therefore repeat one disturbance while
retaining different layouts. Fig.~\ref{fig:headings}A shows the \NSeeds{} draws
and Fig.~\ref{fig:headings}B their complete separation spectrum. \NRepeats{} pairs
lie within \WidestRepeat{} degrees (the closest is \ClosestRepeat{}), resolving to
\NHeadings{} distinct headings with \LargestHeadingGroup{} inside one
\HeadingGroupSpan{}-degree arc. The next pair is \NearestDistinct{} degrees away,
so Proposition~\ref{prop:tol} gives the same repeat set throughout the displayed
tolerance interval.

Table~\ref{tab:repeats} and Fig.~\ref{fig:repeats} report how those
repetitions spread. At the discriminating level, \RepeatDisagreeKnee{}
of the \NRepeats{} repeated disturbances disagree, and the largest disagreement
is \RepeatGapKneeDrones{} of the \NDrones{} airframes. The two seeds whose
headings differ by \ClosestRepeat{} degrees return different outcomes, and one
pair spans the entire range the wind-aware arm ever reached. By Proposition~\ref{prop:repeat} the completion
rate at \KneeLevel{}~N is jointly shaped by the disturbance and the scene
context, and the gap of $\RepeatGapKnee{}$ quantifies the context-dependent
variation that an operating-envelope design can accommodate.
The comparison level matters as much as the number. \NDegenerateLevels{} of the
\NWindLevels{} levels are constant across seeds; Proposition~\ref{prop:vacuous}
therefore identifies their repeat gaps as saturated ceiling or floor calibration
values. The remaining \NInformativeLevels{} levels carry the scene-variation signal:
at \CollapseLevel{}~N the seed spread is \CollapseSeedSpread{} and
\RepeatDisagreeCollapse{} repeat pairs disagree, at \SaturationLevel{}~N
\RepeatDisagreeShoulder{} disagree by \RepeatGapShoulderDrones{} airframe, and at
\KneeLevel{}~N \RepeatDisagreeKnee{} disagree by up to
\RepeatGapKneeDrones{}. Context dependence is consequently concentrated at the
transition where the arms can be distinguished.
Replanning is logged when an allocator changes its assignment mid-episode. At the
discriminating level the wind-agnostic and wind-aware arms replanned
\ReplanKneeAgnostic{} and \ReplanKneeAware{} times, respectively; the
\NReplanLoud{} seeds with at least one event carry $+\ReplanLoudDelta{}$ versus
$+\ReplanQuietDelta{}$ for the \NReplanQuiet{} quiet seeds. With \NSeeds{} seeds
this remains a descriptive location of variance, reported beside the primary test.

Corollary~\ref{cor:notn} clarifies the inferential unit and its operational
reading. The test's units are the \NSeeds{} per-seed differences;
two seeds sharing a heading still differ in their waypoint layouts, and the
advantage of $+\KneeDelta{}$ at an exact $p$ of \KneeP{} stands as reported.
The repeat profile shows that the same force heading can differ by
\RepeatGapKneeDrones{} of \NDrones{} airframes across layouts. The paired mean
locates the opportunity; the episode-level spread supplies a concrete margin for
the next scene-aware deployment study.

\begin{table}[tbp]
\centering
\caption{Seed-level paired record at \KneeLevel{}~N. Each row uses the same force
and waypoint layout in both arms. Heading is drawn once per seed and reused across
force levels; completion is the number of finished airframes, and replanning is the
number of mid-episode assignment changes.}
\label{tab:seedpairs}
\generatedtable{generated_table_seedpairs}
\end{table}

\begin{table}[tbp]
\centering
\caption{Near-matched disturbance repeats. Each row pairs seeds whose headings
differ by less than \RepeatTol{} degrees; entries are the larger within-arm
completion difference in airframes. The final row records whether any seeds differ
at that force, distinguishing informative gaps from saturated levels.}
\label{tab:repeats}
\small
\generatedtable{generated_table_repeats}
\end{table}

\begin{figure}[tbp]
\centering
\includegraphics[width=\linewidth]{fig5_heading_coverage.pdf}
\caption{\CapHeadings}
\label{fig:headings}
\end{figure}

\begin{figure}[tbp]
\centering
\includegraphics[width=\linewidth]{fig6_replicate_gap.pdf}
\caption{\CapRepeats}
\label{fig:repeats}
\end{figure}

\subsection{Heading--layout sweep}

The original sampling record supplies near-matched repetitions. We extend that
record with a designed crossed sweep that resolves the remaining sampling
question: headings are fixed on a grid and every heading is paired with every
waypoint layout. The simulator is deterministic, single-core and about a second
per episode, so this direct coverage test is inexpensive and fully reproducible.
Before comparing the two records we re-ran the first. All \NReproEpisodes{}
episodes of the drawn-heading sweep return the same value in all
\NReproFields{} recorded fields, confirming identical reproducibility of the
baseline results. The crossed record below therefore extends the same validated
endpoint.
The crossed sweep uses a fixed heading grid:
\NAssignedHeadings{} headings \HeadingStep{} degrees apart, each flown with all
\NScenes{} waypoint layouts, at every force level and in both arms
(Definition~\ref{def:crossed}). It leaves \DesignedArc{} degrees unsampled
against \UnsampledArc{}, and it puts \NPairsPerForce{} paired comparisons at
each level against \NSeeds{}. Table~\ref{tab:designed} and
Fig.~\ref{fig:designed}A report the force-level envelope, while
Fig.~\ref{fig:designed}B opens its transition cell.

The operating envelope is reproduced, and the transition remains one level wide.
At \KneeLevel{}~N the paired
advantage is $+\DesignedKneeDelta{}$ with a 95~percent ordinary-bootstrap
interval over the \NCrossedClusters{} layout clusters,
$[+\DesignedKneeLower{}, +\DesignedKneeUpper{}]$. Each layout cluster first
averages the \NAssignedHeadings{} fixed headings; the \NPairsPerForce{} cells
support the design and visualization, while the layout clusters provide the
inferential units. The estimate is from
\DesignedKneeAgnostic{} completed under the wind-agnostic allocator to
\DesignedKneeAware{} under the wind-aware one; the ten-draw estimate of
$+\KneeDelta{}$ sits inside that interval. Exactly \DesignedKneeAdverse{} of the
\NPairsPerForce{} pairs at that level favours the wind-agnostic arm. The
neighbouring levels stay where they were, $+\DesignedShoulderDelta{}$ at
\SaturationLevel{}~N and $+\DesignedCollapseDelta{}$ at \CollapseLevel{}~N, so
the designed sweep confirms the narrowness of the band.
The context dependence is visible at exact heading (zero tolerance). At
\KneeLevel{}~N two episodes at the \emph{same} heading, differing only in
waypoint layout, span $\ZeroSeparationGap{}$ of paired advantage, from
$\ZeroSeparationLow{}$ to $\ZeroSeparationHigh{}$. By
Corollary~\ref{cor:zerotol}, exact-heading repeats establish a scene component
that is separate from the arm and the wind, without a smoothness assumption.
The crossed design therefore identifies the residual scene context directly:
the same directional input expresses different completion outcomes across
layouts, which makes layout a measurable operating axis.
This is a complementary readout of the primary endpoint; Corollary~\ref{cor:zerotol}
keeps the repeated-heading statement tied to the declared design without altering
the force-family test.

Panel B of Fig.~\ref{fig:designed} opens the discriminating level into
its heading-by-layout rectangle, one cell per paired episode. Flat rows would
indicate a heading-only pattern and flat columns a layout-only pattern; the
observed matrix resolves both margins and their interaction. Of the variation in
the advantage at that level, \DesignedShareHeading{} per cent lies in the heading
margin, \DesignedShareScene{} per cent in the layout margin, and
\DesignedShareResidual{} per cent in the interaction. By
Proposition~\ref{prop:crossed}, the interaction is the dominant component of the
measured variation: wind and layout jointly shape the transition, while each
margin contributes a smaller, interpretable share.
Read with the paired test, the two statements are complementary. Over
\NPairsPerForce{} pairs the wind-aware allocator finishes
\DesignedKneeDelta{} more of \NDrones{} aircraft on average at \KneeLevel{}~N,
with a tighter estimate than the original ten draws. Exact wind knowledge still
leaves a scene-dependent spread on an individual flight; the mean locates the
opportunity and the spread supplies the deployment margin.

\begin{table}[tbp]
\centering
\caption{Crossed heading--layout sweep with \NAssignedHeadings{} fixed headings
spaced \HeadingStep{} degrees apart and \NScenes{} layouts per heading. The
reported interval is an ordinary bootstrap over \NCrossedClusters{} layout-level
means after averaging the fixed headings (\NCrossedBootstrap{} draws); the cells are
summarized through their layout-cluster units. The repeat-gap column is the largest
same-heading difference; the final columns give the heading, layout and interaction
shares computed from the displayed crossed cells.}
\label{tab:designed}
\small
\generatedtable{generated_table_designed}
\end{table}

\begin{figure}[tbp]
\centering
\includegraphics[width=\linewidth]{fig7_designed_replicates.pdf}
\caption{\CapDesigned}
\label{fig:designed}
\end{figure}

\subsection{Temporal operating profiles}

The constant-force result establishes where directional information has leverage;
the next question is whether that leverage survives when the vector itself moves.
We therefore ran a primary campaign across \NTemporalProfiles{} profiles,
\NTemporalEstimators{} estimate conditions and three force levels, with the same
paired scene and phase for every arm. The campaign contains \NTemporalEpisodes{}
episode records. Fig.~\ref{fig:temporal}A and Fig.~\ref{fig:temporal}B, together
with Table~\ref{tab:temporal}, report the causal-estimate slice at the transition
alongside the oracle and biased-causal conditions.

At \KneeLevel{}~N the causal estimate gives a paired gain of
$+\TemporalConstantDelta{}$ for the constant profile, with a 95~percent interval
$[+\TemporalConstantLower{},+\TemporalConstantUpper{}]$. The slow gust gives
$+\TemporalSlowDelta{}$ with interval
$[+\TemporalSlowLower{},+\TemporalSlowUpper{}]$, and the fast gust gives
$+\TemporalFastDelta{}$ with interval
$[+\TemporalFastLower{},+\TemporalFastUpper{}]$. The exact two-sided signed-rank
$p$ values for these three rows are printed in Table~\ref{tab:temporal}.
The transition is therefore carried into both temporal regimes: its expression
changes with the rate of the gust, while its location remains a force-scale
phenomenon expressed across dynamic profiles.

The observer diagnostics make that statement operational. Mean relative vector
tracking error is \TemporalTrackingConstant{}~percent for the constant profile,
\TemporalTrackingSlow{}~percent for the slow gust and \TemporalTrackingFast{}~percent
for the fast gust. The neighbouring force cells move by no more than
\TemporalShoulderMax{} in paired gain for the causal condition, so the temporal
extension preserves the narrow shoulder--transition--edge structure of the
original envelope. The \NTemporalLegacyRows{} constant/oracle regression rows
reproduce the legacy control path, tying the new dynamic readout to the established
reference while preserving the original endpoint.

\begin{table}[tbp]
\centering
\caption{Temporal operating-envelope results at the transition force. Each row is a
profile--estimator condition; intervals use the \NTemporalSeeds{} paired seeds,
tracking is mean relative vector error, and $p$ is an exact signed-rank test against
the reference arm. The oracle uses the exact vector, the causal condition uses a
first-order update, and the biased-causal condition adds the stated bias and noise.}
\label{tab:temporal}
\small
\generatedtable{generated_table_temporal}
\end{table}

\begin{figure}[tbp]
\centering
\includegraphics[width=\linewidth]{fig8_temporal_envelope.pdf}
\caption{\CapTemporal}
\label{fig:temporal}
\end{figure}

\subsection{Pathway and gain}

The complete wind-aware arm combines an upwind assignment cost with a
target-position feedforward term. A feedforward-off support campaign retains the
directional assignment cost and selectively disables the target-position bias, producing
\NAblationEpisodes{} paired records under the same temporal profiles, force levels,
seeds and estimate states. At the transition, the allocation-only gains are
$+\AllocOnlyConstant{}$, $+\AllocOnlySlow{}$ and $+\AllocOnlyFast{}$ for the
constant, slow-gust and fast-gust profiles; the corresponding complete-arm gains
are $+\CompleteConstant{}$, $+\CompleteSlow{}$ and $+\CompleteFast{}$. In this
benchmark the comparison localizes the observed task-level change to the
target-position pathway while preserving the directional assignment component,
thereby identifying a concrete tuning lever for the interface; the assignment
cost remains an explicit component to test across the next scene geometries.

The controller setting is then opened as an explicit factor. A
\NGainEpisodes{}-episode sensitivity campaign varies the feedforward gain from zero
to 1.40~m/N on the fixed seven-point grid, with every condition paired to the same
scene and causal estimate. Around the inherited \FeedforwardGainDefault{}~m/N
setting, the tested profiles occupy a broad positive plateau from
\GainPlateauMin{} to \GainPlateauMax{} in paired completion gain. At the upper
tested edge the gains are $+\GainEdgeConstant{}$, $+\GainEdgeSlow{}$ and
$+\GainEdgeFast{}$ for the three profiles. Fig.~\ref{fig:mechanism}A and
Fig.~\ref{fig:mechanism}B, together with Table~\ref{tab:mechanism}, make the
design choice visible: the inherited setting is inside a measured operating
region, while the upper tested edge marks the boundary of the gain grid.

\begin{table}[tbp]
\centering
\caption{Pathway and feedforward-gain sensitivity at the transition force. The
allocation-only arm retains the upwind assignment cost but disables target
feedforward; the complete arm uses the inherited gain. Plateau and upper-edge
columns summarize the tested gain grid. Entries are paired completion gains.}
\label{tab:mechanism}
\small
\generatedtable{generated_table_mechanism}
\end{table}

\begin{figure}[tbp]
\centering
\includegraphics[width=\linewidth]{fig9_mechanism_gain.pdf}
\caption{\CapMechanism}
\label{fig:mechanism}
\end{figure}

\subsection{Scene-geometry factors}

The crossed heading--layout result identifies scene context as a major source of
episode variation. We therefore opened that context with a targeted factorised
campaign that resolves ``layout'' into named factors. The
campaign contains \NGeometryEpisodes{} episodes across eight combinations of
target spread, target range and crossing orientation, three force levels and
\NGeometrySeeds{} paired seeds. The nominal--near--aligned cell is the exact legacy
scene, so the extension has a direct compatibility anchor.
Table~\ref{tab:geometry} lists the nine force-by-factor contrasts and their
declared paired units. Fig.~\ref{fig:geometry}A gives
the high-minus-low factor contrasts, each formed by first averaging the four
crossed settings of the other two factors within each of the \NGeometryClusters{}
seed clusters; the inferential vector therefore has \NGeometryClusters{} terms.
The \NGeometryTermSensitivity{} seed-by-setting terms serve as a sensitivity
readout, while the primary support-unit count remains the seed-cluster vector.
Fig.~\ref{fig:geometry}B
shows all condition-by-force means. At the transition force, crossing the target
geometry gives the largest contrast, $+\GeometryKneeCrossing{}$
$[+\GeometryKneeCrossingLower{},+\GeometryKneeCrossingUpper{}]$ with an exact
two-sided signed-rank $p=\GeometryKneeCrossingP{}$ (Holm-adjusted
$p=\GeometryKneeCrossingHolm{}$ within the nine-contrast support family).
Expanding target spread shifts the gain by
$\GeometryKneeSpread{}$ $[\GeometryKneeSpreadLower{},\GeometryKneeSpreadUpper{}]$,
while moving the target range farther shifts it by
$\GeometryKneeRange{}$ $[\GeometryKneeRangeLower{},+\GeometryKneeRangeUpper{}]$.
The crossing contrast is therefore a concrete scene-design lever in this
benchmark, whereas the spread and range contrasts quantify the direction and
precision of two neighbouring changes. The range estimate is a precision
boundary rather than an equivalence claim: its 30-cluster interval reaches the
zero boundary at the displayed resolution, while the 120-term sensitivity
interval is $[\GeometryRangeSensitivityLower{},+\GeometryRangeSensitivityUpper{}]$.
Using the declared nine-contrast family, the planning calculation gives an 80\%
familywise MDE of $\GeometryRangeMDE{}$ completion-rate units for this range
contrast.

These nine factor-by-force contrasts form one declared geometry-support family.
The high-minus-low terms are paired within seed and crossed scene condition, then
averaged within seed before the exact signed-rank test and interval are computed.
The \NGeometryTermSensitivity{} terms at a force level provide a design-level
sensitivity analysis, while the \NGeometryEpisodes{} episode rows document the
full crossed record and its audit trail.

The condition panel keeps the interpretation anchored to the endpoint: the
nominal legacy cell gives a transition gain of +\GeometryKneeNominal{}, and the
strongest tested cell reaches +\GeometryKneeBest{}. The factor effects define a
controlled operating-envelope map within the declared simulator and provide a
repeatable bridge from ``scene matters'' to an experimental program in which target
geometry can be selected, measured and extended along named axes.

\begin{table}[tbp]
\centering
\caption{Factorised scene-geometry campaign. Each row is a high-minus-low contrast
for one factor and force, with the other two factors crossed. Intervals and exact
signed-rank tests use \NGeometryClusters{} seed-level contrasts after averaging
the four crossed settings within seed; the \NGeometryTermSensitivity{}-term
version is a sensitivity readout. Holm adjustment is across the nine completion
contrasts in the support family. Levels are deterministic simulator
transformations; nominal--near--aligned is the legacy compatibility cell.}
\label{tab:geometry}
\small
\generatedtable{generated_table_geometry}
\end{table}

\begin{figure}[tbp]
\centering
\includegraphics[width=\linewidth]{fig10_geometry_factorial.pdf}
\caption{\CapGeometry}
\label{fig:geometry}
\end{figure}

\subsection{Team-size transfer}

The final extension asks whether the transition is retained when the homogeneous
team grows. The scale campaign contributes \NScaleEpisodes{} episode rows across
\NScaleSizes{} team sizes, \NScaleForces{} force levels and \NScaleSeeds{} paired
seeds, with the causal directional estimate and the same endpoint used above.
Fig.~\ref{fig:scale}A shows the complete size--force surface: the transition
remains concentrated at \KneeLevel{}~N while the neighbouring cells stay close to
the envelope baseline. This isolates transfer of the main interface across team
size with the controller and simulator held fixed.

At the transition, the paired gains are $+\ScaleKneeTwo{}$,
$+\ScaleKneeFour{}$, $+\ScaleKneeSix{}$ and $+\ScaleKneeEight{}$ for the
declared sizes from \ScaleMinSize{} to \ScaleMaxSize{} airframes. Every
95~percent transition-cell interval in Table~\ref{tab:scale} remains positive;
the four-airframe slice is $+\ScaleKneeFour{}$
$[+\ScaleKneeFourLower{},+\ScaleKneeFourUpper{}]$, and all exact tests are at
or below the Holm-adjusted bound $\ScaleKneeHolmMax{}$. The transition gain therefore spans a compact
$+\ScaleKneeMin{}$--$+\ScaleKneeMax{}$ range across the declared sizes.

Panel B of Fig.~\ref{fig:scale} separates the endpoint means at the transition.
The wind-aware arm remains above the reference at every tested size. The
four-airframe constant-profile causal slice (30 paired seeds, both arms and three
force levels; 180 rows) is field-by-field identical to the temporal record and is
distinct from the ten-seed oracle values in Table~\ref{tab:band}. The full table
retains the shoulder and high-force cells that delimit this homogeneous-team
transfer claim.

\begin{table}[tbp]
\centering
\caption{Homogeneous swarm-size operating-envelope campaign. Each row is a size--force
cell over \NScaleSeeds{} paired seeds. Entries are mean completion $\pm$ standard
deviation; gain intervals are 95~percent $t$ intervals. The final column gives the
raw exact signed-rank $p$ followed by the Holm-adjusted $p$ across the 12 size--force
cells. Four-airframe rows are the compatibility slice; other sizes are declared
extensions.}
\label{tab:scale}
\small
\generatedtable{generated_table_scale}
\end{table}

\begin{figure}[tbp]
\centering
\includegraphics[width=\linewidth]{fig11_swarm_scale.pdf}
\caption{\CapScale}
\label{fig:scale}
\end{figure}

\subsection{Endpoint-radius sensitivity}

Completion is scored by a capture radius, so every rate in this paper depends on a
distance the experiment chose. The validation sweep reruns the
force ladder at radii of \RadiusGrid{}~m over \NRadiusEpisodes{} episodes, holding
scene, heading, seed pairing and horizon fixed, and the anchor column reproduces
the primary ladder exactly. Table~\ref{tab:radius} gives the full grid and
Fig.~\ref{fig:radius}A the paired advantage at each radius. The record also
carries the underlying continuous distance as its own corrected family on the
same rows, and it moves the same way: the wind-aware arm ends closer to its
target in \RadiusContinuousCloser{} of the \RadiusContinuousCells{} radius-force
cells, with \RadiusContinuousHolmSurvivors{} clearing Holm. The threshold
therefore reports a discretization of a distance difference that the continuous
endpoint records independently of where the threshold is placed.

The result separates two claims that the primary sweep states together. A
transition exists at every tested radius and its size is comparable across them,
from $+\RadiusPeakDeltaMin{}$ to $+\RadiusPeakDeltaMax{}$ completion. Its
location moves with the radius: the largest advantage sits at
\RadiusTightPeakForce{}~N for the \RadiusTightLevel{}~m radius,
\RadiusAnchorPeakForce{}~N for the \CaptureRadius{}~m anchor and
\RadiusWidePeakForce{}~N for the \RadiusWideLevel{}~m radius, a span of
\RadiusPeakForceSpan{}~N in force across a \RadiusSpanM{}~m change in radius.
Fig.~\ref{fig:radius}B shows the mechanism directly, since a tighter radius
makes the same disturbance harder and shifts the whole completion ladder toward
lower forces. The headline force is therefore a property of the operating
point: directional information has its largest leverage at the force where the
task is on the edge of feasibility, and that force moves with the capture
radius.

The power of the family sets the scope of that reading. Within this predeclared
\RadiusHolmFamily{}-test validation family \RadiusHolmSurvivors{} of the cells
clears a Holm-corrected 5~percent threshold, and the anchor cell reaches
\RadiusAnchorHolmP{}; with \NRadiusSeeds{} seeds per cell the family locates
the transition at each radius descriptively and is reported as a sensitivity
readout rather than a second confirmation. The primary seven-force result
retains its own Bonferroni family and is unchanged.

\begin{table}[tbp]
\centering
\caption{Completion by capture radius and wind force for both allocators, over
\NRadiusSeeds{} paired seeds per cell. ``Paired diff.'' is the wind-aware minus
wind-agnostic seed-level mean with its 95~percent interval, and ``Holm $p$'' is
corrected within the \RadiusHolmFamily{}-test endpoint-radius validation family.
Cells with an undefined exact test are all-zero differences.}
\label{tab:radius}
\scriptsize
\generatedtable{generated_table_radius}
\end{table}

\begin{figure}[tbp]
\centering
\includegraphics[width=\linewidth]{fig12_endpoint_radius.pdf}
\caption{\CapRadius}
\label{fig:radius}
\end{figure}

\subsection{Complete pathway factorial}

The ablation reported above removes the target-position feedforward and observes
the gain collapse. That is consistent with the feedforward carrying the effect and
equally consistent with the two pathways working only in combination, because the
cell with assignment disabled and feedforward retained had not yet been run.
Running it completes the $2\times2$ over \NFactorialEpisodes{} episodes and
\NFactorialSeeds{} paired seeds per profile-force cell.

Fig.~\ref{fig:pathway}A and Table~\ref{tab:factorial} give all three terms. At the
transition force the target-feedforward main effect spans
$+\FactorialFeedforwardMin{}$ to $+\FactorialFeedforwardMax{}$ completion across
the three temporal profiles, with a Holm-corrected $p$ of at most
\FactorialFeedforwardHolmMax{}. The assignment main effect and the
assignment-by-feedforward interaction stay within \FactorialOtherAbsMax{}
completion of zero at every profile and force, and \NFactorialHolmSignificant{} of
the \NumFactorialCells{} tested effects clear the corrected threshold, all of them
feedforward terms. In this benchmark the directional estimate therefore acts
through the target-position feedforward, and the upwind assignment cost and its
interaction with the feedforward stay inside that near-zero band at every
profile and force. The continuous endpoint carries its own corrected family and agrees: at
the transition force the feedforward term moves the mean minimum distance to
target by \FactorialContFeedforwardMin{} to \FactorialContFeedforwardMax{}~m,
closer in every profile, while the assignment and interaction terms stay within
\FactorialContOtherAbsMax{}~m of zero and none of them clears its threshold on
that endpoint either. This is a mechanism statement scoped to this interface
and task.

\subsection{Safety and feasibility diagnostics}

The feasibility campaign completes the execution envelope with descriptive
state and clearance diagnostics. It records \NSafetyMetrics{} diagnostics over
\NSafetyEpisodes{} episodes at \SafetyForceGrid{}~N, and every episode in both
arms remains finite. Table~\ref{tab:safety} gives the full diagnostic record and
Fig.~\ref{fig:pathway}B shows the observed clearance range. The closest
airframe pair is \SafetyWorstSeparation{}~m, in the
\SafetyWorstSeparationArm{} arm at \SafetyWorstSeparationForce{}~N; the
transition-force reference values are \SafetyAgnosticSeparationKnee{}~m and
\SafetyAwareSeparationKnee{}~m, and the largest motor-limit fraction is
\SafetyMotorMax{}. \SafetyGeomHullContacts{} of the \NSafetyEpisodes{} episodes
reach the declared hull-contact distance \SafetyGeomHullContactM{}~m. These
measurements define the benchmark's feasibility boundary and identify the
additional separation constraint required by a physical implementation.

That gap is reported as an observation about these episodes rather than as a
property of the allocator, because most of the metric's variance belongs to the
benchmark instead of to either arm. The scene's minimum target separation
correlates with the observed minimum airframe separation at
\SafetyGeomTargetCorr{}, and \SafetyGeomCloseTargetContacts{} of the
\SafetyGeomHullContacts{} sub-contact episodes fall in the half of scenes whose
targets sit closer than twice the capture radius, against
\SafetyGeomFarTargetContacts{} in the other half. Within a scene, where both arms
fly identical targets, the arm that finished more airframes flew closer in
\SafetyGeomPairedCloser{} pairs and farther in \SafetyGeomPairedFarther{}, with
\SafetyGeomPairedTied{} ties: an airframe that reaches its target sits at target
separation, while one carried off by the wind ends far from everything. The metric
tracks arrival, so on this measure the arm that completes more of the task scores
worse. The scale of the numbers is structural for the same reason. Targets are
drawn with no minimum-separation constraint while completion is scored inside the
capture radius, so two airframes can both be scored finished with their hulls
overlapping, in \SafetyGeomOverlapPrimary{}\% of scenes at the primary radius and
between \SafetyGeomOverlapMin{}\% and \SafetyGeomOverlapMax{}\% across the radii
of the endpoint sweep; that sweep therefore varies how physically admissible
success is as well as how strict it is.

The transition result does not rest on these episodes: completion and observed
separation correlate at \SafetyGeomCompletionCorr{} across the campaign. What the
clearance row supports is narrower than a comparison between allocators, and it is
a caution about the task rather than about the feedforward. Contact-scale
separations (reaching an observed minimum of \SafetyWorstSeparation{}~m) were
reached on a four-airframe simulated task whose objective carries no collision
constraint and whose physics branch models no inter-airframe downwash, at mean
separations of roughly two hull diameters, which is inside the regime where rotor
wake would dominate a physical flight. The observed minimum separation is
therefore retained as a baseline failure-mode diagnostic of the unconstrained
allocation--control interface: directional wind awareness alone cannot ensure
safe separation in crossing geometries. Only the scalar three-dimensional
separation is logged, so a pair separated mostly in altitude could clear the hull
at these distances and the counts above are an upper bound on contacts. A
physical multi-vehicle deployment strictly requires composing this feedforward
layer with an explicit outer-loop safety barrier—such as Control Barrier
Functions (CBFs), velocity obstacles, or safety ellipsoids—to guarantee clearance
and prevent downwash entanglement. What was checked is the list in
Table~\ref{tab:safety}; what was not checked includes collision avoidance under an
active safety filter, degraded or failed airframes, communication loss, energy
budgets and physical multi-rotor flight.

\begin{table}[tbp]
\centering
\caption{Complete assignment-by-feedforward factorial on paired completion, over
\NFactorialSeeds{} paired seeds per profile-force cell. Each term is a seed-level
paired contrast within the cell and ``Holm $p$'' is corrected across the
\NumFactorialCells{} tested effects. Cells with an undefined exact test are
all-zero differences.}
\label{tab:factorial}
\scriptsize
\generatedtable{generated_table_factorial}
\end{table}

\begin{table}[tbp]
\centering
\caption{Descriptive safety and feasibility diagnostics over \NSafetyEpisodes{}
episodes at a constant profile, \NSafetySeeds{} seeds per allocator and force.
Each cell gives the mean over seeds with the worst observed episode in
parentheses; ``worst'' is the minimum for separation and the maximum for the
remaining rows. No test is reported and none is implied.}
\label{tab:safety}
\small
\generatedtable{generated_table_safety}
\end{table}

\begin{figure}[tbp]
\centering
\includegraphics[width=\linewidth]{fig13_pathway_safety.pdf}
\caption{\CapPathwaySafety}
\label{fig:pathway}
\end{figure}

\subsection{Evidence}

The campaign results are read on their declared inferential units and presented as
a layered synthesis that keeps each view's unit visible. Table~\ref{tab:evidence} makes
the hierarchy explicit: \NEvidenceBlocks{} synthesis blocks span
\NEvidenceCampaigns{} campaign views and \NEvidenceRows{} episode rows, while
their units remain the ten primary seeds, twelve crossed layout clusters (each
after heading aggregation) or thirty paired seeds specified by each design. This
descriptive architecture keeps every estimate tied to the design that produced it.

Fig.~\ref{fig:evidence}A places the transition estimates on one rate scale while
retaining each campaign's interval. The primary oracle estimate is
$+\KneeDelta{}$; the crossed layout estimate is $+\DesignedKneeDelta{}$;
the causal temporal profiles span $+\TemporalSlowDelta{}$ to
$+\TemporalConstantDelta{}$; and the homogeneous size axis spans
$+\ScaleKneeMin{}$ to $+\ScaleKneeMax{}$. The estimate-noise rows show the
same endpoint from the exact vector through the tested high-noise edge. The
agreement is therefore visible as a set of condition-specific estimates, with
design differences retained rather than compressed into a pooled number. Fig.~\ref{fig:evidence}B reports episode volume on a logarithmic axis, with each
point labelled by the seed or layout-cluster unit supporting its interval. Together the
panels show the breadth of the campaign while keeping episode volume separate from
inferential replication. The three reviewer-validation blocks appear in the volume
panel and in Table~\ref{tab:evidence} but not in the estimate panel, because none
of them estimates the same quantity: the radius sweep varies the endpoint that
defines the estimate, the factorial contrasts pathways rather than allocators, and
the safety block reports no estimate at all.

\begin{table}[tbp]
\centering
\caption{Cross-campaign evidence architecture. ``Rows'' counts episode records and
``Inferential unit'' gives the seed or layout-cluster unit used for intervals. The
pathway-ablation and gain-sensitivity records form one block; transition values
are descriptive estimates that retain each view's inferential unit.}
\label{tab:evidence}
\small
\generatedtable{generated_table_evidence}
\end{table}

\begin{figure}[tbp]
\centering
\includegraphics[width=\linewidth]{fig14_evidence_summary.pdf}
\caption{\CapEvidence}
\label{fig:evidence}
\end{figure}

\subsection{The transition belongs to the baseline's integral action}

The primary ladder is a measurement of the stock controller's integral
action. The stock position layer uses integral gain \StockKi{}~N/m/s,
wind-up time \IntegralTauS{}~s, clamp \IntegralClampXy{}~m/s and a
\EpisodeHorizonS{}~s episode, so the largest steady integral force is
\MaxSteadyIntegralForceN{}~N. Under those settings the bound primary sweep
reproduces \AuditStockMatch{} of \AuditStockMatch{} checked episode fields;
that check confirms the audit flew the same stock controller, not that the
transition is invariant to Ki. Fig.~\ref{fig:integral-gain} shows the
0.15~N paired gain against Ki. Doubling Ki to 0.10~N/m/s collapses the
advantage from \AuditStockGain{} to \AuditKiTenGain{} (Wilcoxon
\AuditKiTenWilcoxonP{}, sign test \AuditKiTenSignP{}, $t$ interval
[\AuditKiTenCiLo{}, \AuditKiTenCiHi{}]), with baseline and aware completions
\AuditKiTenBase{} and \AuditKiTenAware{}. Both arms then first fall below
half-completion at \AuditKiTenKnee{}~N. Clamp 4 and clamp 8 change
\AuditClampChanged{} of the 0.15~N completions at the stock horizon, so the
pre-declared P1 verdict is \AuditPOneOverall{}: the Ki half is confirmed and
the clamp half is not, at \EpisodeHorizonS{}~s.

Lengthening the episode changes that clamp reading. At 12~s the stock
baseline completion at 0.15~N is \AuditAOneBTwelveBase{}, and raising the
clamp to 8 \AuditClampMattersLongHorizon{} change the 0.20~N completion. The
horizon sweep in Fig.~\ref{fig:horizon} makes the same point along time:
paired gain is \AuditGainTFour{} at 4~s, \AuditGainTSix{} at the stock
horizon, \AuditGainTEight{} at 8~s, then \AuditGainTTen{} and
\AuditGainTTwelve{} once $T\ge 10$~s, where the 0.15~N advantage vanishes.
The fine force ladder gives the transition a width rather than a single
rung: the baseline logistic 10--90 width is \AuditWidthTenNinety{}~N
(\AuditWidthCiLo{} to \AuditWidthCiHi{}), narrower than the 0.05~N primary
grid, so ``one of seven levels'' is a sampling artefact of that grid.

The assignment audit locates the gain outside the assignment path. On the
aware arm, the initial assignment differs from Euclidean in \AuditInitialDivergence{} of
\AuditAwareEpisodes{} episodes. Used (replan) assignments diverge in
\AuditUsedDivergence{} episodes, of which \AuditUsedDivergencePrimary{}
fall on the primary-sweep seeds and \AuditUsedDivergenceKnee{} sit at
0.15~N. A pre-declared sweep of the penalty weight over
$\lambda\in\{\LambdaGrid{}\}$~m/N (Fig.~\ref{fig:lambda-sweep};
\LambdaAwareEpisodes{} aware episodes per weight on the same seeds, forces
and frozen sources) bounds that inertness. The $\lambda=0.25$ row reproduces
the audit above exactly (\AuditInitialDivergence{} initial,
\LambdaUsedDivAtQuarter{} used-divergence episodes). Initial divergence is
zero at every tested weight through $\lambda=\LambdaInertMax{}$~m/N and first
appears at $\lambda=\LambdaFirstDivergence{}$~m/N (\LambdaInitialDivAtOne{}
of \LambdaAwareEpisodes{}; \LambdaInitialDivAtTwo{} at
$\lambda=\LambdaGridMax{}$~m/N), while used-assignment divergence rises
monotonically from \LambdaUsedDivAtZero{} at $\lambda=0$ to
\LambdaUsedDivAtTwo{} at $\lambda=\LambdaGridMax{}$~m/N. The aware
completion at 0.15~N is \LambdaCompletionFlat{} at every weight through
$\lambda=\LambdaInertMax{}$~m/N. Under the reading rule stated before the
sweep ran, this outcome keeps the inertness claim scoped to the inherited
weight and reports the tested range over which it holds; the sweep is
descriptive by declaration and carries no confirmatory test.

The high-force edge is a crash regime, not a clean saturation. Pooled at and
above 0.20~N, a fraction \AuditNeverCloseBase{} (agnostic) and
\AuditNeverCloseAware{} (aware) of airframes never come within 0.5~m of
their target. After a completion is recorded, a fraction
\AuditGroundContactAfter{} later contact the ground and
\AuditAltLossAfter{} lose more than 0.10~m of altitude, consistent with a
frozen rotor command rather than a held hover.

The capture-radius fingerprint likewise departs from a simple saturation
curve. Baseline $F_{50}$ against capture radius has slope
\AuditRadiusSlope{}~N/m (interval [\AuditRadiusSlopeLo{},
\AuditRadiusSlopeHi{}]), which excludes 1. A finer force grid over the same
seven radii (\AFiveBForceStepN{}~N steps, \AFiveBEpisodes{} episodes), added
after the pre-declaration, gives slope \AFiveBSlope{}~N/m (interval
[\AFiveBSlopeLo{}, \AFiveBSlopeHi{}]); it lies inside the pre-declared
interval and is reported as a resolution sensitivity beside the pre-declared
estimate, which remains the estimate of record.

\begin{figure}[tbp]
\centering
\includegraphics[width=\linewidth]{fig14_integral_gain.pdf}
\caption{\CapIntegralGain}
\label{fig:integral-gain}
\end{figure}

\begin{figure}[tbp]
\centering
\includegraphics[width=\linewidth]{fig15_horizon.pdf}
\caption{\CapHorizon}
\label{fig:horizon}
\end{figure}

\begin{figure}[tbp]
\centering
\includegraphics[width=\linewidth]{fig16_lambda_sweep.pdf}
\caption{\CapLambdaSweep}
\label{fig:lambda-sweep}
\end{figure}

\section{Discussion}

The identified mechanism is target-position feedforward compensating a slow
integral inside a \EpisodeHorizonS{}~s episode, and the controller audit
fixes the settings on which that identification rests. Prediction P1 is
\AuditPOneOverall{}: doubling Ki collapses the 0.15~N paired advantage to
\AuditKiTenGain{}, confirming the integral-gain half of the
pre-declaration, while clamp 4 and clamp 8 change \AuditClampChanged{} of the
0.15~N completions at the stock horizon, refuting the clamp half at
$T=\EpisodeHorizonS{}$~s. The same clamp
\AuditClampMattersLongHorizon{} matter at 12~s. Initial-assignment inertness
holds at every tested weight through $\lambda=\LambdaInertMax{}$~m/N and
first breaks at $\lambda=\LambdaFirstDivergence{}$~m/N; the bound claim stays
scoped to the inherited weight, and the sweep reports the tested range
descriptively. These bounds make the operating envelope a property of the
stock integral and the short horizon.

The headline gain travels through the feedforward path. Initial assignment
differs from Euclidean in \AuditInitialDivergence{} of \AuditAwareEpisodes{}
episodes, used (replan) assignments diverge in \AuditUsedDivergence{}
episodes, and the complete factorial localizes the completion change to
target-position feedforward, with the assignment main effect and the
assignment-by-feedforward interaction within \FactorialOtherAbsMax{}
completion of zero. The pre-declared $\lambda$ sweep is descriptive and stops
at $\lambda=\LambdaGridMax{}$~m/N; the allocation reading is bounded by that
grid. Support campaigns keep the same paired estimand: relative noise through
\BestSigma{} preserves the transition, and paired gains stay positive from
\ScaleMinSize{} through \ScaleMaxSize{} airframes. Rerunning the ladder
across capture radii keeps a transition of comparable size and moves it from
\RadiusTightPeakForce{} to \RadiusWidePeakForce{}~N, so \KneeLevel{}~N is
the transition of this operating point.

The geometry contrasts and the descriptive safety record bound the envelope
from two further directions. The contrasts map named scene axes within the
same paired design, and the clearance diagnostic marks the feasibility
boundary of the benchmark as a property of the task: the closest recorded
airframe pair is \SafetyWorstSeparation{}~m on targets drawn without a
separation constraint, and a physical deployment addresses that boundary by
composing the feedforward layer with an explicit safety barrier. Every
measurement here is a simulation. The present envelope is a computational
calibration reference against the stock integral, the \EpisodeHorizonS{}~s
horizon and the declared capture radius, and the transfer protocol names the
quantities a hardware or hardware-in-the-loop extension revalidates.
Heterogeneous vehicles, obstacle topology, penalty weights beyond
\LambdaGridMax{}~m/N and field-sensor traces are the named transfer axes for
that next stage; general allocator properties, hardware principles and field
results lie outside the measured scope.

\section*{Declarations}

%% The Declarations heading above must stay the first thing after the
%% Discussion prose (no comment or \ifdefined line before it): the copy-edit
%% checks slice the Discussion up to the next \section command and count
%% blank-line-separated chunks as paragraphs.
%% Funding and competing-interest wording adopted from the owner's packaging
%% script (scripts/build_springer_submission_package.py), shipped in every
%% packet since 2026-08-28 (ruling R16, 2026-09-09). FUNDING_AND_CONFLICTS
%% stays OPEN for the owner's final read.
\subsection*{Funding}

No external funding was received for this study.

\subsection*{Competing interests}

The author declares that they have no competing financial or non-financial
interests that are directly or indirectly related to the work submitted for
publication.

\subsection*{Ethics approval and consent to participate}

Not applicable. The study is computational and involves no human
participants, animal subjects, personally identifying data or field flights.

\subsection*{Consent for publication}

Not applicable; no identifiable participant data are included.

\subsection*{Data and code availability}

All simulation code, trajectory records and analysis scripts supporting the
findings of this study, which are the data and code behind every figure and
table, are available in the public repository and archive; the simulation
substrate is an open-source physics engine with the disturbance injection
described in the Methods.
\input{generated_availability}

\subsection*{Materials availability}

Not applicable; the study generated no physical materials.

\subsection*{Author contributions}

\ifdefined\ANONYMOUS
The sole author performed conceptualization, methodology, software,
validation, formal analysis, investigation, visualization, writing of the
original draft, and review and editing.
\else
Zeyu Fu performed conceptualization, methodology, software, validation,
formal analysis, investigation, visualization, writing of the original draft,
and review and editing.
\fi

\subsection*{Use of generative AI tools}

Generative AI tools were used solely for code engineering, refactoring and
language polishing under the author's direction. All computational
experiments, analyses and manuscript text were reviewed, verified and edited
by the author, who retains full responsibility for the integrity and final
wording of the work.

%% Identified builds only; draft-only until the owner records the sentence.
\ifdefined\ANONYMOUS\else
\ifdefined\DRAFTTODO
\subsection*{Acknowledgements}

\ownertodo{add acknowledgements or delete this heading before submission;
none are recorded.}
\fi
\fi


{\footnotesize
\bibliographystyle{unsrt}
\bibliography{references}
}

\end{document}
"
% Identified output is the default.  hyperref writes the name into the PDF
% /Author key as well as the title page, so both go through this one macro.
\ifdefined\ANONYMOUS\newcommand{\PDFAUTHOR}{}\else\newcommand{\PDFAUTHOR}{Zeyu Fu}\fi
\hypersetup{
  pdftitle={From directional sensing to task completion: the operating envelope of target-position feedforward against a slow integral in wind-disturbed quadrotor waypoint capture},
  pdfauthor={\PDFAUTHOR},
  pdfsubject={Operating envelope of target-position feedforward against a slow integral},
  pdfkeywords={disturbance feedforward, integral control, wind estimation, thresholded endpoint, task allocation, quadrotor swarm, reproducibility}
}
\numberwithin{equation}{section}

\graphicspath{{../figs/out/}}
\setlength{\emergencystretch}{2em}

%% Every float here is small enough to share a page with text; the defaults are
%% conservative enough to push each one onto a page of its own.
\renewcommand{\topfraction}{0.9}
\renewcommand{\bottomfraction}{0.7}
\renewcommand{\textfraction}{0.08}
\renewcommand{\floatpagefraction}{0.8}
\setcounter{topnumber}{2}
\setcounter{totalnumber}{3}
\captionsetup[table]{skip=5pt}
\captionsetup[figure]{skip=4pt}

%% Numbers and tables are emitted by the figure code from hash-bound evidence.
% Generated by the figure script from hash-bound evidence. Do not hand-edit.
\newcommand{\NSeeds}{10}
\newcommand{\NWindLevels}{7}
\newcommand{\NEpisodesOracle}{140}
\newcommand{\NEpisodesNoise}{150}
\newcommand{\NDrones}{4}
\newcommand{\SigmaGrid}{0, 0.1, 0.3, 0.5, 1}
\newcommand{\CaptureRadius}{0.15}
\newcommand{\HoverThrust}{0.265}
\newcommand{\DroneMass}{27}
\newcommand{\KneeLevel}{0.15}
\newcommand{\KneePctHover}{57}
\newcommand{\SaturationLevel}{0.10}
\newcommand{\CollapseLevel}{0.20}
\newcommand{\KneeAgnostic}{0.05}
\newcommand{\KneeAgnosticSD}{0.10}
\newcommand{\KneeAware}{0.60}
\newcommand{\KneeAwareSD}{0.25}
\newcommand{\KneeDelta}{0.55}
\newcommand{\KneeP}{0.00390625}
\newcommand{\BonferroniAlpha}{0.0071}
\newcommand{\PreregSigma}{0.25}
\newcommand{\NearestSigma}{0.30}
\newcommand{\NearestSigmaDelta}{0.55}
\newcommand{\NearestSigmaLower}{0.44}
\newcommand{\NearestSigmaP}{0.00195}
\newcommand{\SigmaZeroDelta}{0.55}
\newcommand{\SigmaZeroAware}{0.60}
\newcommand{\BestSigma}{0.50}
\newcommand{\BestSigmaDelta}{0.65}
\newcommand{\MaxSigma}{1.00}
\newcommand{\MaxSigmaDelta}{0.30}
\newcommand{\MaxSigmaLower}{0.09}
\newcommand{\MaxSigmaP}{0.0312}
\newcommand{\WeakestLower}{0.09}
\newcommand{\NoiseComparisons}{15}
\newcommand{\NoiseBonferroni}{0.0033}
\newcommand{\BridgeLowerPairs}{8}
\newcommand{\BridgeDistanceShift}{-0.018}
\newcommand{\BridgeCompletionGain}{0.55}
\newcommand{\AwareTopSeeds}{7}
\newcommand{\AwareZeroSeeds}{1}
\newcommand{\AgnosticZeroSeeds}{8}
\newcommand{\MissAgnosticKnee}{0.244}
\newcommand{\MissAwareKnee}{0.226}
\newcommand{\MissAgnosticCollapse}{0.597}
\newcommand{\MissAwareCollapse}{0.595}
\newcommand{\NHeadings}{6}
\newcommand{\LargestHeadingGroup}{4}
\newcommand{\HeadingGroupSpan}{1.60}
\newcommand{\UnsampledArc}{152}
\newcommand{\NRepeats}{7}
\newcommand{\RepeatTol}{2}
\newcommand{\ClosestRepeat}{0.05}
\newcommand{\WidestRepeat}{1.60}
\newcommand{\NearestDistinct}{25.2}
\newcommand{\NInformativeLevels}{3}
\newcommand{\NDegenerateLevels}{4}
\newcommand{\RepeatGapKnee}{0.75}
\newcommand{\RepeatGapKneeDrones}{3}
\newcommand{\RepeatGapShoulder}{0.25}
\newcommand{\RepeatGapShoulderDrones}{1}
\newcommand{\NDisagreeingLevels}{2}
\newcommand{\RepeatDisagreeKnee}{6}
\newcommand{\RepeatDisagreeShoulder}{4}
\newcommand{\RepeatDisagreeCollapse}{0}
\newcommand{\CollapseSeedSpread}{0.25}
\newcommand{\NEpisodesDesigned}{4032}
\newcommand{\NAssignedHeadings}{24}
\newcommand{\HeadingStep}{15}
\newcommand{\NScenes}{12}
\newcommand{\NPairsPerForce}{288}
\newcommand{\DesignedArc}{15}
\newcommand{\DesignedKneeDelta}{0.62}
\newcommand{\DesignedKneeLower}{0.57}
\newcommand{\DesignedKneeUpper}{0.69}
\newcommand{\DesignedKneeAgnostic}{0.08}
\newcommand{\DesignedKneeAware}{0.70}
\newcommand{\DesignedKneeAdverse}{1}
\newcommand{\DesignedShareHeading}{20}
\newcommand{\DesignedShareScene}{15}
\newcommand{\DesignedShareResidual}{65}
\newcommand{\ZeroSeparationGap}{1.25}
\newcommand{\ZeroSeparationLow}{-0.25}
\newcommand{\ZeroSeparationHigh}{+1.00}
\newcommand{\HeadingOnlyGap}{1.25}
\newcommand{\DesignedShoulderDelta}{0.030}
\newcommand{\DesignedCollapseDelta}{0.044}
\newcommand{\NReproEpisodes}{140}
\newcommand{\NReproFields}{6}
\newcommand{\NReplanQuiet}{6}
\newcommand{\NReplanLoud}{4}
\newcommand{\ReplanQuietDelta}{0.67}
\newcommand{\ReplanLoudDelta}{0.38}
\newcommand{\ReplanKneeAgnostic}{3}
\newcommand{\ReplanKneeAware}{6}
\newcommand{\NTemporalEpisodes}{1080}
\newcommand{\NTemporalProfiles}{3}
\newcommand{\NTemporalEstimators}{3}
\newcommand{\NTemporalSeeds}{30}
\newcommand{\NTemporalLegacyRows}{60}
\newcommand{\TemporalConstantDelta}{0.63}
\newcommand{\TemporalConstantLower}{0.52}
\newcommand{\TemporalConstantUpper}{0.75}
\newcommand{\TemporalConstantP}{7.45e-9}
\newcommand{\TemporalSlowDelta}{0.42}
\newcommand{\TemporalSlowLower}{0.29}
\newcommand{\TemporalSlowUpper}{0.56}
\newcommand{\TemporalSlowP}{5.13e-6}
\newcommand{\TemporalFastDelta}{0.60}
\newcommand{\TemporalFastLower}{0.48}
\newcommand{\TemporalFastUpper}{0.72}
\newcommand{\TemporalFastP}{7.45e-9}
\newcommand{\TemporalTrackingConstant}{3.9}
\newcommand{\TemporalTrackingSlow}{9.6}
\newcommand{\TemporalTrackingFast}{11.9}
\newcommand{\TemporalShoulderMax}{0.11}
\newcommand{\NAblationEpisodes}{810}
\newcommand{\AllocOnlyConstant}{0.01}
\newcommand{\AllocOnlySlow}{0.00}
\newcommand{\AllocOnlyFast}{0.01}
\newcommand{\CompleteConstant}{0.63}
\newcommand{\CompleteSlow}{0.42}
\newcommand{\CompleteFast}{0.60}
\newcommand{\NGainEpisodes}{1890}
\newcommand{\FeedforwardGainDefault}{0.35}
\newcommand{\GainPlateauMin}{0.50}
\newcommand{\GainPlateauMax}{0.71}
\newcommand{\GainDefaultConstant}{0.63}
\newcommand{\GainDefaultSlow}{0.42}
\newcommand{\GainDefaultFast}{0.60}
\newcommand{\GainEdgeConstant}{0.67}
\newcommand{\GainEdgeSlow}{0.52}
\newcommand{\GainEdgeFast}{0.58}
\newcommand{\NGeometryEpisodes}{1440}
\newcommand{\NGeometryConditions}{8}
\newcommand{\NGeometrySeeds}{30}
\newcommand{\NGeometryClusters}{30}
\newcommand{\NGeometryTermSensitivity}{120}
\newcommand{\NCrossedClusters}{12}
\newcommand{\NCrossedCells}{288}
\newcommand{\NCrossedBootstrap}{10000}
\newcommand{\GeometryKneeCrossingP}{0.00423}
\newcommand{\GeometryKneeCrossingHolm}{0.0381}
\newcommand{\GeometryKneeSpread}{-0.06}
\newcommand{\GeometryKneeSpreadLower}{-0.14}
\newcommand{\GeometryKneeSpreadUpper}{0.01}
\newcommand{\GeometryKneeRange}{-0.05}
\newcommand{\GeometryKneeRangeLower}{-0.10}
\newcommand{\GeometryKneeRangeUpper}{0.00}
\newcommand{\GeometryKneeCrossing}{0.12}
\newcommand{\GeometryKneeCrossingLower}{0.04}
\newcommand{\GeometryKneeCrossingUpper}{0.20}
\newcommand{\GeometryMDEFamily}{9}
\newcommand{\GeometryMDEAlpha}{0.0056}
\newcommand{\GeometryMDEPower}{0.80}
\newcommand{\GeometryRangeClusterN}{30}
\newcommand{\GeometryRangeClusterSD}{0.139}
\newcommand{\GeometryRangeSensitivityLower}{-0.11}
\newcommand{\GeometryRangeSensitivityUpper}{0.01}
\newcommand{\GeometryRangeMDE}{0.098}
\newcommand{\GeometryKneeNominal}{0.54}
\newcommand{\GeometryKneeBest}{0.66}
\newcommand{\NScaleEpisodes}{720}
\newcommand{\NScaleSizes}{4}
\newcommand{\NScaleForces}{3}
\newcommand{\NScaleSeeds}{30}
\newcommand{\ScaleSizeGrid}{2, 4, 6, 8}
\newcommand{\ScaleForceGrid}{0.10, 0.15, 0.20}
\newcommand{\ScaleMinSize}{2}
\newcommand{\ScaleMaxSize}{8}
\newcommand{\ScaleKneeMin}{0.61}
\newcommand{\ScaleKneeMax}{0.67}
\newcommand{\ScaleKneeTwo}{0.67}
\newcommand{\ScaleKneeFour}{0.63}
\newcommand{\ScaleKneeSix}{0.61}
\newcommand{\ScaleKneeEight}{0.62}
\newcommand{\ScaleKneeTwoLower}{0.55}
\newcommand{\ScaleKneeTwoUpper}{0.78}
\newcommand{\ScaleKneeFourLower}{0.52}
\newcommand{\ScaleKneeFourUpper}{0.75}
\newcommand{\ScaleKneeSixLower}{0.51}
\newcommand{\ScaleKneeSixUpper}{0.71}
\newcommand{\ScaleKneeEightLower}{0.53}
\newcommand{\ScaleKneeEightUpper}{0.72}
\newcommand{\ScaleKneePMax}{2.61e-8}
\newcommand{\ScaleKneePMin}{3.73e-9}
\newcommand{\ScaleKneeHolmMax}{2.35e-7}
\newcommand{\NRadiusEpisodes}{420}
\newcommand{\NRadiusSeeds}{10}
\newcommand{\NRadiusLevels}{3}
\newcommand{\RadiusGrid}{0.10, 0.15, 0.20}
\newcommand{\RadiusTightLevel}{0.10}
\newcommand{\RadiusWideLevel}{0.20}
\newcommand{\RadiusSpanM}{0.10}
\newcommand{\RadiusHolmFamily}{21}
\newcommand{\RadiusTightPeakForce}{0.10}
\newcommand{\RadiusTightPeakDelta}{0.55}
\newcommand{\RadiusAnchorPeakForce}{0.15}
\newcommand{\RadiusAnchorPeakDelta}{0.55}
\newcommand{\RadiusWidePeakForce}{0.20}
\newcommand{\RadiusWidePeakDelta}{0.42}
\newcommand{\RadiusPeakDeltaMin}{0.42}
\newcommand{\RadiusPeakDeltaMax}{0.55}
\newcommand{\RadiusPeakForceSpan}{0.10}
\newcommand{\RadiusHolmSurvivors}{1}
\newcommand{\RadiusAnchorHolmP}{0.0781}
\newcommand{\RadiusContinuousCloser}{16}
\newcommand{\RadiusContinuousCells}{21}
\newcommand{\RadiusContinuousHolmSurvivors}{2}
\newcommand{\NFactorialEpisodes}{270}
\newcommand{\NFactorialSeeds}{30}
\newcommand{\FactorialFeedforwardMin}{0.42}
\newcommand{\FactorialFeedforwardMax}{0.63}
\newcommand{\FactorialFeedforwardHolmMax}{1.28e-4}
\newcommand{\FactorialAssignmentAbsMax}{0.004}
\newcommand{\FactorialInteractionAbsMax}{0.008}
\newcommand{\FactorialOtherAbsMax}{0.008}
\newcommand{\NumFactorialCells}{27}
\newcommand{\NFactorialHolmSignificant}{3}
\newcommand{\FactorialContFeedforwardMin}{-0.034}
\newcommand{\FactorialContFeedforwardMax}{-0.010}
\newcommand{\FactorialContOtherAbsMax}{0.004}
\newcommand{\NSafetyEpisodes}{60}
\newcommand{\NSafetySeeds}{10}
\newcommand{\NSafetyMetrics}{6}
\newcommand{\SafetyForceGrid}{0.10, 0.15, 0.20}
\newcommand{\SafetyWorstSeparation}{0.049}
\newcommand{\SafetyWorstSeparationArm}{wind-aware}
\newcommand{\SafetyWorstSeparationForce}{0.20}
\newcommand{\SafetyAgnosticSeparationKnee}{0.100}
\newcommand{\SafetyAwareSeparationKnee}{0.085}
\newcommand{\SafetyMotorMax}{0.85}
\newcommand{\SafetyNonfiniteTotal}{0}
\newcommand{\SafetyStepMaxMs}{13.9}
\newcommand{\SafetyGeomTargetCorr}{0.76}
\newcommand{\SafetyGeomCompletionCorr}{-0.09}
\newcommand{\SafetyGeomHullContacts}{10}
\newcommand{\SafetyGeomHullContactM}{0.12}
\newcommand{\SafetyGeomCloseTargetContacts}{9}
\newcommand{\SafetyGeomFarTargetContacts}{1}
\newcommand{\SafetyGeomPairedCloser}{9}
\newcommand{\SafetyGeomPairedFarther}{2}
\newcommand{\SafetyGeomPairedTied}{19}
\newcommand{\SafetyGeomOverlapPrimary}{30.9}
\newcommand{\SafetyGeomOverlapMin}{11.5}
\newcommand{\SafetyGeomOverlapMax}{54.5}
\newcommand{\NEvidenceBlocks}{10}
\newcommand{\NEvidenceCampaigns}{11}
\newcommand{\NEvidenceRows}{11,012}
\newcommand{\StockKi}{0.05}
\newcommand{\IntegralTauS}{8}
\newcommand{\EpisodeHorizonS}{6}
\newcommand{\IntegralClampXy}{2}
\newcommand{\MaxSteadyIntegralForceN}{0.1}
\newcommand{\AuditStockMatch}{140}
\newcommand{\AuditKiTenGain}{+0.008}
\newcommand{\AuditKiTenWilcoxonP}{0.84}
\newcommand{\AuditKiTenSignP}{0.73}
\newcommand{\AuditKiTenCiLo}{-0.049}
\newcommand{\AuditKiTenCiHi}{+0.066}
\newcommand{\AuditKiTenBase}{0.79}
\newcommand{\AuditKiTenAware}{0.80}
\newcommand{\AuditKiTenKnee}{0.20}
\newcommand{\AuditPOneOverall}{mixed}
\newcommand{\AuditClampChanged}{0}
\newcommand{\AuditStockGain}{+0.54}
\newcommand{\AuditAOneBTwelveBase}{0.82}
\newcommand{\AuditClampMattersLongHorizon}{does}
\newcommand{\AuditGainTFour}{+0.04}
\newcommand{\AuditGainTSix}{+0.54}
\newcommand{\AuditGainTEight}{+0.14}
\newcommand{\AuditGainTTen}{-0.07}
\newcommand{\AuditGainTTwelve}{-0.07}
\newcommand{\AuditWidthTenNinety}{0.057}
\newcommand{\AuditWidthCiLo}{0.045}
\newcommand{\AuditWidthCiHi}{0.068}
\newcommand{\AuditNeverCloseBase}{0.53}
\newcommand{\AuditNeverCloseAware}{0.53}
\newcommand{\AuditGroundContactAfter}{0.87}
\newcommand{\AuditAltLossAfter}{0.94}
\newcommand{\AuditRadiusSlope}{0.68}
\newcommand{\AuditRadiusSlopeLo}{0.54}
\newcommand{\AuditRadiusSlopeHi}{0.82}
\newcommand{\AuditInitialDivergence}{0}
\newcommand{\AuditAwareEpisodes}{210}
\newcommand{\AuditUsedDivergence}{7}
\newcommand{\AuditUsedDivergencePrimary}{2}
\newcommand{\AuditUsedDivergenceKnee}{1}
\newcommand{\LambdaGrid}{0, 0.05, 0.1, 0.25, 0.5, 1, 2}
\newcommand{\LambdaGridSize}{7}
\newcommand{\LambdaGridMax}{2}
\newcommand{\LambdaReference}{0.25}
\newcommand{\LambdaInertMax}{0.5}
\newcommand{\LambdaFirstDivergence}{1.0}
\newcommand{\LambdaInitialDivAtOne}{3}
\newcommand{\LambdaInitialDivAtTwo}{19}
\newcommand{\LambdaUsedDivAtZero}{0}
\newcommand{\LambdaUsedDivAtQuarter}{7}
\newcommand{\LambdaUsedDivAtTwo}{80}
\newcommand{\LambdaAwareEpisodes}{210}
\newcommand{\LambdaEpisodes}{1,680}
\newcommand{\LambdaCompletionFlat}{0.63}
\newcommand{\AFiveBSlope}{0.705}
\newcommand{\AFiveBSlopeLo}{0.63}
\newcommand{\AFiveBSlopeHi}{0.78}
\newcommand{\AFiveBForceStepN}{0.01}
\newcommand{\AFiveBEpisodes}{8,820}

%% Draft-only owner placeholder (ruling R15, 2026-09-09). Prints ONLY when the
%% build defines \DRAFTTODO (e.g. pdflatex "\def\DRAFTTODO{1}%% Working draft. Structure is generated; prose is authored here.
%% Build from the repository root with: bash scripts/build_paper.sh <this paper dir>
\documentclass[11pt]{article}

\usepackage[utf8]{inputenc}
\usepackage[T1]{fontenc}
\usepackage{amsmath,amssymb,amsthm}
\usepackage{array}
\usepackage{booktabs}
\usepackage{graphicx}
\usepackage{tikz}
\usepackage{placeins}
\usepackage{caption}
\usepackage[margin=1in]{geometry}
% For every included figure, pdfTeX records the absolute path it was read from
% in a /PTEX.FileName object.  Those objects are not rendered, so a text dump of
% the PDF does not show them, but they travel with the file and spell out the
% build machine's home directory, user name and worktree layout.  Suppressing
% them costs nothing: they exist for debugging, not for the reader.
\pdfsuppressptexinfo=-1

\usepackage[hidelinks]{hyperref}
% Blind-review copies are compiled as
%   pdflatex "\def\ANONYMOUS{1}\input{main}"
% Identified output is the default.  hyperref writes the name into the PDF
% /Author key as well as the title page, so both go through this one macro.
\ifdefined\ANONYMOUS\newcommand{\PDFAUTHOR}{}\else\newcommand{\PDFAUTHOR}{Zeyu Fu}\fi
\hypersetup{
  pdftitle={From directional sensing to task completion: the operating envelope of target-position feedforward against a slow integral in wind-disturbed quadrotor waypoint capture},
  pdfauthor={\PDFAUTHOR},
  pdfsubject={Operating envelope of target-position feedforward against a slow integral},
  pdfkeywords={disturbance feedforward, integral control, wind estimation, thresholded endpoint, task allocation, quadrotor swarm, reproducibility}
}
\numberwithin{equation}{section}

\graphicspath{{../figs/out/}}
\setlength{\emergencystretch}{2em}

%% Every float here is small enough to share a page with text; the defaults are
%% conservative enough to push each one onto a page of its own.
\renewcommand{\topfraction}{0.9}
\renewcommand{\bottomfraction}{0.7}
\renewcommand{\textfraction}{0.08}
\renewcommand{\floatpagefraction}{0.8}
\setcounter{topnumber}{2}
\setcounter{totalnumber}{3}
\captionsetup[table]{skip=5pt}
\captionsetup[figure]{skip=4pt}

%% Numbers and tables are emitted by the figure code from hash-bound evidence.
\input{generated_numbers}
\input{captions}
\input{theorems}
\newcommand{\generatedtable}[1]{\input{#1}}

\title{From directional sensing to task completion: the operating envelope of
target-position feedforward against a slow integral in wind-disturbed quadrotor
waypoint capture}
\ifdefined\ANONYMOUS
\author{Author and affiliation withheld for blind review}
\else
\author{Zeyu Fu\\
Army Medical University, Chongqing, China\\
\texttt{fuzeyu99@126.com}\\
ORCID: 0009-0001-8329-0108}
\fi
\date{}

\begin{document}
\maketitle

\begin{abstract}
\input{abstract}
\end{abstract}

\paragraph{Keywords} disturbance feedforward, integral control, wind estimation, thresholded endpoint, task allocation, quadrotor swarm, reproducibility

\input{body}

{\footnotesize
\bibliographystyle{unsrt}
\bibliography{references}
}

\end{document}
");
%% every packet build leaves it undefined, so no unconfirmed declaration text
%% reaches a PDF. The owner replaces each use with the confirmed sentence.
\providecommand{\ownertodo}[1]{\ifdefined\DRAFTTODO[TODO-OWNER: #1]\fi}

%% Independent caption file. Every caption is self-contained: readable without
%% the body text, without cross-references, and without float numbers.

\newcommand{\CapSystemOverview}{%
Research object and endpoint bridge for the wind-aware allocation--control
measurement. Panel A is a drawn schematic of the computational contract: four
matched quadrotors, four waypoint targets and a horizontal disturbance vector.
Dashed paths represent the shared
assignment-to-control task; the two labelled paths distinguish the Euclidean
wind-agnostic cost from the wind-aware upwind cost plus target-position bias.
Panel B is a descriptive seed-paired view at the transition force
\KneeLevel{}~N: each arrow connects the same seed under the two arms in mean
closest-approach versus completion-rate coordinates. The vertical reference is
the per-airframe capture radius used to define completion. Displayed arrows use
the paired seed record as the unit; the panel reports a descriptive endpoint view
with the seed-level inferential unit preserved. %
}

\newcommand{\CapBand}{%
Operating envelope of task completion against steady horizontal wind force for a
wind-agnostic and a wind-aware allocator flying the same \NDrones{} quadrotors on
the same \NSeeds{} seeds per force level, with both panels on one shared force
axis. Panel A is the completion ladder: points are means over seeds and bars are
one standard deviation, clipped at the zero and one bounds of the rate, and the
upper axis restates force as a percentage of the \HoverThrust{}~N hover thrust of
one airframe. Panel B is the seed-paired advantage of the wind-aware over the
wind-agnostic allocator, in completion-rate points, with bars giving the
95~percent $t$ interval over the \NSeeds{} paired seeds. Panel B carries the
quantity the significance test is computed on, and its interval clears zero at
\KneeLevel{}~N alone, where the paired gain is $+\KneeDelta{}$ at an exact
two-sided signed-rank $p$ of \KneeP{}. The two panels are drawn together because
the Panel A bars are dispersions of two marginal means and overlap at the
transition even though every seed moves in the same direction, so the paired
contrast is displayed rather than left to be inferred from that overlap. Shading
separates the near-complete shoulder at and below \SaturationLevel{}~N and the
high-force edge at and above \CollapseLevel{}~N from the unshaded transition,
whose \KneeLevel{}~N level the dashed line marks and where directional
information has the largest operational leverage. Panel A's red and blue legend
identifies the two allocation arms; Panel B is a single paired-difference series
and therefore has no second arm legend.%
}

\newcommand{\CapNoise}{%
Profile of the directional signal as wind-estimate quality changes. The horizontal
axis is relative magnitude noise redrawn independently at every control step. Panel
A shows the paired advantage in completion rate at the transition force
\KneeLevel{}~N, with bars giving the 95~percent $t$ interval over \NSeeds{} paired
seeds and labels giving the exact two-sided signed-rank $p$ value; the dashed line
marks the decision level fixed before execution. Panel B places the two neighbouring
force levels on the same vertical scale, showing that the operating transition stays
localized while estimate quality changes. A paired advantage of zero marks equal
completion for the two arms. In both panels, paired advantage means aware-minus-
agnostic completion rate for the same seed. %
}

\newcommand{\CapSpread}{%
Distribution behind the two completion means at the transition force
\KneeLevel{}~N. Each episode assigns \NDrones{} quadrotors to \NDrones{} waypoints,
so a single episode can score zero, one quarter, one half, three quarters or one.
The bars count how many of the \NSeeds{} seeds land on each value for the two arms;
the distribution displays the episode-level margin carried by the aggregate
operating-envelope gain. %
}

\newcommand{\CapHeadings}{%
Coverage of the force headings drawn by the \NSeeds{} seeds and the repeat groups
in the displayed record. Panel A places each seed's heading on the circle, one
spoke per seed; the heading is fixed by the seed and reused at every force level,
in both allocators and in the estimate-noise sweep. The draws resolve to
\NHeadings{} distinct headings, with \NRepeats{} near-matched pairs and a largest
unvisited arc of \UnsampledArc{} degrees. Panel B orders all
$\NSeeds{}(\NSeeds{}-1)/2$ pairwise separations on a logarithmic axis; the empty
gap between the retained group and the next separation makes the stated grouping
stable across the displayed tolerance interval. %
}

\newcommand{\CapRepeats}{%
Context spread under near-matched force headings. At each force level the bar gives
the largest completion difference between two episodes of the same allocator whose
headings differ by less than \RepeatTol{} degrees, expressed as the number of
airframes, out of \NDrones{}, by which the episodes differ. The strip below records
whether any seeds at that force level varied in completion; a zero bar at a fully
uniform level marks a saturated shoulder or floor, while the informative levels
show where waypoint context changes the outcome under nearly identical directional
input. %
}

\newcommand{\CapDesigned}{%
Heading-by-layout operating-envelope map from a crossed design. \NAssignedHeadings{}
headings \HeadingStep{} degrees apart are each flown with all \NScenes{} waypoint
layouts, at every force level and under both allocators, for \NEpisodesDesigned{}
episodes and \NPairsPerForce{} paired comparisons per level. Panel A shows paired advantage
against force, with the band giving an ordinary-bootstrap 95~percent interval over
the \NCrossedClusters{} layout-level means after averaging the fixed headings
(\NCrossedBootstrap{} deterministic resamples); crosses mark the \NSeeds{} seed
estimates from the independently drawn-heading condition.
Panel B opens the transition force into its heading-by-layout rectangle, one cell
per paired episode; rows are assigned headings and columns are scene layouts, and
colour is the paired completion advantage. The heat map displays the same cells
used to compute the heading, layout and interaction shares. %
}

\newcommand{\CapMiss}{%
Continuous tracking quantity behind the thresholded completion endpoint. Points give
the mean closest approach of each quadrotor to its assigned waypoint, averaged over
seeds, with bars of one standard deviation; the dashed line is the
\CaptureRadius{}~m capture radius that defines completion. Both allocators are
measured at the same forces, and the two series are offset horizontally by a fixed
0.0055~N purely so that their overlapping bars remain separable; the axis ticks
mark the true force levels. The two allocators occupy
nearby trajectories across the force ladder, and the measured distance shift at the
transition moves aircraft across the capture boundary to produce the task-level
operating gain. %
}

\newcommand{\CapTemporal}{%
Temporal operating-envelope extension for a causal directional estimate. Panel A
shows paired completion gain against true force for constant, slow-gust and fast-gust
profiles, with exact, causal and biased-causal estimate conditions and 95~percent
intervals over \NTemporalSeeds{} paired seeds. Panel B shows mean relative vector-tracking error
for the same profiles and estimate conditions. The transition remains concentrated
near \KneeLevel{}~N, while the tracking-error profile quantifies the cost of temporal
rate and links estimator dynamics to the measured task-level response. Colours and
symbols identify the three estimate conditions consistently across both panels. %
}

\newcommand{\CapMechanism}{%
Mechanism and gain-sensitivity extension at the transition force \KneeLevel{}~N.
Panel A gives paired completion gain across the tested target-feedforward gains for
each true-force profile; the dashed line marks the inherited gain
\FeedforwardGainDefault{}~m/N. Panel B compares the complete arm with an
allocation-cost-only ablation that retains directional assignment and selectively
removes the target-position bias. The separation identifies the pathway through
which the directional estimate is converted into task completion in this benchmark.
Intervals in both panels are paired 95~percent $t$ intervals over the declared
profile-specific seeds; the upper gain-grid point marks the measured edge of the
tested gain envelope. %
}

\newcommand{\CapGeometry}{%
Factorised scene-geometry extension of the directional operating envelope. Panel A
shows high-minus-low contrasts for target spread, target range and crossing
orientation at each tested force; intervals are 95~percent cluster-aware intervals
over 30 seed-level factor contrasts after averaging the four crossed settings
within seed; the 120-term version is sensitivity-only. Raw exact signed-rank
values and nine-contrast Holm adjustments define the support-family readout, and
the zero line marks no change in the wind-aware advantage. Panel B shows the mean
paired advantage in all eight scene conditions
across the three force levels, with each row defined by one spread, range and
crossing setting. The nominal--near--aligned cell is the legacy compatibility
condition; the remaining cells are deterministic scene transformations. Together,
the panels define the tested simulator geometry envelope and its named transfer
axes. %
}

\newcommand{\CapScale}{%
Homogeneous swarm-size extension of the directional operating envelope. Panel A
shows the paired completion gain against team size for two, four, six and eight
airframes at the three tested force levels; points are means over the same 30
paired seeds per cell and bars are 95~percent $t$ intervals. Raw exact signed-rank
values and their Holm adjustments are reported over the 12 size--force cells.
Panel B opens the transition force at \KneeLevel{}~N and shows the wind-agnostic and
wind-aware completion endpoints separately, with means and 95~percent intervals over
the same seed pairs. The constant-profile four-airframe slice contains 180 rows:
30 seeds, two arms and three force levels with the causal estimator. The ten-seed
exact-vector condition is a distinct oracle slice, and the other team sizes are
declared simulator cells. Together the panels define the measured homogeneous-size
transfer envelope. %
}

\newcommand{\CapRadius}{%
Sensitivity of the operating envelope to the capture radius that defines task
completion, over \NRadiusEpisodes{} episodes at radii of \RadiusGrid{}~m and
\NRadiusSeeds{} paired seeds per radius and force. Panel A gives the paired
completion advantage of the wind-aware over the wind-agnostic allocator against
steady horizontal wind force, one series per radius, with bars showing 95~percent
intervals over the paired seeds and the series offset horizontally so that
overlapping bars remain separable. The labelled point on each series is that
radius' largest advantage. Panel B gives the two completion ladders those
advantages are differences of, one facet per radius, with the dashed rule marking
the same labelled force. The transition survives every tested radius at a
comparable size, between $+\RadiusPeakDeltaMin{}$ and $+\RadiusPeakDeltaMax{}$
completion, and its location moves with the radius, from
\RadiusTightPeakForce{}~N at the tightest radius to \RadiusWidePeakForce{}~N at
the widest, so the force at which directional information pays is a property of
the operating point. Within this \RadiusHolmFamily{}-test validation family
\RadiusHolmSurvivors{} cell clears a Holm-corrected 5~percent threshold, a
resolution set by the \NRadiusSeeds{} seeds per cell; the family is a
sensitivity readout beside the larger primary sweep. %
}

\newcommand{\CapPathwaySafety}{%
Complete pathway decomposition and the descriptive safety diagnostics, from two
separate campaigns of \NFactorialEpisodes{} and \NSafetyEpisodes{} episodes.
Panel A gives all three terms of the assignment-by-feedforward factorial as
paired completion effects against true horizontal force, one facet per temporal
profile, with bars showing 95~percent intervals over \NFactorialSeeds{} paired
seeds and the terms offset horizontally for legibility. The target-feedforward
main effect carries the whole gain at the transition force, between
$+\FactorialFeedforwardMin{}$ and $+\FactorialFeedforwardMax{}$ completion across
profiles; the assignment main effect and the interaction stay within
\FactorialOtherAbsMax{} of zero and neither clears its Holm-corrected threshold at
any force. Panel B is descriptive only and carries no test: for each allocator and
force it gives the range of the per-episode minimum pairwise airframe separation
over \NSafetySeeds{} seeds, with the point at the mean of that range and the
dashed rule at the closest airframe pair recorded anywhere in the campaign,
\SafetyWorstSeparation{}~m. The panel is a benchmark-level feasibility diagnostic;
because target placement and arrival status contribute to the metric, it is not
interpreted as an allocator-causal clearance effect. %
}

\newcommand{\CapEvidence}{%
Cross-campaign evidence synthesis for the directional operating envelope. Panel A
places the transition paired-completion estimates from the primary, estimate-noise,
crossed-layout, causal-temporal and homogeneous-size views on a common rate axis;
bars are the view-specific 95~percent intervals and colours identify the evidence
block. The crossed-layout estimate is shown with its interval because that record
uses a layout-cluster bootstrap. Panel B shows the episode-row volume of each of the
eight views on a logarithmic axis, with the text next to each point naming the seed
or layout-cluster unit that supports its analysis; pathway ablation and gain
sensitivity are two views of one pathway/gain synthesis block. The figure is
descriptive by design and preserves each view's inferential unit. %
}

\newcommand{\CapIntegralGain}{%
Paired completion gain at the \KneeLevel{}~N force against the baseline integral
gain, from the controller-identification sweep at the stock
\EpisodeHorizonS{}~s horizon. Points are the seed-paired aware-minus-agnostic
completion difference; bars are 95~percent $t$ intervals. Clamp settings of
\IntegralClampXy{}, 4 and 8~m/s are drawn as distinct markers at the same Ki,
so a clamp that leaves the 0.15~N completions unchanged sits on top of the
stock marker. Doubling the stock Ki of \StockKi{}~N/m/s collapses the gain from
\AuditStockGain{} to \AuditKiTenGain{} (interval
[\AuditKiTenCiLo{}, \AuditKiTenCiHi{}]). %
}

\newcommand{\CapHorizon}{%
Horizon sweep under the stock controller (Ki~\StockKi{}~N/m/s, clamp
\IntegralClampXy{}~m/s). Panel A is the paired completion gain at
\KneeLevel{}~N against episode duration, with bars giving the 95~percent $t$
interval; the gain is \AuditGainTSix{} at the stock \EpisodeHorizonS{}~s
horizon and \AuditGainTTen{} once the horizon reaches 10~s. Panel B is the
matching wind-agnostic completion at the same force, which rises as the
episode lengthens and the slow integral is given more time to wind up. Both
panels share one horizon axis. %
}

\newcommand{\CapLambdaSweep}{%
Pre-declared sweep of the assignment penalty weight
$\lambda\in\{\LambdaGrid{}\}$~m/N under the stock controller at the
\EpisodeHorizonS{}~s horizon, \LambdaAwareEpisodes{} aware-arm episodes per
weight. Circles count episodes whose initial ($t=0$) assignment differs from
the Euclidean assignment; squares count episodes in which any call,
initial or replan, differs. The dashed line marks the inherited weight
$\lambda=\LambdaReference{}$, at which the counts reproduce the bound
assignment audit exactly. Initial divergence is zero through
$\lambda=\LambdaInertMax{}$ and first appears at
$\lambda=\LambdaFirstDivergence{}$; the horizontal axis is pseudo-logarithmic
so that $\lambda=0$ sits on the same axis. The sweep is descriptive by
declaration and carries no confirmatory test. %
}

%% Numbered environments, shared by every binding of this paper for the same
%% reason the body is: a statement that is Proposition 2 in one binding and
%% Proposition 3 in another cannot be cross-referenced by a reader holding the
%% other one.
%% Equation labels remain authored in methods.tex and are checked for uniqueness;
%% proofs use amsthm's automatic QED marker with no hand-added square.

\theoremstyle{plain}
\newtheorem{proposition}{Proposition}
\newtheorem{corollary}{Corollary}

\theoremstyle{definition}
\newtheorem{definition}{Definition}

\newcommand{\generatedtable}[1]{\input{#1}}

\title{From directional sensing to task completion: the operating envelope of
target-position feedforward against a slow integral in wind-disturbed quadrotor
waypoint capture}
\ifdefined\ANONYMOUS
\author{Author and affiliation withheld for blind review}
\else
\author{Zeyu Fu\\
Army Medical University, Chongqing, China\\
\texttt{fuzeyu99@126.com}\\
ORCID: 0009-0001-8329-0108}
\fi
\date{}

\begin{document}
\maketitle

\begin{abstract}
%% One paragraph, self-contained. No cross-references, section numbers,
%% structural labels or citations.
%% Every numerical value is emitted from the hash-bound evidence pass.

Directional disturbance information can change both which aircraft receives a
waypoint and how it flies there. We map and quantify a reproducible computational
operating envelope for this information-to-completion interface in a
multi-quadrotor allocation--control stack. The system is a Hungarian
allocation--control stack in a rigid-body simulator: a wind-agnostic arm uses
Euclidean assignment and the stock position controller, while a wind-aware arm
adds an upwind assignment cost and target-position feedforward. Shared seeds,
paired contrasts and a seven-level force ladder define the primary envelope, and
estimate-noise, temporal, pathway, gain, geometry, swarm-size, capture-radius
and feasibility campaigns extend the same endpoint. The envelope has a sharp
transition at \KneeLevel{}~N (\KneePctHover{}~percent of hover thrust), where
completion rises from \KneeAgnostic{} to \KneeAware{} (paired gain
+\KneeDelta{}, surviving the pre-specified signed-rank family correction). Relative
noise through \BestSigma{} preserves the transition, and causal estimates retain
positive gains under constant and gusting profiles. A complete
assignment-by-feedforward factorial attributes the gain to target feedforward,
while paired gains stay positive from \ScaleMinSize{} through \ScaleMaxSize{}
airframes. Rerunning the ladder across capture radii keeps a transition of
comparable size but moves it from \RadiusTightPeakForce{} to
\RadiusWidePeakForce{}~N, so force describes an operating point rather than an
interface constant. The study therefore contributes a simulator operating
envelope, a transparent computational validation contract, and a parameterized
transfer protocol for hardware and field validation; clearance and collision
quantities are retained as benchmark-level feasibility diagnostics rather than
as deployment claims.

\end{abstract}

\paragraph{Keywords} disturbance feedforward, integral control, wind estimation, thresholded endpoint, task allocation, quadrotor swarm, reproducibility

%% The manuscript body, shared by every binding of this paper.
%% The class, the front matter and the bibliography style belong to the
%% binding; everything a reader reads belongs here, once.
\section{Introduction}
A quadrotor that knows which way the wind is blowing can fly an assigned waypoint
with a more useful target bias. At the
vehicle level, directional disturbance information is now estimated, modelled and
used for control: unknown-input filters have been validated for wind-gust
compensation~\cite{azid2021windgust}, aerodynamic tracking laws have been tested
under wind perturbations~\cite{perozzi2018wind}, and onboard wind measurement and
simulation practices have been systematized~\cite{abichandani2020wind,wi2024windstress,
wetz2021distributed,schiano2014wind,gonzalezrocha2019sensingwind,
wang2018windestimation}.
Wind-aware planners likewise make direction and reachability explicit in their
costs~\cite{thanellas2019spatialwind,heidari2021windplanning,techy2009minimum};
recent work extends that line to flow prediction, safety and formation adaptation
~\cite{seal2025,wespr2026,knodedw2025}. A natural team-level hypothesis is that
the same directional information has operational leverage when several aircraft
compete for tasks, because it can enter an upwind assignment cost. This paper
puts that hypothesis to a paired test on a multi-quadrotor benchmark and
identifies where the directional information acts: at the inherited penalty
weight the initial assignment coincides with the Euclidean solution in every
audited episode, and the completion gain travels through target-position
feedforward compensating a slow stock integral.

The instance is a multi-quadrotor stack in which \NDrones{} airframes must reach
\NDrones{} waypoints. A wind-agnostic arm uses Euclidean assignment and the stock
position controller; a wind-aware arm can add an upwind assignment cost and
target-position feedforward. Shared seeds, paired contrasts and a
\NWindLevels{}-level force ladder define the primary envelope. The contribution
is a scoped identification of the mechanism: against the stock slow integral
(\StockKi{}~N/m/s) and the \EpisodeHorizonS{}~s horizon, target-position
feedforward changes short-episode completion, and the campaigns map the
operating envelope in which it does so.
Initial assignment is inert at the inherited $\lambda=0.25$ and, on a
pre-declared \LambdaGridSize{}-point sweep, at every tested weight through
$\lambda=\LambdaInertMax{}$~m/N; it first diverges at
$\lambda=\LambdaFirstDivergence{}$~m/N. Doubling that stock Ki collapses the 0.15~N paired advantage to
\AuditKiTenGain{}. A complete assignment-by-feedforward factorial attributes the
gain to the feedforward path. Allocation taxonomies and aerial-swarm
surveys~\cite{gerkey2004formal,dias2006market,korsah2013taxonomy,chung2018swarm,
skaltsis2023uavallocation} frame the team-level allocation question that the
benchmark puts to the test.
Panels A and B of Fig.~\ref{fig:system-overview} make the measured object
explicit: the same directional input enters assignment and feedforward, and the
paired endpoint keeps thresholded completion beside closest approach. Panel A
is a drawn schematic of the research object.

\begin{figure*}[!t]
\centering
\includegraphics[width=\linewidth]{fig0_system_overview.pdf}
\caption{\CapSystemOverview}
\label{fig:system-overview}
\end{figure*}

The study follows the information path from capability to mechanism. It first
locates the transition and tests estimate-noise robustness, then resolves scene
context with seed-level repeats and a crossed heading--layout design. A continuous
tracking endpoint shows how the thresholded gain is formed. Temporal profiles,
pathway and feedforward-gain controls, scene-geometry factors and a homogeneous
team-size axis then open the same estimand along progressively more operational
dimensions. The exact-vector oracle is the reference for every extension, so the
\NEvidenceBlocks{} synthesis blocks form one measured operating envelope with a
continuous information-to-completion trace. The result is a computational
baseline: the tested simulator envelope is the calibration reference against
which a hardware extension is to be validated.

\section{Methods}
%% Independent methods file. Not woven into main.tex prose.
%% Every quantity described here resolves to a hash-bound evidence entry.

\subsection{Simulation and disturbance setup}

All episodes run in an open-source rigid-body quadrotor simulator built on a
general-purpose physics engine~\cite{panerati2021gym,coumans2016pybullet,
ozturk2025mrtasim}. The
rigid-body and state variables follow the standard multirotor modeling
convention~\cite{mahony2012multirotor}; the airframe is the simulator's stock
nano-quadrotor model: \DroneMass{}~g, with a
per-airframe hover thrust of \HoverThrust{}~N. An episode places \NDrones{} of these
airframes and \NDrones{} randomly drawn waypoints in a shared arena and runs for six
simulated seconds with physics integrated at 240~Hz and control issued at 48~Hz.
Everything runs headless on one CPU core; no graphics or accelerator is involved.

The physical and control grids are fixed throughout the study by
Eq.~\eqref{eq:time-grid}:
\begin{equation}
  T=6~\mathrm{s},\qquad f_{\mathrm{phys}}=240~\mathrm{Hz},\qquad
  f_{\mathrm{ctrl}}=48~\mathrm{Hz},\qquad
  \Delta t_{\mathrm{ctrl}}=f_{\mathrm{ctrl}}^{-1} .
  \label{eq:time-grid}
\end{equation}
Each control action is held for five rigid-body integration steps, and every
campaign uses this same horizon and sampling contract.

The simulation substrate intentionally isolates the rigid-body dynamics of the
allocation--control interface. Four physical and aerodynamic phenomena are
deliberately unmodeled in this computational baseline:
(1) inter-rotor and inter-vehicle aerodynamic downwash or turbulent wake vortex
interactions;
(2) ground and ceiling proximity aerodynamic effects;
(3) blade flapping dynamics and translational rotor drag at forward airspeed; and
(4) motor thermal torque limits and battery voltage sag under sustained counter-wind
thrust.
The empirical transition window (\KneeLevel{}~N, or \KneePctHover{}\% of hover thrust)
is therefore characterized strictly as the operating envelope of the declared
rigid-body simulation contract, serving as an uncorrupted computational
benchmark rather than an uncalibrated claim of open-air flight robustness.

The primary endpoint is the completion rate. An airframe is marked complete at its
first entry into the capture ball, rather than only at the final time. For airframe
$i$ with active assignment $a_i(t)$, the completion time, rate and continuous
closest-approach endpoint are defined in Eq.~\eqref{eq:endpoints}:
\begin{equation}
  \begin{aligned}
  d_i(t)&=\lVert\mathbf{x}_i(t)-\mathbf{y}_{a_i(t)}\rVert_2,\\
  \tau_i&=\inf\{t\in[0,T]:d_i(t)\le r_c\},\\
  C&=\frac{1}{N}\sum_{i=1}^{N}\mathbb{1}(\tau_i\le T),\qquad
  d_i^{\min}=\min_{0\le t\le T}d_i(t) .
  \end{aligned}
  \label{eq:endpoints}
\end{equation}
Equation~\eqref{eq:endpoints} defines the thresholded and continuous endpoints
used throughout the study. Here $N=\NDrones{}$ and $r_c=\CaptureRadius{}$~m. Because there are four
airframes, each episode takes one of five discrete completion values. The
continuous endpoint remains informative when the thresholded rate is pinned at
zero and explains how a small tracking shift can move an airframe across the
capture boundary.

Each seed $s$ initializes one reproducible scene stream. With
$G=\lceil\sqrt{N}\rceil$, the launch positions and raw waypoint draws are
generated according to Eq.~\eqref{eq:scene-draw}:
\begin{equation}
  \begin{aligned}
  x^{0}_{i,x}&=0.4(i\bmod G)-0.2,&
  x^{0}_{i,y}&=0.4\lfloor i/G\rfloor-0.2,&
  x^{0}_{i,z}&=0.1,\\
  \mathbf{y}^{0}_j&\sim\mathcal{U}(\mathcal{B}),&
  \mathcal{B}&=[-0.8,0.8]^2\times[0.4,0.9].
  \end{aligned}
  \label{eq:scene-draw}
\end{equation}
Equation~\eqref{eq:scene-draw} fixes the launch grid, waypoint support and seeded
draw for $i,j=0,\ldots,N-1$. Coordinates remain in metres and forces in newtons;
the preprocessing stage applies no image or learned transformation.

The scene-geometry extension is a deterministic post-draw transform. Let
$\mathbf{c}_x$ and $\mathbf{c}_y$ be the launch and raw-target horizontal
centroids, let $\rho\in\{1,1.30\}$ denote nominal or expanded spread, and let
$\chi\in\{0,1\}$ denote aligned or crossed orientation. The transformed horizontal
target coordinates are defined by Eq.~\eqref{eq:geometry-transform}:
\begin{equation}
  \begin{aligned}
  \mathbf{u}_j&=\mathbf{c}_y+\rho(\mathbf{y}^{0}_{j,xy}-\mathbf{c}_y),\\
  \mathbf{v}_j&=(1-\chi)\mathbf{u}_j+\chi\{\mathbf{c}_x-(\mathbf{u}_j-\mathbf{c}_x)\},\\
  \mathbf{y}_{j,xy}&=\mathbf{v}_j+\mathbb{1}(\text{far})[0.22,-0.18]^{\top},
  \qquad y_{j,z}=y^{0}_{j,z} .
\end{aligned}
  \label{eq:geometry-transform}
\end{equation}

The transform consumes no additional random numbers, making the
nominal--near--aligned cell an exact scene anchor. The wind vector applied to
the rigid-body solver during control interval $k$ is specified by
Eq.~\eqref{eq:force-vector}:
\begin{equation}
  \mathbf{w}_k=m_k[\cos\theta_k,\;\sin\theta_k,\;0]^{\top},
  \qquad t_k=k\Delta t_{\mathrm{ctrl}},
  \label{eq:force-vector}
\end{equation}
The value is held over the five physics sub-steps, so scene, heading, phase and
treatment pairing are explicit before evaluation.

The simulator supplies the rigid-body dynamics; we add a transparent force
profile to expose the allocation--control interface while complementing established
wind-measurement, identification and
disturbance-observer practice~\cite{abichandani2020wind,wi2024windstress,
chen2016dobc,rodriguezmata2018highgain,wang2021esoquadrotor}. At every physics
sub-step, the horizontal force in
Eq.~\eqref{eq:force-vector} is applied to each airframe through the engine's
external-force interface. For the constant profile, $m_k=m$ and
$\theta_k=\theta$; for temporal profiles the profile advances only at the
control grid. This is a real force entering the rigid-body solver before the state
is logged.

A trajectory-validation checkpoint identified and corrected the force-application
reference point before the reported campaigns were admitted. The engine's
world-frame external-force call expects an absolute world point; using the live
centre-of-mass position at every sub-step therefore applies a pure translational
disturbance without an unintended moment arm. A raw-trajectory check confirmed
the correction through the expected attitude and displacement response. Every
number reported here comes from the validated implementation, and the trajectory
check is retained as a reusable gate for external-force experiments.

\subsection{Evaluation protocol}

Every evaluation preserves the same simulator, seeded scene and endpoint while
changing one declared information or geometry factor at a time. The primary
comparison is the paired aware--agnostic contrast in Eq.~\eqref{eq:paired-estimand};
the support campaigns reuse that unit so their results can be read on the same
operating-envelope scale.

The fixed simulator and endpoint contract is summarized in
Table~\ref{tab:parameter-contract}; the opened factors and their inferential
units are separated into the campaign design contract in
Table~\ref{tab:campaign-contract}. Keeping these roles distinct makes the first
table a compact reproduction anchor and the second a map of how the evidence
expands. Neither table is a second source of numerical results: values reported
in the Results remain programmatically generated from the evaluated campaign data.
\input{parameter_and_compatibility_table}
\input{campaign_design_table}

\subsection{Allocation and control}

The wind-agnostic arm assigns airframes to waypoints by solving a linear
assignment problem on Euclidean distance with the Hungarian algorithm
~\cite{kuhn1955hungarian,gerkey2004formal,dias2006market,choi2009cbba}, and flies
the resulting targets with one standard proportional-integral-derivative (PID)
position controller per airframe.

The controller gains, motor model, pulse-width-modulation to rotor-speed (PWM/RPM) conversion and saturation path are
unchanged from the vendored simulator implementation and are listed in the
parameter contract. This stock position layer is consistent with the standard
rigid-body tracking formulation for quadrotors~\cite{lee2010se3}; the present
study varies the information supplied to the allocation--control interface while
holding the underlying stabilizing controller fixed.

For the wind-aware arm, let $\widehat{\mathbf w}_k$ be the supplied estimate and let
$\mathbf{u}_{ij,k}$ point from airframe $i$ to waypoint $j$. The assignment cost
and target used by the two arms are defined in Eq.~\eqref{eq:assignment-cost}:
\begin{equation}
  J_{ij,k}=\lVert\mathbf{y}_j-\mathbf{x}_{i,k}\rVert_2
  +\lambda\max\!\left(0,-\mathbf{u}_{ij,k}^{\top}\widehat{\mathbf w}_k\right),
  \qquad
  \mathbf{u}_{ij,k}=\frac{\mathbf{y}_j-\mathbf{x}_{i,k}}
  {\lVert\mathbf{y}_j-\mathbf{x}_{i,k}\rVert_2},
  \label{eq:assignment-cost}
\end{equation}
Equation~\eqref{eq:assignment-cost} adds the directional penalty to the Euclidean
distance while retaining a normalized direction vector for each candidate pair.

The assignment returned at control step $k$ is the one-to-one minimizer in
Eq.~\eqref{eq:assignment}:
\begin{equation}
  \begin{aligned}
  \mathbf{Z}_k&=\underset{\mathbf{Z}\in\mathcal{P}_N}{\arg\min}\;
  \sum_{i=1}^{N}\sum_{j=1}^{N}J_{ij,k}z_{ij,k},\\
  \mathcal{P}_N&=\left\{\mathbf{Z}\in\{0,1\}^{N\times N}:\;
  \sum_jz_{ij}=1,\;\sum_iz_{ij}=1\right\}.
  \end{aligned}
  \label{eq:assignment}
\end{equation}
Equation~\eqref{eq:assignment} applies the one-to-one constraints at each
re-planning step. Here $z_{ij,k}=1$ records that airframe $i$ takes waypoint $j$; the Hungarian
algorithm solves this minimization at each re-planning step. This comparator
follows the assignment--trajectory decomposition used in cooperative UAV studies
~\cite{beard2002target,turpin2014goal} and in steady-wind task/path formulations
~\cite{luo2018integrated}. It retains the same scene, stock controller, horizon
and assignment solver while changing one information input, so the measured
contribution is the paired endpoint change under a controlled contrast.

The target-position feedforward rule is defined separately for the two arms in
Eq.~\eqref{eq:target-feedforward}:
\begin{equation}
  \mathbf{r}_{j,k}=\mathbf{y}_j-\gamma\widehat{\mathbf w}_k
  \quad\text{(aware)},
  \qquad
  \mathbf{r}_{j,k}=\mathbf{y}_j
  \quad\text{(agnostic)},
  \label{eq:target-feedforward}
\end{equation}
Equation~\eqref{eq:target-feedforward} defines the target supplied to the stock
position controller in each arm. Here $\lambda=0.25$ in the aware arm and $\gamma$ is the tested target
feedforward gain (the inherited value is \FeedforwardGainDefault{}~m/N). The
assignment is re-solved once per
simulated second over unfinished pairs; the replanning count is defined by
Eq.~\eqref{eq:replanning}:
\begin{equation}
  R=\sum_{k\in\mathcal{K}}\mathbb{1}(A_k\ne A_{k-1}),
  \qquad \mathcal{K}=\{48,96,144,192,240\},
  \label{eq:replanning}
\end{equation}
Equation~\eqref{eq:replanning} counts assignment changes on the declared control
steps. Here $A_k$ is the active assignment. The penalty weight and feedforward gain
are fixed for the primary envelope and opened only in the declared sensitivity
campaign.

A causal estimator starts from the zero vector, consumes the current measurement
and its prior state, and updates at 48~Hz; the controller clips commands through
the stock simulator saturation path. The aware arm therefore receives a
current-and-past signal only.

\subsection{Paired design and metrics}

Both arms are run on the same \NSeeds{} seeds at each disturbance level, and a seed
fixes the waypoint layout and the disturbance heading. The comparison is therefore
paired within seed, not a comparison of independent samples. Across
\NWindLevels{} disturbance levels, two arms and \NSeeds{} seeds, the primary sweep is
\NEpisodesOracle{} episodes.

The paired seed is the inferential unit for every primary and support comparison;
episode rows are retained for audit and visualization but are not treated as
independent samples. The exact test, interval construction and multiplicity rule
are stated once in the Statistical analysis subsection below.

For clarity, the estimand at force $w$ is the seed-level paired contrast defined
in Eq.~\eqref{eq:paired-estimand}:
\begin{equation}
  \Delta_s(w)=C_{s,\,\mathrm{aware}}(w)-C_{s,\,\mathrm{agnostic}}(w),
  \qquad
  \bar\Delta(w)=\frac{1}{S}\sum_{s=1}^{S}\Delta_s(w),
  \label{eq:paired-estimand}
\end{equation}
where $C$ is Eq.~\eqref{eq:endpoints} and $S$ is the number of paired seeds.
The continuous companion is summarized by the seed-level mean of $d_i^{\min}$
from Eq.~\eqref{eq:endpoints}. For a supplied wind estimate, the diagnostic
reported beside the task endpoint is the relative vector error defined by
Eq.~\eqref{eq:relative-error}:
\begin{equation}
  e_k=\frac{\lVert\widehat{\mathbf w}_k-\mathbf w_k\rVert_2}
  {\max(\lVert\mathbf w_k\rVert_2,\,\epsilon)},
  \qquad \epsilon>0,
  \label{eq:relative-error}
\end{equation}
Equation~\eqref{eq:relative-error} defines the relative vector-error diagnostic.
Temporal summaries average $e_k$ over the control grid before averaging over
seeds. The replanning count is the event metric in Eq.~\eqref{eq:replanning}; it
is descriptive and is never substituted for the primary endpoint.

\subsection{Statistical analysis}

The intervals in the force, temporal, ablation and gain tables are $t$ intervals
on the seed-level contrasts in Eq.~\eqref{eq:paired-estimand}. The resulting
95~percent interval and its seed-level variance are given by
Eq.~\eqref{eq:paired-ci}:
\begin{equation}
  \begin{aligned}
  \mathrm{CI}_{95\%}(w)&=\bar\Delta(w)\pm
  t_{0.975,S-1}\frac{s_{\Delta}(w)}{\sqrt{S}},\\
  s_{\Delta}^2(w)&=\frac{1}{S-1}\sum_{s=1}^{S}
  [\Delta_s(w)-\bar\Delta(w)]^2 .
  \end{aligned}
  \label{eq:paired-ci}
\end{equation}
Each level is also tested with an exact two-sided Wilcoxon signed-rank
test~\cite{wilcoxon1945individual} implemented in a standard scientific computing
library~\cite{virtanen2020scipy}, with zero differences dropped. Where every
paired difference at a level is exactly zero the signed-rank statistic is
undefined; the table records a dash, distinct from a $p$ value of one. To
guard against selecting the best of \NWindLevels{} levels after the fact, in
the spirit of family-wise multiple-comparison control~\cite{dunn1961multiple},
the threshold is defined by Eq.~\eqref{eq:bonferroni}:
\begin{equation}
  \alpha_{\mathrm{Bonf}}=\frac{0.05}{\NWindLevels{}}
  =\BonferroniAlpha{}.
  \label{eq:bonferroni}
\end{equation}

The exact signed-rank values use the same contrasts with zero differences removed;
the $t$ interval is reported alongside the exact test for the seed-level families.
The revised support record keeps its inferential unit explicit: crossed-layout
intervals use an ordinary bootstrap over layout clusters, whereas geometry and
scale tests use their declared seed-level contrasts. Episode rows and bootstrap
draws are never counted as independent experimental units.

For a support family with $m$ testable contrasts, we use a deterministic Holm
step-down adjustment. If $p_{(r)}$ is the $r$th raw value after sorting by
increasing $p$ (ties are broken by the contrast identifier), its adjusted value is
defined by Eq.~\eqref{eq:holm}:

\begin{equation}
  p^{\mathrm{Holm}}_{(r)}=
  \min\!\left\{1,\;\max_{1\leq j\leq r}(m-j+1)p_{(j)}\right\},
  \qquad r=1,\ldots,m .
  \label{eq:holm}
\end{equation}
This deterministic step-down value is used for each named support family.

To report what the geometry design can resolve, we add a minimum-detectable-effect
(MDE) planning diagnostic. It is calculated on the declared seed-level
inferential units, uses the familywise level $\GeometryMDEAlpha{}$ for the
\GeometryMDEFamily{} geometry contrasts and
does not constitute an equivalence test. For a two-sided familywise level
$\alpha_F$, target power $\pi$, $n$ paired units and observed contrast standard
deviation $s$, the planning approximation is defined in Eq.~\eqref{eq:mde}:
\begin{equation}
  \operatorname{MDE}_{\pi}=
  \left[t_{1-\alpha_F/2,\,n-1}+\Phi^{-1}(\pi)\right]\frac{s}{\sqrt{n}} .
  \label{eq:mde}
\end{equation}
The MDE is a sensitivity calculation rather than a claim that effects below it
are absent; no equivalence margin was prespecified. It records the resolution of
the executed design and identifies the effect range for a follow-up campaign.

The raw exact value and the corresponding adjusted value are retained together in
the recorded result table. The primary seven-force family remains governed by the
predeclared Bonferroni threshold above; Holm is used only for the explicitly named
support families in the campaign contract.

The seven-force ladder is the primary confirmatory family. Estimate-noise,
temporal, pathway, gain, crossed, geometry and scale analyses are separately
named support or descriptive families; their correction rules and cluster units
are printed in Table~\ref{tab:campaign-contract}. This keeps the unit of
inference explicit and prevents the episode count from being mistaken for an
independent-sample size. This simulation design follows reporting practice that
separates the data-generating design, estimand and analysis family
~\cite{morris2019simulation}. Every generated table is re-derived from its
recorded result file, checked for its declared row count and pairing, and emitted by the
same R pass that emits the manuscript macros; the prose reads those values directly.

For the estimate-noise table, the complete disturbance-by-noise family contains
15 comparisons and uses the declared threshold 0.0033. The 0.30 and 0.50 rows
cross that threshold; the 0.10 row is retained as a positive descriptive estimate,
and the 1.00 row has exact $p=\MaxSigmaP{}$ and is reported without a
multiplicity-adjusted significance claim. The crossed, geometry and scale
intervals use their declared layout or seed-level units, and their support status
is reported separately from the primary force-family claim.

\subsection{Information quality and temporal profiles}

The two arms above both give the wind-aware controller the exact disturbance
vector. This is the key information condition in the reference experiment, and
the manipulation below relaxes it. The wind-aware arm is re-run with the
disturbance vector it receives corrupted by relative magnitude noise
$\sigma/|w| \in \{\SigmaGrid\}$ and matched angular noise, redrawn
independently at every 48~Hz control step. In this independent-noise condition,
the supplied estimate is generated from the true magnitude $m_k$ and heading
$\theta_k$ according to Eq.~\eqref{eq:independent-noise}:
\begin{equation}
  \begin{aligned}
  \widetilde m_k&=\max\{0,m_k(1+\sigma\xi_k)\},&
  \widetilde\theta_k&=\theta_k+\sigma\eta_k,\\
  \widehat{\mathbf w}_k&=\widetilde m_k
  [\cos\widetilde\theta_k,\sin\widetilde\theta_k,0]^{\top},&
  \xi_k,\eta_k&\stackrel{\mathrm{iid}}{\sim}\mathcal{N}(0,1).
  \end{aligned}
  \label{eq:independent-noise}
\end{equation}
The magnitude and angular perturbations are independent at each control step.
The manipulation is applied at the three disturbance levels that bracket the
region of interest, giving \NEpisodesNoise{} episodes.

Because the wind-agnostic arm never reads the disturbance estimate and is
deterministic given a seed, its recorded outcomes serve as the fixed reference for
every noise condition. A regression check confirms this: at zero estimate noise the
re-run reproduces the original completion of \SigmaZeroAware{} exactly.

Two aspects of this manipulation are recorded with the data. First, the decision
point fixed before the run was $\sigma/|w| = \PreregSigma{}$; the nearest executed
grid level is \NearestSigma{}, so the pre-set decision is evaluated on that
grid-aligned condition. Second, redrawing the noise
independently at each control step is a controlled first robustness layer: it is
zero-mean around the truth and a feedforward term integrating over roughly three
hundred control steps per episode can average much of it away. The temporal
extension below makes bias and correlation explicit and tests how the operating
envelope responds when the estimate has memory.

The next experiment makes the disturbance and its estimate evolve together. Wind
gust estimation and wind-perturbation tracking studies motivate treating the
estimator state and vehicle response as separate observable quantities
~\cite{azid2021windgust,perozzi2018wind}. Three profiles are used: a constant
force, a slow sinusoidal gust and a faster sinusoidal gust. With $P$ the selected
period and phase $\phi$, the centred profile is defined by
Eq.~\eqref{eq:temporal-profile}:
\begin{equation}
  \begin{aligned}
  q(t)&=\sin(2\pi t/P+\phi)-\sin\phi,\\
  m(t)&=m_0[1+a_mq(t)],\\
  \theta(t)&=\theta_0+a_\theta q(t).
  \end{aligned}
  \label{eq:temporal-profile}
\end{equation}
The slow profile uses
$(a_m,a_\theta)=(0.18,0.22)$, the fast profile uses $(0.12,0.16)$, and the
constant profile sets both amplitudes to zero. Substitution into
Eq.~\eqref{eq:force-vector} gives the true force
presented to the simulator. The estimator is a first-order causal filter with
configurable response time, angular bias, magnitude bias and measurement noise.
Writing $b_m$ for the magnitude bias, $b_\theta$ for the angular bias and
$\boldsymbol\nu_k$ for zero-mean measurement noise, the current measurement is
defined by Eq.~\eqref{eq:causal-measurement}:
\begin{equation}
  \mathbf{z}_k=R(b_\theta)\,m_k(1+b_m)[\cos\theta_k,\sin\theta_k,0]^{\top}
  +\boldsymbol\nu_k,
  \label{eq:causal-measurement}
\end{equation}
where $R(b_\theta)$ rotates the horizontal components by $b_\theta$. Its state
update is given by Eq.~\eqref{eq:causal-estimator}:
\begin{equation}
  \widehat{\mathbf w}_k=(1-\alpha)\widehat{\mathbf w}_{k-1}+\alpha\mathbf{z}_k,
  \qquad
  \alpha=\begin{cases}1,&\tau=0,\\
  1-\exp(-\Delta t_{\mathrm{ctrl}}/\tau),&\tau>0,
  \end{cases}
  \label{eq:causal-estimator}
\end{equation}
Here $\tau=0$ is the exact-vector oracle condition. The filter uses the current
measurement and its previous state; this causal vector is the signal supplied to
the wind-aware arm.


The named presets are fixed before reading the outcomes: the causal condition has
response time $\tau=0.18$~s, zero angular and magnitude bias, and 0.04 measurement
noise fraction; the biased-causal condition has $\tau=0.32$~s, $12^\circ$ angular
bias, 0.06 magnitude bias and 0.08 noise fraction. Both states start at the zero
vector, use the same 48~Hz control grid, and draw their independent measurement
stream from a seed-derived generator. These are controlled simulator estimators
designed to expose timing, bias and noise effects before a physical-sensor replay.


For each profile we use the shoulder and transition forces (0.10, 0.15 and
0.20~N) and \NTemporalSeeds{} paired seeds. The baseline episode and every wind-aware condition
share the seed, waypoint layout, disturbance phase and simulator configuration. The
three estimator conditions are an exact-vector oracle, a causal estimate and a
causal estimate with controlled bias. The resulting campaign contains
\NTemporalEpisodes{} episode records, with the baseline retained once per paired
condition.

Each record stores the true and estimated vectors, relative magnitude error, angular
error and profile metadata. The constant/oracle subset is also a regression of the
legacy control path over \NTemporalLegacyRows{} rows; it must remain identical before
the temporal results are admitted to the manuscript.

\subsection{Pathway, geometry and team-size campaigns}

The wind-aware stack has two directional pathways: an upwind assignment cost and a
target-position feedforward term. To identify the pathway carrying the measured task
gain, we repeat the temporal design with the feedforward term selectively disabled
while retaining the directional assignment cost. This support campaign contains
\NAblationEpisodes{} paired records and uses the same profiles, forces, seeds and
estimator states as the primary campaign. Its allocation-only contrast is read
alongside the complete-arm contrast to localize the pathway while retaining the
same endpoint.

We then vary only the target-position feedforward gain in a fixed seven-point grid
from zero through 1.40~m/N, including the inherited \FeedforwardGainDefault{}~m/N
setting. The causal estimator, profiles, forces, seeds and phases are held fixed,
and every gain condition is paired to the same baseline record. The gain campaign
contains \NGainEpisodes{} episode records. The grid describes the tested
controller envelope; its upper edge is the measured boundary of that grid, and the
inherited setting remains the reference used by the primary temporal comparison.

The seed-level spread and the crossed heading--layout sweep identify scene context
as a first-class axis, so we open that axis with a small, fully crossed campaign.
It contains \NGeometryEpisodes{} episode records: two levels each of target spread,
target range and crossing orientation, crossed with the three force levels around
the transition, \NGeometrySeeds{} seeds and the two allocator arms. The declared
conditions are nominal versus expanded target spread, near versus far target range,
and aligned versus crossed target geometry. The nominal--near--aligned cell is the
legacy scene exactly and serves as the compatibility anchor for the extension.

The expanded spread scales the target cloud by 1.30 about its centroid. The far
condition applies a fixed $[0.22,-0.18]$~m horizontal offset, and the crossed
condition reflects the target cloud through the launch-centroid frame. These
deterministic transformations define the tested operational factor space.
The transforms are applied after the seeded layout is drawn and consume no random
numbers, so the same seed, heading, phase, controller and endpoint are paired
across every cell.

At force $w_k$ and factor $f$, the four settings of the other two geometry factors
are first averaged within each seed. Let $q=1,\ldots,4$ index those settings and
let $s=1,\ldots,S$ index the paired seeds. The seed-level contrast and its grand
mean are defined in Eq.~\eqref{eq:geometry-contrast}:
\begin{equation}
  \begin{aligned}
  d_{s,q,k}(f)&=\Delta_{s,q,k}^{\mathrm{high}}(f)-
    \Delta_{s,q,k}^{\mathrm{low}}(f),\\
  D_{s,k}(f)&=\frac{1}{4}\sum_{q=1}^{4}d_{s,q,k}(f),\\
  \bar D_k(f)&=\frac{1}{S}\sum_{s=1}^{S}D_{s,k}(f),
  \qquad S=30.
  \end{aligned}
  \label{eq:geometry-contrast}
\end{equation}
where $k$ indexes the three force levels and $\Delta$ is the paired estimand in
Eq.~\eqref{eq:paired-estimand}. The manuscript-facing interval and exact
signed-rank test use the 30 values $D_{s,k}(f)$, so one seed contributes one
inferential unit after the crossed settings have been averaged. The 120 terms
$d_{s,q,k}(f)$ remain a named sensitivity readout, and the 1,440 episode rows provide
the full audit record; inferential support remains at the seed level. The campaign is a targeted
support experiment that extends the operating envelope along scene factors while
preserving the primary seven-level sweep.

To test whether the transition is retained when the homogeneous team grows, we
repeat the same paired endpoint on the declared swarm-size grid
\ScaleSizeGrid{} airframes and force grid \ScaleForceGrid{}~N. For each of the
\NScaleSeeds{} seeds, the waypoint stream, disturbance heading, controller,
allocator and six-second horizon are held fixed while the number of airframes and
matched waypoints defines the team-size factor. The wind-aware arm uses the causal
directional estimate; the wind-agnostic arm is the same Euclidean reference.
The scale campaign therefore adds a team-size axis while preserving the definition
of completion and the unit of inference.

For swarm size $n$, force level $w$ and seed $s$, the scale estimand is defined in
Eq.~\eqref{eq:scale-estimand}:
\begin{equation}
  \Delta_{n,w,s}=C^{\mathrm{aware}}_{n,w,s}-C^{\mathrm{agnostic}}_{n,w,s},
  \qquad
  \bar\Delta_{n,w}=\frac{1}{S}\sum_{s=1}^{S}\Delta_{n,w,s},
  \qquad S=\NScaleSeeds{},
  \label{eq:scale-estimand}
\end{equation}
Equation~\eqref{eq:scale-estimand} defines the size--force contrast and its
seed-level mean. The confidence interval and exact signed-rank test are applied to the $S$
seed-level differences at each size--force cell. The disturbance heading is
assigned from the primary seed-indexed heading record, so the comparison tests
team size rather than introducing a second heading draw.

The $n=4$ rows at the constant profile are checked field by field against the
causal temporal campaign's baseline and causal rows for seeds 0--29, forces 0.10,
0.15 and 0.20~N, and both treatments (180 rows in total). The seed sets overlap
exactly; this compatibility regression is reported alongside, and separately from,
the ten-seed oracle row in Table~\ref{tab:band}. The full record contains
\NScaleEpisodes{} episode rows,
with each row retained for audit, continuous miss-distance summaries and the
figure-level operating-envelope map.

\subsection{Scene-context campaigns}

The disturbance heading is a seed-indexed condition. It is drawn once per seed
and then reused at every force level, in both arms, and in the second sweep; the
figure code verifies that one seed carries one heading in both records. This
reuse makes a direct repeat question available: two seeds with nearly the same
heading run the same disturbance twice, and their remaining differences are the
other factors fixed by the seed.

\begin{definition}[Disturbance repeat]\label{def:repeat}
Write an episode as $(a, w, \theta, z)$: the arm $a$, the force magnitude $w$,
the heading $\theta$ on the circle, and $z$ for everything else a seed fixes,
here the waypoint layout. For a tolerance $\varepsilon > 0$, two episodes are a
\emph{disturbance repeat} at $w$ when they share $a$ and $w$ and their headings
satisfy $d(\theta,\theta') \le \varepsilon$, where $d$ is the shorter of the two
arcs between them. The \emph{repeat gap} $D_\varepsilon(w)$ is the largest
difference in completion rate over all disturbance repeats at $w$, taken over
both arms.
\end{definition}

\begin{proposition}[What a repeat measures]\label{prop:repeat}
Suppose completion at $w$ is a function $f_w(a,\theta)$ of the arm and the
heading alone, and that $f_w$ is constant on every arc of width $\varepsilon$.
Then $D_\varepsilon(w) = 0$. Consequently $D_\varepsilon(w) > 0$ identifies a
scene-dependent component in completion at $w$, and its value is a lower bound on
the range over which completion at $w$ varies with $z$.
\end{proposition}

\begin{proof}
Both members of a repeat share $a$ and $w$, and the arc between their headings
has width at most $\varepsilon$, so $f_w$ assigns them the same value and every
repeat difference vanishes. The bound follows because each repeat difference is
realised by two episodes differing only in $z$.
\end{proof}

\begin{proposition}[Saturation-aware repeat gap]\label{prop:vacuous}
If at level $w$ every episode of every arm records the same completion, then
$D_\varepsilon(w) = 0$ irrespective of what completion depends on.
\end{proposition}

\begin{proof}
Every difference taken in Definition~\ref{def:repeat} is a difference of two
equal numbers.
\end{proof}

Proposition~\ref{prop:vacuous} establishes how the repeat gap should be read
alongside the question of whether the seeds differed at all. A saturated ceiling
and a saturated floor both return a gap of zero, so the pair of measures separates
endpoint convergence from scene variation. Of the \NWindLevels{} force levels
swept, the \NInformativeLevels{} variable levels carry the scene-dependence signal,
while the remaining levels define the saturated calibration boundaries.

\begin{proposition}[Tolerance-stable repeat set]\label{prop:tol}
$D_\varepsilon$ depends on $\varepsilon$ only through the set of pairs it
admits. If no pairwise heading separation of the executed design lies in
$(\alpha,\beta]$, then every $\varepsilon \in [\alpha,\beta)$ admits the same
pairs and returns the same $D_\varepsilon$ at every level.
\end{proposition}

\begin{proof}
The admitted set is $\{(i,j) : d(\theta_i,\theta_j) \le \varepsilon\}$, which
changes only as $\varepsilon$ crosses a realised separation.
\end{proof}

The executed design satisfies the hypothesis of Proposition~\ref{prop:tol} with
a margin large enough to make the repeat set stable: \NRepeats{} pairs of
seeds lie within \WidestRepeat{} degrees of one another and the next-closest
pair is \NearestDistinct{} degrees apart, so every tolerance between those two
numbers admits the same \NRepeats{} pairs. We state $\varepsilon = \RepeatTol{}$
degrees and report the interval beside it as the pre-specified operating rule.

\begin{corollary}[Paired units are preserved]\label{cor:notn}
The comparison is paired within seed, so the units of the signed-rank test are
the \NSeeds{} per-seed differences. Two seeds forming a disturbance repeat still
differ in $z$ and remain distinct units, so a positive $D_\varepsilon$ leaves
both the number of pairs and the exactness of the test untouched.
\end{corollary}

The repeat gap and the paired mean answer different questions and should be read
together. The gap measures how much episode outcomes vary at nearly identical
directional input, whereas the paired mean estimates the arm contrast; reporting
both preserves the distinction between scene heterogeneity and treatment effect.

The repeat design localizes variation outside the disturbance while keeping the
waypoint layout and other seed-fixed factors coupled. The crossed heading--layout
sweep below then resolves those factors into named margins and their interaction.

The original record reads the heading drawn by each seed. The designed measurement
uses an explicit heading argument; when none is supplied the routine consumes the
seeded draw, so an assigned heading leaves the waypoint layout at a given seed
exactly as the drawn heading would have left it. Two checks support this contract:
every episode of the original record re-runs to the same numbers, and an episode
assigned the heading its seed would have drawn reproduces that seed's episode
field for field.

On that basis we ran a crossed sweep at the same \NWindLevels{} force levels and
in the same two arms: \NAssignedHeadings{} headings spaced \HeadingStep{}
degrees apart around the full circle, each flown with every one of \NScenes{}
waypoint layouts, for \NEpisodesDesigned{} episodes and \NPairsPerForce{} paired
comparisons at every level. The simulator, controller, allocator, layout
generator, capture radius and episode length remain fixed. This direct coverage
turns the ten-draw context question into a full crossed record while preserving
the original endpoint.

\begin{definition}[Crossed disturbance design]\label{def:crossed}
A design is \emph{crossed} at force $w$ if its episodes at $w$ are indexed by
pairs $(\theta_i, z_j)$ with $i \le H$ and $j \le S$, carrying exactly one
episode per arm at each pair, with the headings $\theta_i$ fixed in advance and
the layouts $z_j$ drawn independently of them. Write
$\delta_{ij}$ for the paired difference in completion at $(\theta_i, z_j)$,
$\bar\delta_{i\cdot}$ and $\bar\delta_{\cdot j}$ for its margins, and
$\bar\delta$ for its mean.
\end{definition}

\begin{proposition}[What crossing separates]\label{prop:crossed}
The corresponding variance decomposition identity in a crossed design is given by
Eq.~\eqref{eq:crossed-decomp}:
\begin{equation}
  \begin{aligned}
  \sum_{i,j}(\delta_{ij}-\bar\delta)^2
  \;=\; S\sum_i(\bar\delta_{i\cdot}-\bar\delta)^2
  \;+\; H\sum_j(\bar\delta_{\cdot j}-\bar\delta)^2
  \\
  \;+\; \sum_{i,j}(\delta_{ij}-\bar\delta_{i\cdot}-\bar\delta_{\cdot j}
  +\bar\delta)^2 .
\end{aligned}
\label{eq:crossed-decomp}
\end{equation}
Equation~\eqref{eq:crossed-decomp} partitions the observed variation into heading,
layout and joint interaction components.
If $\delta_{ij} = g(\theta_i)$ for some $g$, the second and third terms vanish;
if $\delta_{ij} = h(z_j)$, the first and third do. A non-zero third term
therefore identifies a joint heading--layout component beyond the two marginal
effects, making the interaction contribution explicit.
\end{proposition}

\begin{proof}
Expand $\delta_{ij}-\bar\delta$ as the sum of the two centred margins and the
remainder; the three are orthogonal under summation over the full rectangle, so
the cross terms vanish and the identity is arithmetic, with nothing estimated.
The two special cases make the corresponding margins constant.
\end{proof}

\begin{corollary}[The zero-tolerance repeat set]\label{cor:zerotol}
A crossed design contains, at every heading, $\binom{S}{2}$ pairs of episodes
per arm whose headings are equal. Definition~\ref{def:repeat} at
$\varepsilon = 0$ therefore admits a non-empty set, and $D_0(w)>0$ directly
identifies scene dependence at fixed arm and heading, with no tolerance or
smoothness assumption.
\end{corollary}

\begin{proof}
Take $\varepsilon=0$ in Definition~\ref{def:repeat}: two episodes at the same
$\theta_i$ satisfy $d(\theta,\theta')=0$, so they form a repeat. If completion
at $w$ were $f_w(a,\theta)$ then every such pair would agree, so a positive
$D_0(w)$ directly identifies the additional scene component without the arc
hypothesis used in Proposition~\ref{prop:repeat}.
\end{proof}

Corollary~\ref{cor:zerotol} provides a zero-tolerance complement to
Proposition~\ref{prop:tol} for the crossed record. The repeated headings make the
scene component directly observable, without relying on a tolerance interval.

Intervals on the crossed record are taken over layout clusters, not over cells. The
\NAssignedHeadings{} headings are a fixed grid that covers the circle, so they
are averaged within each layout rather than treated as sampled units. The
\NCrossedClusters{} layouts are the independent clusters; the interval at each
force level is an ordinary percentile bootstrap over their layout-level means.
With $H=\NAssignedHeadings{}$ and $S=\NCrossedClusters{}$, the estimand and one
bootstrap replicate are defined in Eq.~\eqref{eq:crossed-bootstrap}:
\begin{equation}
  \begin{aligned}
  \bar\delta_j(w)&=\frac{1}{H}\sum_{i=1}^{H}\delta_{ij}(w),&
  \bar\delta(w)&=\frac{1}{S}\sum_{j=1}^{S}\bar\delta_j(w),\\
  \bar\delta^{*(b)}(w)&=\frac{1}{S}\sum_{j=1}^{S}\bar\delta_{J_{bj}}(w),&
  J_{bj}&\stackrel{\mathrm{iid}}{\sim}\operatorname{Unif}\{1,\ldots,S\},
  \end{aligned}
  \label{eq:crossed-bootstrap}
\end{equation}
Equation~\eqref{eq:crossed-bootstrap} defines the layout-cluster mean and its
deterministic bootstrap replicate. Here $b=1,\ldots,B$ and $B=\NCrossedBootstrap{}$. The reported 95\% interval
is the empirical 2.5th--97.5th percentile range of the bootstrap means. Treating
all \NPairsPerForce{} cells as independent would understate uncertainty and is
not done. The crossed-design record does not add a signed-rank $p$ value to this
interval readout; its Holm family metadata is retained for the separate endpoint
audit.

\subsection{Reproducibility}

Every numerical result reported in this manuscript is programmatically derived
from the recorded simulation logs and statistical analysis code. Verification
manifests bind all analytical outputs and figures to the underlying data before
rendering, ensuring end-to-end consistency across tables and text. This
separation of source, computation and rendered claim follows reproducible-research
and simulation-reporting guidance~\cite{peng2011repro,sandve2013repro,
monks2019stress,ozturk2025mrtasim}.

Figures are generated in R with patchwork from the same bound tables: each panel
is constructed as an independent plot, then composed with the width or height
ratio recorded in the figure-layout contract. Panel tags are added as plot tags at
the upper-left and titles remain centred within their own panel; the exported PDF
is checked for the declared physical canvas before it enters the manuscript.

\subsection{Computational validation and transfer protocol}

The validation target is the declared simulator operating envelope, not an
implicit claim of completed field flight. Four checks are applied before a result
enters the manuscript. First, numerical reproducibility is enforced by the
SHA-256 evidence manifest, deterministic seeds and a clean-room rerun of the
legacy receipt. Second, internal system validity is established with paired
aware--agnostic episodes: the simulator, scene, horizon, controller, allocator
interface and endpoint are held fixed while one information or geometry factor
is opened at a time. Third, endpoint validity is checked by reading thresholded
completion together with closest approach and by repeating the force ladder over
capture radii. Fourth, transfer validity is represented by a declared parameter
correspondence rather than by an unreported hardware claim.

The correspondence is operational. A target vehicle is specified by mass,
available thrust, controller update rate, wind-estimator update rate and capture
criterion; the simulator quantities in the parameter contract are the initial
values for those fields. A hardware or hardware-in-the-loop extension should
re-estimate mass and thrust, time-align measured wind with vehicle state, replay
the same constant and temporal profiles at the declared update rate, and run
both arms on identical scenes or recorded traces. The primary acceptance
quantity remains the paired completion contrast at the transition, with the
continuous closest-approach endpoint and the clearance diagnostic reported
alongside it. This protocol makes the current computation a reproducible
pre-flight screening and mechanism-isolation study while leaving physical
transfer as a separately testable stage.

The distinction is part of the contribution. The present paper establishes what
can be measured, reproduced and stress-tested before hardware is committed; it
does not substitute simulator evidence for field evidence. All conclusions are
therefore indexed to the stated rigid-body simulator, force profiles, geometry
cells and homogeneous team sizes, and the transfer protocol identifies exactly
which quantities must be revalidated when those conditions change.


\section{Results}

\subsection{Operating envelope}

Table~\ref{tab:band} and Fig.~\ref{fig:band}A give the full sweep. The continuous
closest-approach companion is reported below with the tracking endpoint and is
read alongside the thresholded rate.
At \KneeLevel{}~N, which is \KneePctHover{} percent of one airframe's hover thrust, the
wind-agnostic allocator completes \KneeAgnostic{} of waypoints on average and the
wind-aware allocator completes \KneeAware{}. The paired difference is
$+\KneeDelta{}$ with an exact two-sided signed-rank $p$ of \KneeP{}, which clears
the Bonferroni-corrected threshold of \BonferroniAlpha{} for \NWindLevels{} levels.
The two marginal standard deviations overlap at that level, so the ladder on its
own understates the separation: every seed moves in the same direction even though
the arms' dispersions do not separate. Fig.~\ref{fig:band}B therefore plots the
seed-paired difference the test is actually computed on, and its 95~percent
interval excludes zero at \KneeLevel{}~N and at no other force in the sweep.
That one-level reading is scoped to the stock integral gain \StockKi{}~N/m/s,
the \EpisodeHorizonS{}~s horizon, the \CaptureRadius{}~m capture radius and the
declared 0.05~N force grid; the controller-identification audit shows how the
same endpoint moves when those settings change.
This force is specific to the \CaptureRadius{}~m capture radius that defines
completion here; rerunning the ladder at two further radii finds a transition
of comparable size at a different force in each, so \KneeLevel{}~N is the
transition of this operating point and moves with the capture radius.

The full sweep gives the envelope around that transition. At and below
\SaturationLevel{}~N both allocators complete nearly every waypoint, establishing a
saturated shoulder; at and above \CollapseLevel{}~N both arms occupy the measured
high-force edge. The four all-zero levels and the two levels with $p=1$ provide
explicit convergence boundaries for the endpoint, alongside the transition where
the arms separate. Taken together, the shoulder, transition and high-force edge
map where directional information has operational leverage; the complete ladder
makes each regime visible.

\begin{table}[tbp]
\centering
\caption{Completion rate by wind force for the wind-agnostic and wind-aware arms,
over \NSeeds{} paired seeds per level. Entries are mean $\pm$ standard deviation;
intervals and exact signed-rank $p$ values use seed-level paired differences. A dash
identifies a saturated all-zero level where no signed-rank statistic is defined.}
\label{tab:band}
\generatedtable{generated_table_band}
\end{table}

\begin{figure}[tbp]
\centering
\includegraphics[width=\linewidth]{fig1_wind_band.pdf}
\caption{\CapBand}
\label{fig:band}
\end{figure}

\subsection{Estimate-noise robustness}

Table~\ref{tab:noise}, Fig.~\ref{fig:noise}A and Fig.~\ref{fig:noise}B report the
controlled estimate perturbation that motivates this robustness analysis.
Panel A tracks the transition force, while Panel B verifies its neighbouring
force levels. Introducing controlled perturbations to the wind vector supplied to
the wind-aware allocator preserves the band. The paired advantage is
$+\SigmaZeroDelta{}$ with an exact input, and remains statistically aligned with
that value across relative noise of \NearestSigma{} and \BestSigma{}, where
the point estimates are $+\NearestSigmaDelta{}$ and $+\BestSigmaDelta{}$.

At the largest tested relative noise, \MaxSigma{}, where the corruption is as large
as the signal and the heading error has a standard deviation of roughly a radian,
the measured advantage reaches its robustness edge at $+\MaxSigmaDelta{}$ with an
exact $p$ of \MaxSigmaP{}. Every 95~percent interval in the table excludes zero, the
weakest lower bound being \WeakestLower{} at that edge.

The decision point fixed before the runs was relative noise of \PreregSigma{};
the nearest executed grid level is \NearestSigma{}, so we evaluate the pre-set
decision on that grid-aligned condition, where the interval is bounded away from
zero. Counting the complete disturbance-by-noise family gives a strict
Bonferroni threshold of \NoiseBonferroni{}. The 0.30 and 0.50 relative-noise
conditions clear that family threshold; the 0.10 condition is retained as a
positive descriptive estimate below that threshold, and the 1.00 condition
($p=\MaxSigmaP{}$) marks the measured edge without a multiplicity-adjusted
significance claim. Read together, the curve is a profile of how directional information translates
into allocation benefit as estimate quality changes.
The neighbouring disturbance levels stay flat at every noise level, which answers a
question the point estimates alone would leave open: degrading the estimate leaves
the transition location stable and changes the advantage within the measured band.
That stability is a useful calibration result for selecting the next estimator
quality regime.

\begin{table}[tbp]
\centering
\caption{Effect of relative wind-estimate noise on the paired advantage at
\KneeLevel{}~N over \NSeeds{} paired seeds. The reference arm is deterministic given
a seed, so its recorded outcome is reused exactly for each noise level. Intervals are 95~percent $t$ intervals on seed-level differences;
$p$ values are exact two-sided signed-rank tests.}
\label{tab:noise}
\generatedtable{generated_table_noise}
\end{table}

\begin{figure}[tbp]
\centering
\includegraphics[width=\linewidth]{fig2_noise_sensitivity.pdf}
\caption{\CapNoise}
\label{fig:noise}
\end{figure}

\subsection{Episode-level distribution}

Because an episode assigns \NDrones{} quadrotors to \NDrones{} waypoints, its
completion rate takes one of five values, and a mean over \NSeeds{} seeds can hide
very different distributions. Fig.~\ref{fig:spread} opens the two means at the
discriminating level. The wind-aware mean of \KneeAware{} is concentrated rather
than diffuse: \AwareTopSeeds{} of \NSeeds{} seeds finish three of the four
waypoints, while \AwareZeroSeeds{} seed finishes none. The wind-agnostic mean of
\KneeAgnostic{} is similarly concentrated, with \AgnosticZeroSeeds{} of
\NSeeds{} seeds finishing nothing.

This distribution is part of the contribution. The gain is a near-categorical
change in outcome for most seeds, with a smaller set of matched scenes producing
the same endpoint under both arms. Reporting the mean together with this
distribution makes the operating envelope useful for aggregate planning and
episode-level expectations.

\begin{figure}[tbp]
\centering
\includegraphics[width=\linewidth]{fig3_seed_spread.pdf}
\caption{\CapSpread}
\label{fig:spread}
\end{figure}

\subsection{Continuous tracking endpoint}

The completion rate is a step function of a continuous quantity: an aircraft counts
as finished when its closest approach falls inside the \CaptureRadius{}~m capture
radius. Fig.~\ref{fig:miss} plots that continuous quantity directly and shows
how the operational gain is formed. At the discriminating level
the mean closest approach is \MissAgnosticKnee{}~m for the wind-agnostic allocator
and \MissAwareKnee{}~m for the wind-aware one, a difference of under two
centimetres, while the completion rates at the same level differ by
$\KneeDelta{}$. At the largest disturbance the two arms are closer still,
\MissAgnosticCollapse{}~m against \MissAwareCollapse{}~m, and neither completes
anything.

The wind-aware allocator flies slightly more accurately at the disturbance where
that small improvement moves most aircraft across the capture boundary. This is the
operational leverage of the method: target-position feedforward turns a modest
continuous tracking shift into a substantial task-level change at the transition.
The continuous endpoint makes that mechanism explicit and links the completion
jump to its physical displacement. This yields a general measurement principle
that extends beyond wind: when an endpoint thresholds a continuous quantity, the
reported effect is determined by where the boundary sits relative to the two
distributions as well as by the intervention. Pairing the rate with its continuous
companion therefore shows both the operational payoff and the physical scale of
the tracking change.

The object-level bridge in Fig.~\ref{fig:system-overview}B makes that relationship
visible at the seed level. At the transition, \BridgeLowerPairs{} of the
\NSeeds{} matched seeds move to a lower mean closest approach under the aware
arm, with a mean shift of $\BridgeDistanceShift{}$~m and a completion gain of
$+\BridgeCompletionGain{}$ rate units. This bridge is a descriptive reuse of
the primary record; the inferential unit and the single paired test are
unchanged.

\begin{figure}[tbp]
\centering
\includegraphics[width=\linewidth]{fig4_miss_distance.pdf}
\caption{\CapMiss}
\label{fig:miss}
\end{figure}

\subsection{Scene-context variation}

The episode-level distribution leaves an important systems question. If the
wind-aware mean is a mixture over seeds, which part of the episode context
carries that spread? A seed fixes both the force heading and the waypoint layout, so near
repetitions provide a direct probe of how the operating envelope depends on
context. The record supports a joint reading: the force level locates the
transition, while the scene determines how strongly an individual episode
expresses it.
Table~\ref{tab:seedpairs} prints the discriminating cell one seed at a time.
The heading and replanning columns expose context that is absent from aggregate
summaries.
The heading is drawn once per seed and reused at every force level, in both arms
and in the second sweep; near-matched seeds therefore repeat one disturbance while
retaining different layouts. Fig.~\ref{fig:headings}A shows the \NSeeds{} draws
and Fig.~\ref{fig:headings}B their complete separation spectrum. \NRepeats{} pairs
lie within \WidestRepeat{} degrees (the closest is \ClosestRepeat{}), resolving to
\NHeadings{} distinct headings with \LargestHeadingGroup{} inside one
\HeadingGroupSpan{}-degree arc. The next pair is \NearestDistinct{} degrees away,
so Proposition~\ref{prop:tol} gives the same repeat set throughout the displayed
tolerance interval.

Table~\ref{tab:repeats} and Fig.~\ref{fig:repeats} report how those
repetitions spread. At the discriminating level, \RepeatDisagreeKnee{}
of the \NRepeats{} repeated disturbances disagree, and the largest disagreement
is \RepeatGapKneeDrones{} of the \NDrones{} airframes. The two seeds whose
headings differ by \ClosestRepeat{} degrees return different outcomes, and one
pair spans the entire range the wind-aware arm ever reached. By Proposition~\ref{prop:repeat} the completion
rate at \KneeLevel{}~N is jointly shaped by the disturbance and the scene
context, and the gap of $\RepeatGapKnee{}$ quantifies the context-dependent
variation that an operating-envelope design can accommodate.
The comparison level matters as much as the number. \NDegenerateLevels{} of the
\NWindLevels{} levels are constant across seeds; Proposition~\ref{prop:vacuous}
therefore identifies their repeat gaps as saturated ceiling or floor calibration
values. The remaining \NInformativeLevels{} levels carry the scene-variation signal:
at \CollapseLevel{}~N the seed spread is \CollapseSeedSpread{} and
\RepeatDisagreeCollapse{} repeat pairs disagree, at \SaturationLevel{}~N
\RepeatDisagreeShoulder{} disagree by \RepeatGapShoulderDrones{} airframe, and at
\KneeLevel{}~N \RepeatDisagreeKnee{} disagree by up to
\RepeatGapKneeDrones{}. Context dependence is consequently concentrated at the
transition where the arms can be distinguished.
Replanning is logged when an allocator changes its assignment mid-episode. At the
discriminating level the wind-agnostic and wind-aware arms replanned
\ReplanKneeAgnostic{} and \ReplanKneeAware{} times, respectively; the
\NReplanLoud{} seeds with at least one event carry $+\ReplanLoudDelta{}$ versus
$+\ReplanQuietDelta{}$ for the \NReplanQuiet{} quiet seeds. With \NSeeds{} seeds
this remains a descriptive location of variance, reported beside the primary test.

Corollary~\ref{cor:notn} clarifies the inferential unit and its operational
reading. The test's units are the \NSeeds{} per-seed differences;
two seeds sharing a heading still differ in their waypoint layouts, and the
advantage of $+\KneeDelta{}$ at an exact $p$ of \KneeP{} stands as reported.
The repeat profile shows that the same force heading can differ by
\RepeatGapKneeDrones{} of \NDrones{} airframes across layouts. The paired mean
locates the opportunity; the episode-level spread supplies a concrete margin for
the next scene-aware deployment study.

\begin{table}[tbp]
\centering
\caption{Seed-level paired record at \KneeLevel{}~N. Each row uses the same force
and waypoint layout in both arms. Heading is drawn once per seed and reused across
force levels; completion is the number of finished airframes, and replanning is the
number of mid-episode assignment changes.}
\label{tab:seedpairs}
\generatedtable{generated_table_seedpairs}
\end{table}

\begin{table}[tbp]
\centering
\caption{Near-matched disturbance repeats. Each row pairs seeds whose headings
differ by less than \RepeatTol{} degrees; entries are the larger within-arm
completion difference in airframes. The final row records whether any seeds differ
at that force, distinguishing informative gaps from saturated levels.}
\label{tab:repeats}
\small
\generatedtable{generated_table_repeats}
\end{table}

\begin{figure}[tbp]
\centering
\includegraphics[width=\linewidth]{fig5_heading_coverage.pdf}
\caption{\CapHeadings}
\label{fig:headings}
\end{figure}

\begin{figure}[tbp]
\centering
\includegraphics[width=\linewidth]{fig6_replicate_gap.pdf}
\caption{\CapRepeats}
\label{fig:repeats}
\end{figure}

\subsection{Heading--layout sweep}

The original sampling record supplies near-matched repetitions. We extend that
record with a designed crossed sweep that resolves the remaining sampling
question: headings are fixed on a grid and every heading is paired with every
waypoint layout. The simulator is deterministic, single-core and about a second
per episode, so this direct coverage test is inexpensive and fully reproducible.
Before comparing the two records we re-ran the first. All \NReproEpisodes{}
episodes of the drawn-heading sweep return the same value in all
\NReproFields{} recorded fields, confirming identical reproducibility of the
baseline results. The crossed record below therefore extends the same validated
endpoint.
The crossed sweep uses a fixed heading grid:
\NAssignedHeadings{} headings \HeadingStep{} degrees apart, each flown with all
\NScenes{} waypoint layouts, at every force level and in both arms
(Definition~\ref{def:crossed}). It leaves \DesignedArc{} degrees unsampled
against \UnsampledArc{}, and it puts \NPairsPerForce{} paired comparisons at
each level against \NSeeds{}. Table~\ref{tab:designed} and
Fig.~\ref{fig:designed}A report the force-level envelope, while
Fig.~\ref{fig:designed}B opens its transition cell.

The operating envelope is reproduced, and the transition remains one level wide.
At \KneeLevel{}~N the paired
advantage is $+\DesignedKneeDelta{}$ with a 95~percent ordinary-bootstrap
interval over the \NCrossedClusters{} layout clusters,
$[+\DesignedKneeLower{}, +\DesignedKneeUpper{}]$. Each layout cluster first
averages the \NAssignedHeadings{} fixed headings; the \NPairsPerForce{} cells
support the design and visualization, while the layout clusters provide the
inferential units. The estimate is from
\DesignedKneeAgnostic{} completed under the wind-agnostic allocator to
\DesignedKneeAware{} under the wind-aware one; the ten-draw estimate of
$+\KneeDelta{}$ sits inside that interval. Exactly \DesignedKneeAdverse{} of the
\NPairsPerForce{} pairs at that level favours the wind-agnostic arm. The
neighbouring levels stay where they were, $+\DesignedShoulderDelta{}$ at
\SaturationLevel{}~N and $+\DesignedCollapseDelta{}$ at \CollapseLevel{}~N, so
the designed sweep confirms the narrowness of the band.
The context dependence is visible at exact heading (zero tolerance). At
\KneeLevel{}~N two episodes at the \emph{same} heading, differing only in
waypoint layout, span $\ZeroSeparationGap{}$ of paired advantage, from
$\ZeroSeparationLow{}$ to $\ZeroSeparationHigh{}$. By
Corollary~\ref{cor:zerotol}, exact-heading repeats establish a scene component
that is separate from the arm and the wind, without a smoothness assumption.
The crossed design therefore identifies the residual scene context directly:
the same directional input expresses different completion outcomes across
layouts, which makes layout a measurable operating axis.
This is a complementary readout of the primary endpoint; Corollary~\ref{cor:zerotol}
keeps the repeated-heading statement tied to the declared design without altering
the force-family test.

Panel B of Fig.~\ref{fig:designed} opens the discriminating level into
its heading-by-layout rectangle, one cell per paired episode. Flat rows would
indicate a heading-only pattern and flat columns a layout-only pattern; the
observed matrix resolves both margins and their interaction. Of the variation in
the advantage at that level, \DesignedShareHeading{} per cent lies in the heading
margin, \DesignedShareScene{} per cent in the layout margin, and
\DesignedShareResidual{} per cent in the interaction. By
Proposition~\ref{prop:crossed}, the interaction is the dominant component of the
measured variation: wind and layout jointly shape the transition, while each
margin contributes a smaller, interpretable share.
Read with the paired test, the two statements are complementary. Over
\NPairsPerForce{} pairs the wind-aware allocator finishes
\DesignedKneeDelta{} more of \NDrones{} aircraft on average at \KneeLevel{}~N,
with a tighter estimate than the original ten draws. Exact wind knowledge still
leaves a scene-dependent spread on an individual flight; the mean locates the
opportunity and the spread supplies the deployment margin.

\begin{table}[tbp]
\centering
\caption{Crossed heading--layout sweep with \NAssignedHeadings{} fixed headings
spaced \HeadingStep{} degrees apart and \NScenes{} layouts per heading. The
reported interval is an ordinary bootstrap over \NCrossedClusters{} layout-level
means after averaging the fixed headings (\NCrossedBootstrap{} draws); the cells are
summarized through their layout-cluster units. The repeat-gap column is the largest
same-heading difference; the final columns give the heading, layout and interaction
shares computed from the displayed crossed cells.}
\label{tab:designed}
\small
\generatedtable{generated_table_designed}
\end{table}

\begin{figure}[tbp]
\centering
\includegraphics[width=\linewidth]{fig7_designed_replicates.pdf}
\caption{\CapDesigned}
\label{fig:designed}
\end{figure}

\subsection{Temporal operating profiles}

The constant-force result establishes where directional information has leverage;
the next question is whether that leverage survives when the vector itself moves.
We therefore ran a primary campaign across \NTemporalProfiles{} profiles,
\NTemporalEstimators{} estimate conditions and three force levels, with the same
paired scene and phase for every arm. The campaign contains \NTemporalEpisodes{}
episode records. Fig.~\ref{fig:temporal}A and Fig.~\ref{fig:temporal}B, together
with Table~\ref{tab:temporal}, report the causal-estimate slice at the transition
alongside the oracle and biased-causal conditions.

At \KneeLevel{}~N the causal estimate gives a paired gain of
$+\TemporalConstantDelta{}$ for the constant profile, with a 95~percent interval
$[+\TemporalConstantLower{},+\TemporalConstantUpper{}]$. The slow gust gives
$+\TemporalSlowDelta{}$ with interval
$[+\TemporalSlowLower{},+\TemporalSlowUpper{}]$, and the fast gust gives
$+\TemporalFastDelta{}$ with interval
$[+\TemporalFastLower{},+\TemporalFastUpper{}]$. The exact two-sided signed-rank
$p$ values for these three rows are printed in Table~\ref{tab:temporal}.
The transition is therefore carried into both temporal regimes: its expression
changes with the rate of the gust, while its location remains a force-scale
phenomenon expressed across dynamic profiles.

The observer diagnostics make that statement operational. Mean relative vector
tracking error is \TemporalTrackingConstant{}~percent for the constant profile,
\TemporalTrackingSlow{}~percent for the slow gust and \TemporalTrackingFast{}~percent
for the fast gust. The neighbouring force cells move by no more than
\TemporalShoulderMax{} in paired gain for the causal condition, so the temporal
extension preserves the narrow shoulder--transition--edge structure of the
original envelope. The \NTemporalLegacyRows{} constant/oracle regression rows
reproduce the legacy control path, tying the new dynamic readout to the established
reference while preserving the original endpoint.

\begin{table}[tbp]
\centering
\caption{Temporal operating-envelope results at the transition force. Each row is a
profile--estimator condition; intervals use the \NTemporalSeeds{} paired seeds,
tracking is mean relative vector error, and $p$ is an exact signed-rank test against
the reference arm. The oracle uses the exact vector, the causal condition uses a
first-order update, and the biased-causal condition adds the stated bias and noise.}
\label{tab:temporal}
\small
\generatedtable{generated_table_temporal}
\end{table}

\begin{figure}[tbp]
\centering
\includegraphics[width=\linewidth]{fig8_temporal_envelope.pdf}
\caption{\CapTemporal}
\label{fig:temporal}
\end{figure}

\subsection{Pathway and gain}

The complete wind-aware arm combines an upwind assignment cost with a
target-position feedforward term. A feedforward-off support campaign retains the
directional assignment cost and selectively disables the target-position bias, producing
\NAblationEpisodes{} paired records under the same temporal profiles, force levels,
seeds and estimate states. At the transition, the allocation-only gains are
$+\AllocOnlyConstant{}$, $+\AllocOnlySlow{}$ and $+\AllocOnlyFast{}$ for the
constant, slow-gust and fast-gust profiles; the corresponding complete-arm gains
are $+\CompleteConstant{}$, $+\CompleteSlow{}$ and $+\CompleteFast{}$. In this
benchmark the comparison localizes the observed task-level change to the
target-position pathway while preserving the directional assignment component,
thereby identifying a concrete tuning lever for the interface; the assignment
cost remains an explicit component to test across the next scene geometries.

The controller setting is then opened as an explicit factor. A
\NGainEpisodes{}-episode sensitivity campaign varies the feedforward gain from zero
to 1.40~m/N on the fixed seven-point grid, with every condition paired to the same
scene and causal estimate. Around the inherited \FeedforwardGainDefault{}~m/N
setting, the tested profiles occupy a broad positive plateau from
\GainPlateauMin{} to \GainPlateauMax{} in paired completion gain. At the upper
tested edge the gains are $+\GainEdgeConstant{}$, $+\GainEdgeSlow{}$ and
$+\GainEdgeFast{}$ for the three profiles. Fig.~\ref{fig:mechanism}A and
Fig.~\ref{fig:mechanism}B, together with Table~\ref{tab:mechanism}, make the
design choice visible: the inherited setting is inside a measured operating
region, while the upper tested edge marks the boundary of the gain grid.

\begin{table}[tbp]
\centering
\caption{Pathway and feedforward-gain sensitivity at the transition force. The
allocation-only arm retains the upwind assignment cost but disables target
feedforward; the complete arm uses the inherited gain. Plateau and upper-edge
columns summarize the tested gain grid. Entries are paired completion gains.}
\label{tab:mechanism}
\small
\generatedtable{generated_table_mechanism}
\end{table}

\begin{figure}[tbp]
\centering
\includegraphics[width=\linewidth]{fig9_mechanism_gain.pdf}
\caption{\CapMechanism}
\label{fig:mechanism}
\end{figure}

\subsection{Scene-geometry factors}

The crossed heading--layout result identifies scene context as a major source of
episode variation. We therefore opened that context with a targeted factorised
campaign that resolves ``layout'' into named factors. The
campaign contains \NGeometryEpisodes{} episodes across eight combinations of
target spread, target range and crossing orientation, three force levels and
\NGeometrySeeds{} paired seeds. The nominal--near--aligned cell is the exact legacy
scene, so the extension has a direct compatibility anchor.
Table~\ref{tab:geometry} lists the nine force-by-factor contrasts and their
declared paired units. Fig.~\ref{fig:geometry}A gives
the high-minus-low factor contrasts, each formed by first averaging the four
crossed settings of the other two factors within each of the \NGeometryClusters{}
seed clusters; the inferential vector therefore has \NGeometryClusters{} terms.
The \NGeometryTermSensitivity{} seed-by-setting terms serve as a sensitivity
readout, while the primary support-unit count remains the seed-cluster vector.
Fig.~\ref{fig:geometry}B
shows all condition-by-force means. At the transition force, crossing the target
geometry gives the largest contrast, $+\GeometryKneeCrossing{}$
$[+\GeometryKneeCrossingLower{},+\GeometryKneeCrossingUpper{}]$ with an exact
two-sided signed-rank $p=\GeometryKneeCrossingP{}$ (Holm-adjusted
$p=\GeometryKneeCrossingHolm{}$ within the nine-contrast support family).
Expanding target spread shifts the gain by
$\GeometryKneeSpread{}$ $[\GeometryKneeSpreadLower{},\GeometryKneeSpreadUpper{}]$,
while moving the target range farther shifts it by
$\GeometryKneeRange{}$ $[\GeometryKneeRangeLower{},+\GeometryKneeRangeUpper{}]$.
The crossing contrast is therefore a concrete scene-design lever in this
benchmark, whereas the spread and range contrasts quantify the direction and
precision of two neighbouring changes. The range estimate is a precision
boundary rather than an equivalence claim: its 30-cluster interval reaches the
zero boundary at the displayed resolution, while the 120-term sensitivity
interval is $[\GeometryRangeSensitivityLower{},+\GeometryRangeSensitivityUpper{}]$.
Using the declared nine-contrast family, the planning calculation gives an 80\%
familywise MDE of $\GeometryRangeMDE{}$ completion-rate units for this range
contrast.

These nine factor-by-force contrasts form one declared geometry-support family.
The high-minus-low terms are paired within seed and crossed scene condition, then
averaged within seed before the exact signed-rank test and interval are computed.
The \NGeometryTermSensitivity{} terms at a force level provide a design-level
sensitivity analysis, while the \NGeometryEpisodes{} episode rows document the
full crossed record and its audit trail.

The condition panel keeps the interpretation anchored to the endpoint: the
nominal legacy cell gives a transition gain of +\GeometryKneeNominal{}, and the
strongest tested cell reaches +\GeometryKneeBest{}. The factor effects define a
controlled operating-envelope map within the declared simulator and provide a
repeatable bridge from ``scene matters'' to an experimental program in which target
geometry can be selected, measured and extended along named axes.

\begin{table}[tbp]
\centering
\caption{Factorised scene-geometry campaign. Each row is a high-minus-low contrast
for one factor and force, with the other two factors crossed. Intervals and exact
signed-rank tests use \NGeometryClusters{} seed-level contrasts after averaging
the four crossed settings within seed; the \NGeometryTermSensitivity{}-term
version is a sensitivity readout. Holm adjustment is across the nine completion
contrasts in the support family. Levels are deterministic simulator
transformations; nominal--near--aligned is the legacy compatibility cell.}
\label{tab:geometry}
\small
\generatedtable{generated_table_geometry}
\end{table}

\begin{figure}[tbp]
\centering
\includegraphics[width=\linewidth]{fig10_geometry_factorial.pdf}
\caption{\CapGeometry}
\label{fig:geometry}
\end{figure}

\subsection{Team-size transfer}

The final extension asks whether the transition is retained when the homogeneous
team grows. The scale campaign contributes \NScaleEpisodes{} episode rows across
\NScaleSizes{} team sizes, \NScaleForces{} force levels and \NScaleSeeds{} paired
seeds, with the causal directional estimate and the same endpoint used above.
Fig.~\ref{fig:scale}A shows the complete size--force surface: the transition
remains concentrated at \KneeLevel{}~N while the neighbouring cells stay close to
the envelope baseline. This isolates transfer of the main interface across team
size with the controller and simulator held fixed.

At the transition, the paired gains are $+\ScaleKneeTwo{}$,
$+\ScaleKneeFour{}$, $+\ScaleKneeSix{}$ and $+\ScaleKneeEight{}$ for the
declared sizes from \ScaleMinSize{} to \ScaleMaxSize{} airframes. Every
95~percent transition-cell interval in Table~\ref{tab:scale} remains positive;
the four-airframe slice is $+\ScaleKneeFour{}$
$[+\ScaleKneeFourLower{},+\ScaleKneeFourUpper{}]$, and all exact tests are at
or below the Holm-adjusted bound $\ScaleKneeHolmMax{}$. The transition gain therefore spans a compact
$+\ScaleKneeMin{}$--$+\ScaleKneeMax{}$ range across the declared sizes.

Panel B of Fig.~\ref{fig:scale} separates the endpoint means at the transition.
The wind-aware arm remains above the reference at every tested size. The
four-airframe constant-profile causal slice (30 paired seeds, both arms and three
force levels; 180 rows) is field-by-field identical to the temporal record and is
distinct from the ten-seed oracle values in Table~\ref{tab:band}. The full table
retains the shoulder and high-force cells that delimit this homogeneous-team
transfer claim.

\begin{table}[tbp]
\centering
\caption{Homogeneous swarm-size operating-envelope campaign. Each row is a size--force
cell over \NScaleSeeds{} paired seeds. Entries are mean completion $\pm$ standard
deviation; gain intervals are 95~percent $t$ intervals. The final column gives the
raw exact signed-rank $p$ followed by the Holm-adjusted $p$ across the 12 size--force
cells. Four-airframe rows are the compatibility slice; other sizes are declared
extensions.}
\label{tab:scale}
\small
\generatedtable{generated_table_scale}
\end{table}

\begin{figure}[tbp]
\centering
\includegraphics[width=\linewidth]{fig11_swarm_scale.pdf}
\caption{\CapScale}
\label{fig:scale}
\end{figure}

\subsection{Endpoint-radius sensitivity}

Completion is scored by a capture radius, so every rate in this paper depends on a
distance the experiment chose. The validation sweep reruns the
force ladder at radii of \RadiusGrid{}~m over \NRadiusEpisodes{} episodes, holding
scene, heading, seed pairing and horizon fixed, and the anchor column reproduces
the primary ladder exactly. Table~\ref{tab:radius} gives the full grid and
Fig.~\ref{fig:radius}A the paired advantage at each radius. The record also
carries the underlying continuous distance as its own corrected family on the
same rows, and it moves the same way: the wind-aware arm ends closer to its
target in \RadiusContinuousCloser{} of the \RadiusContinuousCells{} radius-force
cells, with \RadiusContinuousHolmSurvivors{} clearing Holm. The threshold
therefore reports a discretization of a distance difference that the continuous
endpoint records independently of where the threshold is placed.

The result separates two claims that the primary sweep states together. A
transition exists at every tested radius and its size is comparable across them,
from $+\RadiusPeakDeltaMin{}$ to $+\RadiusPeakDeltaMax{}$ completion. Its
location moves with the radius: the largest advantage sits at
\RadiusTightPeakForce{}~N for the \RadiusTightLevel{}~m radius,
\RadiusAnchorPeakForce{}~N for the \CaptureRadius{}~m anchor and
\RadiusWidePeakForce{}~N for the \RadiusWideLevel{}~m radius, a span of
\RadiusPeakForceSpan{}~N in force across a \RadiusSpanM{}~m change in radius.
Fig.~\ref{fig:radius}B shows the mechanism directly, since a tighter radius
makes the same disturbance harder and shifts the whole completion ladder toward
lower forces. The headline force is therefore a property of the operating
point: directional information has its largest leverage at the force where the
task is on the edge of feasibility, and that force moves with the capture
radius.

The power of the family sets the scope of that reading. Within this predeclared
\RadiusHolmFamily{}-test validation family \RadiusHolmSurvivors{} of the cells
clears a Holm-corrected 5~percent threshold, and the anchor cell reaches
\RadiusAnchorHolmP{}; with \NRadiusSeeds{} seeds per cell the family locates
the transition at each radius descriptively and is reported as a sensitivity
readout rather than a second confirmation. The primary seven-force result
retains its own Bonferroni family and is unchanged.

\begin{table}[tbp]
\centering
\caption{Completion by capture radius and wind force for both allocators, over
\NRadiusSeeds{} paired seeds per cell. ``Paired diff.'' is the wind-aware minus
wind-agnostic seed-level mean with its 95~percent interval, and ``Holm $p$'' is
corrected within the \RadiusHolmFamily{}-test endpoint-radius validation family.
Cells with an undefined exact test are all-zero differences.}
\label{tab:radius}
\scriptsize
\generatedtable{generated_table_radius}
\end{table}

\begin{figure}[tbp]
\centering
\includegraphics[width=\linewidth]{fig12_endpoint_radius.pdf}
\caption{\CapRadius}
\label{fig:radius}
\end{figure}

\subsection{Complete pathway factorial}

The ablation reported above removes the target-position feedforward and observes
the gain collapse. That is consistent with the feedforward carrying the effect and
equally consistent with the two pathways working only in combination, because the
cell with assignment disabled and feedforward retained had not yet been run.
Running it completes the $2\times2$ over \NFactorialEpisodes{} episodes and
\NFactorialSeeds{} paired seeds per profile-force cell.

Fig.~\ref{fig:pathway}A and Table~\ref{tab:factorial} give all three terms. At the
transition force the target-feedforward main effect spans
$+\FactorialFeedforwardMin{}$ to $+\FactorialFeedforwardMax{}$ completion across
the three temporal profiles, with a Holm-corrected $p$ of at most
\FactorialFeedforwardHolmMax{}. The assignment main effect and the
assignment-by-feedforward interaction stay within \FactorialOtherAbsMax{}
completion of zero at every profile and force, and \NFactorialHolmSignificant{} of
the \NumFactorialCells{} tested effects clear the corrected threshold, all of them
feedforward terms. In this benchmark the directional estimate therefore acts
through the target-position feedforward, and the upwind assignment cost and its
interaction with the feedforward stay inside that near-zero band at every
profile and force. The continuous endpoint carries its own corrected family and agrees: at
the transition force the feedforward term moves the mean minimum distance to
target by \FactorialContFeedforwardMin{} to \FactorialContFeedforwardMax{}~m,
closer in every profile, while the assignment and interaction terms stay within
\FactorialContOtherAbsMax{}~m of zero and none of them clears its threshold on
that endpoint either. This is a mechanism statement scoped to this interface
and task.

\subsection{Safety and feasibility diagnostics}

The feasibility campaign completes the execution envelope with descriptive
state and clearance diagnostics. It records \NSafetyMetrics{} diagnostics over
\NSafetyEpisodes{} episodes at \SafetyForceGrid{}~N, and every episode in both
arms remains finite. Table~\ref{tab:safety} gives the full diagnostic record and
Fig.~\ref{fig:pathway}B shows the observed clearance range. The closest
airframe pair is \SafetyWorstSeparation{}~m, in the
\SafetyWorstSeparationArm{} arm at \SafetyWorstSeparationForce{}~N; the
transition-force reference values are \SafetyAgnosticSeparationKnee{}~m and
\SafetyAwareSeparationKnee{}~m, and the largest motor-limit fraction is
\SafetyMotorMax{}. \SafetyGeomHullContacts{} of the \NSafetyEpisodes{} episodes
reach the declared hull-contact distance \SafetyGeomHullContactM{}~m. These
measurements define the benchmark's feasibility boundary and identify the
additional separation constraint required by a physical implementation.

That gap is reported as an observation about these episodes rather than as a
property of the allocator, because most of the metric's variance belongs to the
benchmark instead of to either arm. The scene's minimum target separation
correlates with the observed minimum airframe separation at
\SafetyGeomTargetCorr{}, and \SafetyGeomCloseTargetContacts{} of the
\SafetyGeomHullContacts{} sub-contact episodes fall in the half of scenes whose
targets sit closer than twice the capture radius, against
\SafetyGeomFarTargetContacts{} in the other half. Within a scene, where both arms
fly identical targets, the arm that finished more airframes flew closer in
\SafetyGeomPairedCloser{} pairs and farther in \SafetyGeomPairedFarther{}, with
\SafetyGeomPairedTied{} ties: an airframe that reaches its target sits at target
separation, while one carried off by the wind ends far from everything. The metric
tracks arrival, so on this measure the arm that completes more of the task scores
worse. The scale of the numbers is structural for the same reason. Targets are
drawn with no minimum-separation constraint while completion is scored inside the
capture radius, so two airframes can both be scored finished with their hulls
overlapping, in \SafetyGeomOverlapPrimary{}\% of scenes at the primary radius and
between \SafetyGeomOverlapMin{}\% and \SafetyGeomOverlapMax{}\% across the radii
of the endpoint sweep; that sweep therefore varies how physically admissible
success is as well as how strict it is.

The transition result does not rest on these episodes: completion and observed
separation correlate at \SafetyGeomCompletionCorr{} across the campaign. What the
clearance row supports is narrower than a comparison between allocators, and it is
a caution about the task rather than about the feedforward. Contact-scale
separations (reaching an observed minimum of \SafetyWorstSeparation{}~m) were
reached on a four-airframe simulated task whose objective carries no collision
constraint and whose physics branch models no inter-airframe downwash, at mean
separations of roughly two hull diameters, which is inside the regime where rotor
wake would dominate a physical flight. The observed minimum separation is
therefore retained as a baseline failure-mode diagnostic of the unconstrained
allocation--control interface: directional wind awareness alone cannot ensure
safe separation in crossing geometries. Only the scalar three-dimensional
separation is logged, so a pair separated mostly in altitude could clear the hull
at these distances and the counts above are an upper bound on contacts. A
physical multi-vehicle deployment strictly requires composing this feedforward
layer with an explicit outer-loop safety barrier—such as Control Barrier
Functions (CBFs), velocity obstacles, or safety ellipsoids—to guarantee clearance
and prevent downwash entanglement. What was checked is the list in
Table~\ref{tab:safety}; what was not checked includes collision avoidance under an
active safety filter, degraded or failed airframes, communication loss, energy
budgets and physical multi-rotor flight.

\begin{table}[tbp]
\centering
\caption{Complete assignment-by-feedforward factorial on paired completion, over
\NFactorialSeeds{} paired seeds per profile-force cell. Each term is a seed-level
paired contrast within the cell and ``Holm $p$'' is corrected across the
\NumFactorialCells{} tested effects. Cells with an undefined exact test are
all-zero differences.}
\label{tab:factorial}
\scriptsize
\generatedtable{generated_table_factorial}
\end{table}

\begin{table}[tbp]
\centering
\caption{Descriptive safety and feasibility diagnostics over \NSafetyEpisodes{}
episodes at a constant profile, \NSafetySeeds{} seeds per allocator and force.
Each cell gives the mean over seeds with the worst observed episode in
parentheses; ``worst'' is the minimum for separation and the maximum for the
remaining rows. No test is reported and none is implied.}
\label{tab:safety}
\small
\generatedtable{generated_table_safety}
\end{table}

\begin{figure}[tbp]
\centering
\includegraphics[width=\linewidth]{fig13_pathway_safety.pdf}
\caption{\CapPathwaySafety}
\label{fig:pathway}
\end{figure}

\subsection{Evidence}

The campaign results are read on their declared inferential units and presented as
a layered synthesis that keeps each view's unit visible. Table~\ref{tab:evidence} makes
the hierarchy explicit: \NEvidenceBlocks{} synthesis blocks span
\NEvidenceCampaigns{} campaign views and \NEvidenceRows{} episode rows, while
their units remain the ten primary seeds, twelve crossed layout clusters (each
after heading aggregation) or thirty paired seeds specified by each design. This
descriptive architecture keeps every estimate tied to the design that produced it.

Fig.~\ref{fig:evidence}A places the transition estimates on one rate scale while
retaining each campaign's interval. The primary oracle estimate is
$+\KneeDelta{}$; the crossed layout estimate is $+\DesignedKneeDelta{}$;
the causal temporal profiles span $+\TemporalSlowDelta{}$ to
$+\TemporalConstantDelta{}$; and the homogeneous size axis spans
$+\ScaleKneeMin{}$ to $+\ScaleKneeMax{}$. The estimate-noise rows show the
same endpoint from the exact vector through the tested high-noise edge. The
agreement is therefore visible as a set of condition-specific estimates, with
design differences retained rather than compressed into a pooled number. Fig.~\ref{fig:evidence}B reports episode volume on a logarithmic axis, with each
point labelled by the seed or layout-cluster unit supporting its interval. Together the
panels show the breadth of the campaign while keeping episode volume separate from
inferential replication. The three reviewer-validation blocks appear in the volume
panel and in Table~\ref{tab:evidence} but not in the estimate panel, because none
of them estimates the same quantity: the radius sweep varies the endpoint that
defines the estimate, the factorial contrasts pathways rather than allocators, and
the safety block reports no estimate at all.

\begin{table}[tbp]
\centering
\caption{Cross-campaign evidence architecture. ``Rows'' counts episode records and
``Inferential unit'' gives the seed or layout-cluster unit used for intervals. The
pathway-ablation and gain-sensitivity records form one block; transition values
are descriptive estimates that retain each view's inferential unit.}
\label{tab:evidence}
\small
\generatedtable{generated_table_evidence}
\end{table}

\begin{figure}[tbp]
\centering
\includegraphics[width=\linewidth]{fig14_evidence_summary.pdf}
\caption{\CapEvidence}
\label{fig:evidence}
\end{figure}

\subsection{The transition belongs to the baseline's integral action}

The primary ladder is a measurement of the stock controller's integral
action. The stock position layer uses integral gain \StockKi{}~N/m/s,
wind-up time \IntegralTauS{}~s, clamp \IntegralClampXy{}~m/s and a
\EpisodeHorizonS{}~s episode, so the largest steady integral force is
\MaxSteadyIntegralForceN{}~N. Under those settings the bound primary sweep
reproduces \AuditStockMatch{} of \AuditStockMatch{} checked episode fields;
that check confirms the audit flew the same stock controller, not that the
transition is invariant to Ki. Fig.~\ref{fig:integral-gain} shows the
0.15~N paired gain against Ki. Doubling Ki to 0.10~N/m/s collapses the
advantage from \AuditStockGain{} to \AuditKiTenGain{} (Wilcoxon
\AuditKiTenWilcoxonP{}, sign test \AuditKiTenSignP{}, $t$ interval
[\AuditKiTenCiLo{}, \AuditKiTenCiHi{}]), with baseline and aware completions
\AuditKiTenBase{} and \AuditKiTenAware{}. Both arms then first fall below
half-completion at \AuditKiTenKnee{}~N. Clamp 4 and clamp 8 change
\AuditClampChanged{} of the 0.15~N completions at the stock horizon, so the
pre-declared P1 verdict is \AuditPOneOverall{}: the Ki half is confirmed and
the clamp half is not, at \EpisodeHorizonS{}~s.

Lengthening the episode changes that clamp reading. At 12~s the stock
baseline completion at 0.15~N is \AuditAOneBTwelveBase{}, and raising the
clamp to 8 \AuditClampMattersLongHorizon{} change the 0.20~N completion. The
horizon sweep in Fig.~\ref{fig:horizon} makes the same point along time:
paired gain is \AuditGainTFour{} at 4~s, \AuditGainTSix{} at the stock
horizon, \AuditGainTEight{} at 8~s, then \AuditGainTTen{} and
\AuditGainTTwelve{} once $T\ge 10$~s, where the 0.15~N advantage vanishes.
The fine force ladder gives the transition a width rather than a single
rung: the baseline logistic 10--90 width is \AuditWidthTenNinety{}~N
(\AuditWidthCiLo{} to \AuditWidthCiHi{}), narrower than the 0.05~N primary
grid, so ``one of seven levels'' is a sampling artefact of that grid.

The assignment audit locates the gain outside the assignment path. On the
aware arm, the initial assignment differs from Euclidean in \AuditInitialDivergence{} of
\AuditAwareEpisodes{} episodes. Used (replan) assignments diverge in
\AuditUsedDivergence{} episodes, of which \AuditUsedDivergencePrimary{}
fall on the primary-sweep seeds and \AuditUsedDivergenceKnee{} sit at
0.15~N. A pre-declared sweep of the penalty weight over
$\lambda\in\{\LambdaGrid{}\}$~m/N (Fig.~\ref{fig:lambda-sweep};
\LambdaAwareEpisodes{} aware episodes per weight on the same seeds, forces
and frozen sources) bounds that inertness. The $\lambda=0.25$ row reproduces
the audit above exactly (\AuditInitialDivergence{} initial,
\LambdaUsedDivAtQuarter{} used-divergence episodes). Initial divergence is
zero at every tested weight through $\lambda=\LambdaInertMax{}$~m/N and first
appears at $\lambda=\LambdaFirstDivergence{}$~m/N (\LambdaInitialDivAtOne{}
of \LambdaAwareEpisodes{}; \LambdaInitialDivAtTwo{} at
$\lambda=\LambdaGridMax{}$~m/N), while used-assignment divergence rises
monotonically from \LambdaUsedDivAtZero{} at $\lambda=0$ to
\LambdaUsedDivAtTwo{} at $\lambda=\LambdaGridMax{}$~m/N. The aware
completion at 0.15~N is \LambdaCompletionFlat{} at every weight through
$\lambda=\LambdaInertMax{}$~m/N. Under the reading rule stated before the
sweep ran, this outcome keeps the inertness claim scoped to the inherited
weight and reports the tested range over which it holds; the sweep is
descriptive by declaration and carries no confirmatory test.

The high-force edge is a crash regime, not a clean saturation. Pooled at and
above 0.20~N, a fraction \AuditNeverCloseBase{} (agnostic) and
\AuditNeverCloseAware{} (aware) of airframes never come within 0.5~m of
their target. After a completion is recorded, a fraction
\AuditGroundContactAfter{} later contact the ground and
\AuditAltLossAfter{} lose more than 0.10~m of altitude, consistent with a
frozen rotor command rather than a held hover.

The capture-radius fingerprint likewise departs from a simple saturation
curve. Baseline $F_{50}$ against capture radius has slope
\AuditRadiusSlope{}~N/m (interval [\AuditRadiusSlopeLo{},
\AuditRadiusSlopeHi{}]), which excludes 1. A finer force grid over the same
seven radii (\AFiveBForceStepN{}~N steps, \AFiveBEpisodes{} episodes), added
after the pre-declaration, gives slope \AFiveBSlope{}~N/m (interval
[\AFiveBSlopeLo{}, \AFiveBSlopeHi{}]); it lies inside the pre-declared
interval and is reported as a resolution sensitivity beside the pre-declared
estimate, which remains the estimate of record.

\begin{figure}[tbp]
\centering
\includegraphics[width=\linewidth]{fig14_integral_gain.pdf}
\caption{\CapIntegralGain}
\label{fig:integral-gain}
\end{figure}

\begin{figure}[tbp]
\centering
\includegraphics[width=\linewidth]{fig15_horizon.pdf}
\caption{\CapHorizon}
\label{fig:horizon}
\end{figure}

\begin{figure}[tbp]
\centering
\includegraphics[width=\linewidth]{fig16_lambda_sweep.pdf}
\caption{\CapLambdaSweep}
\label{fig:lambda-sweep}
\end{figure}

\section{Discussion}

The identified mechanism is target-position feedforward compensating a slow
integral inside a \EpisodeHorizonS{}~s episode, and the controller audit
fixes the settings on which that identification rests. Prediction P1 is
\AuditPOneOverall{}: doubling Ki collapses the 0.15~N paired advantage to
\AuditKiTenGain{}, confirming the integral-gain half of the
pre-declaration, while clamp 4 and clamp 8 change \AuditClampChanged{} of the
0.15~N completions at the stock horizon, refuting the clamp half at
$T=\EpisodeHorizonS{}$~s. The same clamp
\AuditClampMattersLongHorizon{} matter at 12~s. Initial-assignment inertness
holds at every tested weight through $\lambda=\LambdaInertMax{}$~m/N and
first breaks at $\lambda=\LambdaFirstDivergence{}$~m/N; the bound claim stays
scoped to the inherited weight, and the sweep reports the tested range
descriptively. These bounds make the operating envelope a property of the
stock integral and the short horizon.

The headline gain travels through the feedforward path. Initial assignment
differs from Euclidean in \AuditInitialDivergence{} of \AuditAwareEpisodes{}
episodes, used (replan) assignments diverge in \AuditUsedDivergence{}
episodes, and the complete factorial localizes the completion change to
target-position feedforward, with the assignment main effect and the
assignment-by-feedforward interaction within \FactorialOtherAbsMax{}
completion of zero. The pre-declared $\lambda$ sweep is descriptive and stops
at $\lambda=\LambdaGridMax{}$~m/N; the allocation reading is bounded by that
grid. Support campaigns keep the same paired estimand: relative noise through
\BestSigma{} preserves the transition, and paired gains stay positive from
\ScaleMinSize{} through \ScaleMaxSize{} airframes. Rerunning the ladder
across capture radii keeps a transition of comparable size and moves it from
\RadiusTightPeakForce{} to \RadiusWidePeakForce{}~N, so \KneeLevel{}~N is
the transition of this operating point.

The geometry contrasts and the descriptive safety record bound the envelope
from two further directions. The contrasts map named scene axes within the
same paired design, and the clearance diagnostic marks the feasibility
boundary of the benchmark as a property of the task: the closest recorded
airframe pair is \SafetyWorstSeparation{}~m on targets drawn without a
separation constraint, and a physical deployment addresses that boundary by
composing the feedforward layer with an explicit safety barrier. Every
measurement here is a simulation. The present envelope is a computational
calibration reference against the stock integral, the \EpisodeHorizonS{}~s
horizon and the declared capture radius, and the transfer protocol names the
quantities a hardware or hardware-in-the-loop extension revalidates.
Heterogeneous vehicles, obstacle topology, penalty weights beyond
\LambdaGridMax{}~m/N and field-sensor traces are the named transfer axes for
that next stage; general allocator properties, hardware principles and field
results lie outside the measured scope.

\section*{Declarations}

%% The Declarations heading above must stay the first thing after the
%% Discussion prose (no comment or \ifdefined line before it): the copy-edit
%% checks slice the Discussion up to the next \section command and count
%% blank-line-separated chunks as paragraphs.
%% Funding and competing-interest wording adopted from the owner's packaging
%% script (scripts/build_springer_submission_package.py), shipped in every
%% packet since 2026-08-28 (ruling R16, 2026-09-09). FUNDING_AND_CONFLICTS
%% stays OPEN for the owner's final read.
\subsection*{Funding}

No external funding was received for this study.

\subsection*{Competing interests}

The author declares that they have no competing financial or non-financial
interests that are directly or indirectly related to the work submitted for
publication.

\subsection*{Ethics approval and consent to participate}

Not applicable. The study is computational and involves no human
participants, animal subjects, personally identifying data or field flights.

\subsection*{Consent for publication}

Not applicable; no identifiable participant data are included.

\subsection*{Data and code availability}

All simulation code, trajectory records and analysis scripts supporting the
findings of this study, which are the data and code behind every figure and
table, are available in the public repository and archive; the simulation
substrate is an open-source physics engine with the disturbance injection
described in the Methods.
%% Generated by scripts/release_targets.py from the release ledger.
%% Do not edit: an identifier typed here would not be checked against
%% anything, which is exactly the failure the ledger exists to prevent.
%% Ledger state: github=CREATED, zenodo=DEPOSITED
\ifdefined\ANONYMOUS That material is deposited in a public archive under a persistent identifier, and the development repository is public; both addresses are withheld from this review copy and appear in the identified version.\else That material is deposited under the persistent identifier \href{https://doi.org/10.5281/zenodo.22644574}{10.5281/zenodo.22644574}, and the development repository is at \url{https://github.com/PeterPonyu/wind-aware-allocation-band}.\fi


\subsection*{Materials availability}

Not applicable; the study generated no physical materials.

\subsection*{Author contributions}

\ifdefined\ANONYMOUS
The sole author performed conceptualization, methodology, software,
validation, formal analysis, investigation, visualization, writing of the
original draft, and review and editing.
\else
Zeyu Fu performed conceptualization, methodology, software, validation,
formal analysis, investigation, visualization, writing of the original draft,
and review and editing.
\fi

\subsection*{Use of generative AI tools}

Generative AI tools were used solely for code engineering, refactoring and
language polishing under the author's direction. All computational
experiments, analyses and manuscript text were reviewed, verified and edited
by the author, who retains full responsibility for the integrity and final
wording of the work.

%% Identified builds only; draft-only until the owner records the sentence.
\ifdefined\ANONYMOUS\else
\ifdefined\DRAFTTODO
\subsection*{Acknowledgements}

\ownertodo{add acknowledgements or delete this heading before submission;
none are recorded.}
\fi
\fi


{\footnotesize
\bibliographystyle{unsrt}
\bibliography{references}
}

\end{document}
");
%% every packet build leaves it undefined, so no unconfirmed declaration text
%% reaches a PDF. The owner replaces each use with the confirmed sentence.
\providecommand{\ownertodo}[1]{\ifdefined\DRAFTTODO[TODO-OWNER: #1]\fi}

%% Independent caption file. Every caption is self-contained: readable without
%% the body text, without cross-references, and without float numbers.

\newcommand{\CapSystemOverview}{%
Research object and endpoint bridge for the wind-aware allocation--control
measurement. Panel A is a drawn schematic of the computational contract: four
matched quadrotors, four waypoint targets and a horizontal disturbance vector.
Dashed paths represent the shared
assignment-to-control task; the two labelled paths distinguish the Euclidean
wind-agnostic cost from the wind-aware upwind cost plus target-position bias.
Panel B is a descriptive seed-paired view at the transition force
\KneeLevel{}~N: each arrow connects the same seed under the two arms in mean
closest-approach versus completion-rate coordinates. The vertical reference is
the per-airframe capture radius used to define completion. Displayed arrows use
the paired seed record as the unit; the panel reports a descriptive endpoint view
with the seed-level inferential unit preserved. %
}

\newcommand{\CapBand}{%
Operating envelope of task completion against steady horizontal wind force for a
wind-agnostic and a wind-aware allocator flying the same \NDrones{} quadrotors on
the same \NSeeds{} seeds per force level, with both panels on one shared force
axis. Panel A is the completion ladder: points are means over seeds and bars are
one standard deviation, clipped at the zero and one bounds of the rate, and the
upper axis restates force as a percentage of the \HoverThrust{}~N hover thrust of
one airframe. Panel B is the seed-paired advantage of the wind-aware over the
wind-agnostic allocator, in completion-rate points, with bars giving the
95~percent $t$ interval over the \NSeeds{} paired seeds. Panel B carries the
quantity the significance test is computed on, and its interval clears zero at
\KneeLevel{}~N alone, where the paired gain is $+\KneeDelta{}$ at an exact
two-sided signed-rank $p$ of \KneeP{}. The two panels are drawn together because
the Panel A bars are dispersions of two marginal means and overlap at the
transition even though every seed moves in the same direction, so the paired
contrast is displayed rather than left to be inferred from that overlap. Shading
separates the near-complete shoulder at and below \SaturationLevel{}~N and the
high-force edge at and above \CollapseLevel{}~N from the unshaded transition,
whose \KneeLevel{}~N level the dashed line marks and where directional
information has the largest operational leverage. Panel A's red and blue legend
identifies the two allocation arms; Panel B is a single paired-difference series
and therefore has no second arm legend.%
}

\newcommand{\CapNoise}{%
Profile of the directional signal as wind-estimate quality changes. The horizontal
axis is relative magnitude noise redrawn independently at every control step. Panel
A shows the paired advantage in completion rate at the transition force
\KneeLevel{}~N, with bars giving the 95~percent $t$ interval over \NSeeds{} paired
seeds and labels giving the exact two-sided signed-rank $p$ value; the dashed line
marks the decision level fixed before execution. Panel B places the two neighbouring
force levels on the same vertical scale, showing that the operating transition stays
localized while estimate quality changes. A paired advantage of zero marks equal
completion for the two arms. In both panels, paired advantage means aware-minus-
agnostic completion rate for the same seed. %
}

\newcommand{\CapSpread}{%
Distribution behind the two completion means at the transition force
\KneeLevel{}~N. Each episode assigns \NDrones{} quadrotors to \NDrones{} waypoints,
so a single episode can score zero, one quarter, one half, three quarters or one.
The bars count how many of the \NSeeds{} seeds land on each value for the two arms;
the distribution displays the episode-level margin carried by the aggregate
operating-envelope gain. %
}

\newcommand{\CapHeadings}{%
Coverage of the force headings drawn by the \NSeeds{} seeds and the repeat groups
in the displayed record. Panel A places each seed's heading on the circle, one
spoke per seed; the heading is fixed by the seed and reused at every force level,
in both allocators and in the estimate-noise sweep. The draws resolve to
\NHeadings{} distinct headings, with \NRepeats{} near-matched pairs and a largest
unvisited arc of \UnsampledArc{} degrees. Panel B orders all
$\NSeeds{}(\NSeeds{}-1)/2$ pairwise separations on a logarithmic axis; the empty
gap between the retained group and the next separation makes the stated grouping
stable across the displayed tolerance interval. %
}

\newcommand{\CapRepeats}{%
Context spread under near-matched force headings. At each force level the bar gives
the largest completion difference between two episodes of the same allocator whose
headings differ by less than \RepeatTol{} degrees, expressed as the number of
airframes, out of \NDrones{}, by which the episodes differ. The strip below records
whether any seeds at that force level varied in completion; a zero bar at a fully
uniform level marks a saturated shoulder or floor, while the informative levels
show where waypoint context changes the outcome under nearly identical directional
input. %
}

\newcommand{\CapDesigned}{%
Heading-by-layout operating-envelope map from a crossed design. \NAssignedHeadings{}
headings \HeadingStep{} degrees apart are each flown with all \NScenes{} waypoint
layouts, at every force level and under both allocators, for \NEpisodesDesigned{}
episodes and \NPairsPerForce{} paired comparisons per level. Panel A shows paired advantage
against force, with the band giving an ordinary-bootstrap 95~percent interval over
the \NCrossedClusters{} layout-level means after averaging the fixed headings
(\NCrossedBootstrap{} deterministic resamples); crosses mark the \NSeeds{} seed
estimates from the independently drawn-heading condition.
Panel B opens the transition force into its heading-by-layout rectangle, one cell
per paired episode; rows are assigned headings and columns are scene layouts, and
colour is the paired completion advantage. The heat map displays the same cells
used to compute the heading, layout and interaction shares. %
}

\newcommand{\CapMiss}{%
Continuous tracking quantity behind the thresholded completion endpoint. Points give
the mean closest approach of each quadrotor to its assigned waypoint, averaged over
seeds, with bars of one standard deviation; the dashed line is the
\CaptureRadius{}~m capture radius that defines completion. Both allocators are
measured at the same forces, and the two series are offset horizontally by a fixed
0.0055~N purely so that their overlapping bars remain separable; the axis ticks
mark the true force levels. The two allocators occupy
nearby trajectories across the force ladder, and the measured distance shift at the
transition moves aircraft across the capture boundary to produce the task-level
operating gain. %
}

\newcommand{\CapTemporal}{%
Temporal operating-envelope extension for a causal directional estimate. Panel A
shows paired completion gain against true force for constant, slow-gust and fast-gust
profiles, with exact, causal and biased-causal estimate conditions and 95~percent
intervals over \NTemporalSeeds{} paired seeds. Panel B shows mean relative vector-tracking error
for the same profiles and estimate conditions. The transition remains concentrated
near \KneeLevel{}~N, while the tracking-error profile quantifies the cost of temporal
rate and links estimator dynamics to the measured task-level response. Colours and
symbols identify the three estimate conditions consistently across both panels. %
}

\newcommand{\CapMechanism}{%
Mechanism and gain-sensitivity extension at the transition force \KneeLevel{}~N.
Panel A gives paired completion gain across the tested target-feedforward gains for
each true-force profile; the dashed line marks the inherited gain
\FeedforwardGainDefault{}~m/N. Panel B compares the complete arm with an
allocation-cost-only ablation that retains directional assignment and selectively
removes the target-position bias. The separation identifies the pathway through
which the directional estimate is converted into task completion in this benchmark.
Intervals in both panels are paired 95~percent $t$ intervals over the declared
profile-specific seeds; the upper gain-grid point marks the measured edge of the
tested gain envelope. %
}

\newcommand{\CapGeometry}{%
Factorised scene-geometry extension of the directional operating envelope. Panel A
shows high-minus-low contrasts for target spread, target range and crossing
orientation at each tested force; intervals are 95~percent cluster-aware intervals
over 30 seed-level factor contrasts after averaging the four crossed settings
within seed; the 120-term version is sensitivity-only. Raw exact signed-rank
values and nine-contrast Holm adjustments define the support-family readout, and
the zero line marks no change in the wind-aware advantage. Panel B shows the mean
paired advantage in all eight scene conditions
across the three force levels, with each row defined by one spread, range and
crossing setting. The nominal--near--aligned cell is the legacy compatibility
condition; the remaining cells are deterministic scene transformations. Together,
the panels define the tested simulator geometry envelope and its named transfer
axes. %
}

\newcommand{\CapScale}{%
Homogeneous swarm-size extension of the directional operating envelope. Panel A
shows the paired completion gain against team size for two, four, six and eight
airframes at the three tested force levels; points are means over the same 30
paired seeds per cell and bars are 95~percent $t$ intervals. Raw exact signed-rank
values and their Holm adjustments are reported over the 12 size--force cells.
Panel B opens the transition force at \KneeLevel{}~N and shows the wind-agnostic and
wind-aware completion endpoints separately, with means and 95~percent intervals over
the same seed pairs. The constant-profile four-airframe slice contains 180 rows:
30 seeds, two arms and three force levels with the causal estimator. The ten-seed
exact-vector condition is a distinct oracle slice, and the other team sizes are
declared simulator cells. Together the panels define the measured homogeneous-size
transfer envelope. %
}

\newcommand{\CapRadius}{%
Sensitivity of the operating envelope to the capture radius that defines task
completion, over \NRadiusEpisodes{} episodes at radii of \RadiusGrid{}~m and
\NRadiusSeeds{} paired seeds per radius and force. Panel A gives the paired
completion advantage of the wind-aware over the wind-agnostic allocator against
steady horizontal wind force, one series per radius, with bars showing 95~percent
intervals over the paired seeds and the series offset horizontally so that
overlapping bars remain separable. The labelled point on each series is that
radius' largest advantage. Panel B gives the two completion ladders those
advantages are differences of, one facet per radius, with the dashed rule marking
the same labelled force. The transition survives every tested radius at a
comparable size, between $+\RadiusPeakDeltaMin{}$ and $+\RadiusPeakDeltaMax{}$
completion, and its location moves with the radius, from
\RadiusTightPeakForce{}~N at the tightest radius to \RadiusWidePeakForce{}~N at
the widest, so the force at which directional information pays is a property of
the operating point. Within this \RadiusHolmFamily{}-test validation family
\RadiusHolmSurvivors{} cell clears a Holm-corrected 5~percent threshold, a
resolution set by the \NRadiusSeeds{} seeds per cell; the family is a
sensitivity readout beside the larger primary sweep. %
}

\newcommand{\CapPathwaySafety}{%
Complete pathway decomposition and the descriptive safety diagnostics, from two
separate campaigns of \NFactorialEpisodes{} and \NSafetyEpisodes{} episodes.
Panel A gives all three terms of the assignment-by-feedforward factorial as
paired completion effects against true horizontal force, one facet per temporal
profile, with bars showing 95~percent intervals over \NFactorialSeeds{} paired
seeds and the terms offset horizontally for legibility. The target-feedforward
main effect carries the whole gain at the transition force, between
$+\FactorialFeedforwardMin{}$ and $+\FactorialFeedforwardMax{}$ completion across
profiles; the assignment main effect and the interaction stay within
\FactorialOtherAbsMax{} of zero and neither clears its Holm-corrected threshold at
any force. Panel B is descriptive only and carries no test: for each allocator and
force it gives the range of the per-episode minimum pairwise airframe separation
over \NSafetySeeds{} seeds, with the point at the mean of that range and the
dashed rule at the closest airframe pair recorded anywhere in the campaign,
\SafetyWorstSeparation{}~m. The panel is a benchmark-level feasibility diagnostic;
because target placement and arrival status contribute to the metric, it is not
interpreted as an allocator-causal clearance effect. %
}

\newcommand{\CapEvidence}{%
Cross-campaign evidence synthesis for the directional operating envelope. Panel A
places the transition paired-completion estimates from the primary, estimate-noise,
crossed-layout, causal-temporal and homogeneous-size views on a common rate axis;
bars are the view-specific 95~percent intervals and colours identify the evidence
block. The crossed-layout estimate is shown with its interval because that record
uses a layout-cluster bootstrap. Panel B shows the episode-row volume of each of the
eight views on a logarithmic axis, with the text next to each point naming the seed
or layout-cluster unit that supports its analysis; pathway ablation and gain
sensitivity are two views of one pathway/gain synthesis block. The figure is
descriptive by design and preserves each view's inferential unit. %
}

\newcommand{\CapIntegralGain}{%
Paired completion gain at the \KneeLevel{}~N force against the baseline integral
gain, from the controller-identification sweep at the stock
\EpisodeHorizonS{}~s horizon. Points are the seed-paired aware-minus-agnostic
completion difference; bars are 95~percent $t$ intervals. Clamp settings of
\IntegralClampXy{}, 4 and 8~m/s are drawn as distinct markers at the same Ki,
so a clamp that leaves the 0.15~N completions unchanged sits on top of the
stock marker. Doubling the stock Ki of \StockKi{}~N/m/s collapses the gain from
\AuditStockGain{} to \AuditKiTenGain{} (interval
[\AuditKiTenCiLo{}, \AuditKiTenCiHi{}]). %
}

\newcommand{\CapHorizon}{%
Horizon sweep under the stock controller (Ki~\StockKi{}~N/m/s, clamp
\IntegralClampXy{}~m/s). Panel A is the paired completion gain at
\KneeLevel{}~N against episode duration, with bars giving the 95~percent $t$
interval; the gain is \AuditGainTSix{} at the stock \EpisodeHorizonS{}~s
horizon and \AuditGainTTen{} once the horizon reaches 10~s. Panel B is the
matching wind-agnostic completion at the same force, which rises as the
episode lengthens and the slow integral is given more time to wind up. Both
panels share one horizon axis. %
}

\newcommand{\CapLambdaSweep}{%
Pre-declared sweep of the assignment penalty weight
$\lambda\in\{\LambdaGrid{}\}$~m/N under the stock controller at the
\EpisodeHorizonS{}~s horizon, \LambdaAwareEpisodes{} aware-arm episodes per
weight. Circles count episodes whose initial ($t=0$) assignment differs from
the Euclidean assignment; squares count episodes in which any call,
initial or replan, differs. The dashed line marks the inherited weight
$\lambda=\LambdaReference{}$, at which the counts reproduce the bound
assignment audit exactly. Initial divergence is zero through
$\lambda=\LambdaInertMax{}$ and first appears at
$\lambda=\LambdaFirstDivergence{}$; the horizontal axis is pseudo-logarithmic
so that $\lambda=0$ sits on the same axis. The sweep is descriptive by
declaration and carries no confirmatory test. %
}
